\chapter{Objects in the permanent $h_6^2$ statement}
\label{chap:h6-statement}

This chapter constructs the objects needed to state the permanent-cycle
target.  The reduction to $d_{12}$ and the conditions used in its exclusion
belong to the later proof, not to these definitions.  The ambient category
of spectra is abstract.  Its mod--$2$ Eilenberg--Mac Lane object and the
explicit Milnor cooperation coordinates below are foundational inputs.
The generic implementation still takes those coordinates as an argument.
The final targets now select one fixed foundation and coordinates through
explicitly named development axioms; removing parameters does not discharge
those dependencies.
The next comparison stage keeps the abstract background and derives them
from the ring structure, cooperations, and tensor compatibility of $H$;
that derivation is not yet implemented.

% SOURCE: PROJECT_BOUNDARY.md, Confirmed design decision 1; abstract foundation
% and Milnor cooperation boundary explicitly confirmed for this implementation.
\begin{definition}[Mod--2 Eilenberg--Mac Lane object and unit]
  \label{def:h6-em-object}
  \lean{KIP126.StableHomotopy.Cohomology.Mod2EilenbergMacLane}
  \lean{KIP126.StableHomotopy.Cohomology.Mod2EilenbergMacLane.unit}
  \uses{def:stable-homotopy-group,def:mathlib-f2}
  \notready
  Fix an object $H$ with an additive isomorphism
  $\pi_0H\cong\mathbb F_2$ and $\pi_nH=0$ for $n\ne0$.
  Define $u:S^0\to H$ to be the map corresponding to $1\in\mathbb F_2$.
  The equation $\pi_nH=0$ means that every element of that group equals its
  zero element; it does not mean that the group is empty.
\end{definition}

\begin{definition}[Ring structure and cooperation maps of the chosen object]
  \label{def:h6-ring-cooperations}
  \lean{KIP126.StableHomotopy.Cohomology.Mod2RingStructure}
  \lean{KIP126.StableHomotopy.Cohomology.Mod2Cooperations}
  \lean{KIP126.StableHomotopy.Cohomology.cooperationDiagonalMap}
  \lean{KIP126.StableHomotopy.Cohomology.cooperationCounit}
  \uses{def:h6-em-object,def:unital-multiplication-data}
  \leanok
  Equip the specified $H$ with an associative unital algebra-object
  structure whose unit is the specified $u$. Define the graded cooperation
  groups to be $\pi_n(H\wedge H)$. Inserting $u$ between the two factors
  gives a map to $\pi_n(H\wedge(H\wedge H))$, and multiplication gives
  the counit to $\pi_nH$. No K\"unneth identification or Milnor polynomial
  coordinates are part of these definitions.
\end{definition}

% SOURCE: the Adams resolution in def:hf2-adams-tower, made explicit by
% iterating the unit and its fiber; aimpaper/main.tex:253-263.
\begin{definition}[Iterated-fiber Adams tower]
  \label{def:h6-fiber-tower}
  \lean{KIP126.Classical.Adams.fiber}
  \lean{KIP126.Classical.Adams.fiberι}
  \lean{KIP126.Classical.Adams.adamsUnit}
  \lean{KIP126.Classical.Adams.adamsTower}
  \lean{KIP126.Classical.Adams.adamsTowerStep}
  \lean{KIP126.Classical.Adams.adamsTowerComposite}
  \lean{KIP126.Classical.Adams.adamsTowerMapAt}
  \lean{KIP126.Classical.Adams.adamsLayerAt}
  \uses{def:h6-em-object,def:cofiber-triangle,def:stable-sphere-suspension}
  \notready
  For $X$, put $T_0=X$ and
  \[
    T_{s+1}=\operatorname{fib}(T_s\longrightarrow H\wedge T_s).
  \]
  The displayed map is the unit $u$ smashed with $T_s$, after the unit
  isomorphism.  Define fibers by desuspending the chosen cofibers.
  Compose the fiber inclusions to obtain $T_b\to T_a$ for $a\le b$,
  extend the tower constantly to negative stages, and put
  $Q_s=\operatorname{cofib}(T_{s+1}\to T_s)$.
\end{definition}

\begin{proposition}[The Adams unit on the coefficient object splits]
  \label{prop:h6-unit-splits}
  \lean{KIP126.StableHomotopy.Cohomology.adamsUnit_mul}
  \lean{KIP126.StableHomotopy.Cohomology.cooperationDiagonal_counit_right}
  \uses{def:h6-ring-cooperations,def:h6-fiber-tower}
  \leanok
  Multiplication $H\wedge H\to H$ retracts the map
  $H\to H\wedge H$ used in the Adams tower. Multiplying the final two
  factors also retracts the middle-unit insertion
  $H\wedge H\to H\wedge(H\wedge H)$.
\end{proposition}
\begin{proof}
  Apply the left unit identity of the algebra object, then tensor that
  identity with $H$.
  \leanok
\end{proof}

\begin{proposition}[The tower layers are smashed coefficient objects]
  \label{prop:h6-tower-layer-comparison}
  \lean{KIP126.Classical.Adams.cofiberFiberIso,KIP126.Classical.Adams.adamsLayerTriangleIso,KIP126.Classical.Adams.adamsLayerIso,KIP126.Classical.Adams.adamsLayerIso_ι,KIP126.Classical.Adams.adamsLayerIso_δ}
  \uses{def:h6-fiber-tower,def:cofiber-triangle}
  \leanok
  In any of the specified stable contexts, for every unit $u:S^0\to H$,
  spectrum $X$, and $s\ge0$, the actual tower layer $Q_s$ is isomorphic to
  $H\wedge T_s$. The comparison sends the chosen map $T_s\to Q_s$ to
  the Adams unit $T_s\to H\wedge T_s$ and is part of an isomorphism of
  distinguished triangles, so it also controls the connecting map.
  This statement uses no Milnor coordinates, Lin data, multiplication on
  $H$, or convergence assertion.
\end{proposition}
\begin{proof}
  Inverse-rotate the cofiber triangle of the Adams unit. Its first map is
  $-\mathrm{fiber}\iota$, whereas the implemented tower uses the positive
  $\mathrm{fiber}\iota$. Complete the negative identity on the fiber and
  identity on $T_s$ to an isomorphism of distinguished triangles. Its third
  component is the required layer comparison. The commutative squares give
  the unit-map equality and the connecting-map equality with its required
  minus sign.
  \leanok
\end{proof}

\begin{definition}[The iterated-smash comparison tower]
  \label{def:h6-smash-comparison-tower}
  \lean{KIP126.Classical.Adams.adamsSmashTower,KIP126.Classical.Adams.adamsSmashTowerStep,KIP126.Classical.Adams.adamsSmashTowerComposite}
  \uses{def:h6-fiber-tower}
  \leanok
  Write $\overline H=\operatorname{fib}(S\to H)$ with its existing inclusion
  $b:\overline H\to S$. Set $P_0X=X$ and $P_{s+1}X=\overline H\wedge P_sX$.
  The transition $P_{s+1}X\to P_sX$ is $b\wedge\mathrm{id}$ followed by
  the left unit isomorphism. This comparison model does not replace the
  actual iterated-fiber tower or its internal spectral sequence.
\end{definition}

\begin{proposition}[Tensor exactness identifies a fiber step, including its maps]
  \label{prop:h6-fiber-smash-step}
  \lean{KIP126.Classical.Adams.adamsTensorFiberTriangleIso,KIP126.Classical.Adams.adamsFiberTensorIso,KIP126.Classical.Adams.adamsTensorFiberTriangleIso_hom₂,KIP126.Classical.Adams.adamsTensorFiberTriangleIso_hom₃,KIP126.Classical.Adams.adamsFiberTensorIso_inv_δ,KIP126.Classical.Adams.adamsFiberTensorIso_inv_ι,KIP126.Classical.Adams.adamsFiberTensorIso_hom_ι}
  \uses{def:h6-fiber-tower,def:cofiber-triangle}
  \leanok
  If $-\wedge X$ is shift-compatible and preserves distinguished triangles,
  smashing the unit fiber triangle gives an isomorphism of triangles with
  the fiber triangle of $X\to H\wedge X$. Its second and third components
  are the left unit isomorphism and identity, respectively. Its first
  component gives
  $e_X:\operatorname{fib}(X\to H\wedge X)\simeq\overline H\wedge X$.
  Under monoidal additivity,
  $e_X\,(b\wedge\mathrm{id})\,\lambda_X=\mathrm{fiber}\iota$.
  The connecting-map square is also proved, with the functor's suspension
  comparison retained. These are consequences of exactness, not extra
  fiber/tensor comparison assumptions; no ring, Milnor, or Lin data enter.
\end{proposition}
\begin{proof}
  Rotate the two distinguished triangles, complete the unitor and identity
  to an isomorphism, and inverse-rotate back. The first square and
  additivity cancel the two inverse-rotation minus signs. The third square
  supplies the connecting-map identity.
  \leanok
\end{proof}

\begin{proposition}[The actual tower agrees with the iterated-smash tower]
  \label{prop:h6-tower-smash-comparison}
  \lean{KIP126.Classical.Adams.adamsTowerSmashIso,KIP126.Classical.Adams.adamsTowerSmashIso_step,KIP126.Classical.Adams.adamsTowerSmashIso_composite}
  \uses{def:h6-smash-comparison-tower,prop:h6-fiber-smash-step}
  \leanok
  Assume the preceding exactness and shift structures for every right tensor
  functor, and monoidal additivity. There are recursively constructed
  isomorphisms $e_s:T_sX\simeq P_sX$ satisfying
  $e_{s+1}P_{s+1,s}=T_{s+1,s}e_s$, and the analogous equality for every
  finite composite transition. Thus this is a comparison of towers and
  their maps, not just a list of object isomorphisms. These generic
  statements do not supply the tensor-exactness witnesses on the fixed
  foundation. These comparisons alone do not establish a multiplicative
  tower, naturality of the chosen comparisons in $X$, or Leibniz.
\end{proposition}
\begin{proof}
  Iterate $e_{T_sX}$ and tensor the previously constructed comparison by
  $\overline H$. The transition equality follows from the one-step inclusion
  comparison and unitor naturality. Induct on the composite's length.
  \leanok
\end{proof}

\begin{proposition}[Finite-stage homotopy lifts transfer to the smash model]
  \label{prop:h6-smash-liftability}
  \lean{KIP126.Classical.Adams.adamsTowerSmashIso_lift_iff}
  \uses{prop:h6-tower-smash-comparison}
  \leanok
  For $x\in\pi_n(T_sX)$ and $\ell\ge0$, a preimage of $x$ under
  $\pi_n(T_{s+\ell}X)\to\pi_n(T_sX)$ exists if and only if the image
  $(e_s)_*x$ has a preimage under
  $\pi_n(P_{s+\ell}X)\to\pi_n(P_sX)$.
  This transfers the finite lifting questions used in the cycle definitions;
  it does not assert that any such lift exists.
\end{proposition}
\begin{proof}
  Compose a lift with $e_{s+\ell}$ in one direction and its inverse in the
  other. Use preservation of composite transitions and cancel the target
  isomorphism.
  \leanok
\end{proof}

\begin{definition}[Concatenation of sphere smash stages]
  \label{def:h6-smash-sphere-pairing}
  \lean{KIP126.Classical.Adams.adamsSmashSphereTensorIso,KIP126.Classical.Adams.adamsSmashTowerNestingIso,KIP126.Classical.Adams.adamsSmashSpherePairingIso}
  \uses{def:h6-smash-comparison-tower}
  \leanok
  Using only associators and the sphere unitor, construct
  $P_sS\wedge Y\simeq P_sY$ and $P_s(P_tX)\simeq P_{t+s}X$.
  Their composite gives the concrete concatenation pairing
  $m_{s,t}:P_sS\wedge P_tS\simeq P_{t+s}S$.
  The reversed display order of the sum supports recursion in $s$; it is
  ordinary addition of filtration degrees, not a different grading.
\end{definition}

\begin{proposition}[Concatenation commutes with the first-input transition]
  \label{prop:h6-smash-pairing-left-transition}
  \lean{KIP126.Classical.Adams.adamsSmashSpherePairingIso_zero,KIP126.Classical.Adams.adamsSmashSpherePairingIso_succ,KIP126.Classical.Adams.adamsSmashSpherePairingIso_step_left}
  \uses{def:h6-smash-sphere-pairing}
  \leanok
  The zero-stage pairing is the left unitor; the successor pairing is the
  associator followed by $\overline H\wedge m_{s,t}$.
  Writing $p_s:P_{s+1}S\to P_sS$ for the specified transition,
  $m_{s+1,t}p_{t+s}=(p_s\wedge\mathrm{id})m_{s,t}$.
  This statement concerns only the first input; the second-input transition
  law is not claimed here.
\end{proposition}
\begin{proof}
  Reassociate, use the interchange law for the inclusion
  $\overline H\to S$, and apply associator and unitor naturality.
  \leanok
\end{proof}

\begin{definition}[The transported pairing on the existing sphere tower]
  \label{def:h6-existing-tower-pairing}
  \lean{KIP126.Classical.Adams.adamsTowerSpherePairingIso}
  \uses{def:h6-smash-sphere-pairing,prop:h6-tower-smash-comparison}
  \leanok
  Under the explicit tensor-exactness and shift structures, transport
  concatenation through the already constructed comparisons $e_s:T_sS\simeq P_sS$:
  $\mu_{s,t}=(e_s\wedge e_t)m_{s,t}e_{t+s}^{-1}$.
  This is a concrete map on the existing iterated-fiber tower; no additional
  pairing, ring structure, or spectral sequence is postulated.
\end{definition}

\begin{proposition}[The transported pairing preserves the first-input transition]
  \label{prop:h6-existing-tower-pairing-left-transition}
  \lean{KIP126.Classical.Adams.adamsTowerSpherePairingIso_comparison,KIP126.Classical.Adams.adamsTowerSpherePairingIso_step_left}
  \uses{def:h6-existing-tower-pairing,prop:h6-smash-pairing-left-transition}
  \leanok
  With monoidal additivity, the pairing satisfies the defining comparison
  square and
  $\mu_{s+1,t}T_{t+s+1,t+s}=(T_{s+1,s}\wedge\mathrm{id})\mu_{s,t}$
  for the actual sphere tower maps. This does not yet assert a full
  multiplicative tower: the second-input transition is proved only under the
  additional condition below; its fixed witness, the full two-term boundary
  product formula, and comparison with Lin multiplication remain unproved.
  In particular it is not a witness of \texttt{SecondDifferentialLeibniz}.
\end{proposition}
\begin{proof}
  Postcompose with the target comparison isomorphism. Use its preservation
  of tower transitions, the smash pairing's first-input law, and tensor
  functoriality. Cancel the target isomorphism.
  \leanok
\end{proof}

\begin{definition}[The two-factor inclusion condition]
  \label{def:h6-unit-fiber-inclusion-commutes}
  \lean{KIP126.Classical.Adams.UnitFiberInclusionCommutes}
  \uses{def:h6-smash-comparison-tower}
  \leanok
  Write $B=\overline H$ and $b:B\to S$ for the actual unit-fiber inclusion.
  The condition is the equality of maps $B\wedge B\to B$
  \[
    (b\wedge\mathrm{id}_B)\lambda_B
      =(\mathrm{id}_B\wedge b)\rho_B.
  \]
  This defines a predicate, not an axiom or a witness. It is not currently
  derived from the specified multiplication on $H$ or tensor exactness.
  In particular the definition's formalization status does not claim that
  the condition holds for the fixed sphere foundation.
\end{definition}

\begin{proposition}[Conditional right-input compatibility in the smash model]
  \label{prop:h6-smash-pairing-right-transition}
  \lean{KIP126.Classical.Adams.unitFiberInclusionCommutes_tensor,KIP126.Classical.Adams.adamsSmashTowerStep_succ_of_inclusion_commutes,KIP126.Classical.Adams.adamsSmashSpherePairingNextRight,KIP126.Classical.Adams.adamsSmashSpherePairingNextRight_zero,KIP126.Classical.Adams.adamsSmashSpherePairingNextRight_succ,KIP126.Classical.Adams.adamsSmashSpherePairingIso_step_right}
  \uses{def:h6-unit-fiber-inclusion-commutes,def:h6-smash-sphere-pairing}
  \leanok
  Assume the two-factor inclusion condition. After tensoring by any $X$,
  \[
    (b\wedge\mathrm{id}_{B\wedge X})\lambda_{B\wedge X}
      =B\wedge((b\wedge\mathrm{id}_X)\lambda_X).
  \]
  Thus $p_{s+1}=B\wedge p_s$ on the smash tower. Reindex $m_{s,t+1}$
  along $(t+1)+s=(t+s)+1$, calling the resulting map $m^+_{s,t}$.
  Then $m^+_{s,t}p_{t+s}=(\mathrm{id}_{P_sS}\wedge p_t)m_{s,t}$
  for every $s,t$. No new pairing is chosen in this reindexing.
\end{proposition}
\begin{proof}
  Tensor the assumed two-factor equality with $X$ and use associator
  naturality and the triangle identity. Induct on $s$, using the zero and
  successor formulas for the reindexed pairing and $p_{s+1}=B\wedge p_s$.
  These are generic conditional proofs; they do not supply the hypothesis.
  \leanok
\end{proof}

\begin{proposition}[Conditional right-input compatibility on the actual tower]
  \label{prop:h6-existing-tower-pairing-right-transition}
  \lean{KIP126.Classical.Adams.adamsTowerSpherePairingNextRight,KIP126.Classical.Adams.adamsTowerSpherePairingNextRight_comparison,KIP126.Classical.Adams.adamsTowerSpherePairingIso_step_right}
  \uses{def:h6-existing-tower-pairing,prop:h6-smash-pairing-right-transition}
  \leanok
  Under the preceding tensor-exactness, shift and additivity structures and
  the two-factor inclusion condition, the existing actual tower pairing
  satisfies
  $\mu^+_{s,t}T_{t+s+1,t+s}=(\mathrm{id}_{T_sS}\wedge T_{t+1,t})\mu_{s,t}$.
  Here $\mu^+$ is the existing pairing with the same successor reindexing,
  and its comparison square with $m^+$ commutes.
\end{proposition}
\begin{proof}
  Use naturality of the equality transport with respect to $e_s$. Then
  postcompose with the target comparison, apply its transition law and the
  conditional smash-model law, and cancel that comparison.
  \leanok
\end{proof}

\begin{proposition}[Fiber inclusions are injective on maps under orthogonality]
  \label{prop:h6-fiber-mapping-injective}
  \lean{KIP126.Classical.Adams.fiberι_postcomp_injective_of_vanishing}
  \uses{def:h6-fiber-tower}
  \leanok
  For a map $f:X\to Y$ and an object $Z$, suppose every map
  $Z\to\Sigma^{-1}Y$ is zero. Postcomposition with the actual inclusion
  $\operatorname{fib}(f)\to X$ is injective on maps out of $Z$.
  This does not assert that the fiber inclusion is a monomorphism for all
  source objects.
\end{proposition}
\begin{proof}
  If two maps have equal composites, their difference maps to zero in $X$.
  The twice inverse-rotated cofiber triangle factors that difference through
  $\Sigma^{-1}Y$. Apply the vanishing hypothesis. The negative sign on the
  first map of the fiber triangle is explicitly retained and cancelled.
  \leanok
\end{proof}

\begin{proposition}[Deriving the two-factor equality from mapping-group vanishing]
  \label{prop:h6-unit-fiber-commutes-from-vanishing}
  \lean{KIP126.Classical.Adams.unitFiberInclusion_postcomp_eq,KIP126.Classical.Adams.unitFiberInclusionCommutes_of_vanishing,KIP126.Classical.Adams.adamsTowerSpherePairingIso_step_right_of_vanishing}
  \uses{prop:h6-fiber-mapping-injective,def:h6-unit-fiber-inclusion-commutes,prop:h6-existing-tower-pairing-right-transition}
  \leanok
  The two elimination maps $B\wedge B\to B$ always have equal composites
  with $b:B\to S$. If $[B\wedge B,\Sigma^{-1}H]=0$, they are equal
  already as maps to $B$. Consequently the actual tower's right-input law
  follows from this vanishing and the previously stated tensor structures,
  without taking the two-factor equality itself as an input.
\end{proposition}
\begin{proof}
  Unitor naturality and tensor interchange identify the two composites to
  $S$. Apply the injectivity theorem to $f:S\to H$ and $Z=B\wedge B$,
  then use the already proved conditional right-input law.
  \leanok
\end{proof}

\begin{proposition}[A t-structure criterion for the right-input law]
  \label{prop:h6-unit-fiber-commutes-from-connectivity}
  \lean{KIP126.Classical.Adams.unitFiber_mapping_vanishing_of_tStructure,KIP126.Classical.Adams.unitFiberInclusionCommutes_of_tStructure,KIP126.Classical.Adams.adamsTowerSpherePairingIso_step_right_of_tStructure}
  \uses{prop:h6-unit-fiber-commutes-from-vanishing}
  \leanok
  Give $C$ a t-structure in Mathlib's cohomological indexing. If
  $B\wedge B\in C^{\le0}$ and $H\in C^{\ge0}$, then
  $[B\wedge B,\Sigma^{-1}H]=0$. In particular, it suffices that
  $B\in C^{\le0}$ and $C^{\le0}$ is closed under the tensor product.
  These lower conditions imply the two-factor equality and, with the
  preceding tensor structures, the right-input law on the actual tower.
  No t-structure or connectivity witness for the fixed foundation is
  constructed or postulated by these theorems.
\end{proposition}
\begin{proof}
  The shift convention gives $\Sigma^{-1}H\in C^{\ge1}$. Orthogonality
  to $C^{\le0}$ makes all maps zero. Apply the preceding proposition.
  \leanok
\end{proof}

\begin{proposition}[Homotopy injectivity of the actual unit-fiber inclusion]
  \label{prop:h6-unit-fiber-homotopy-injective}
  \lean{KIP126.Classical.Adams.fiberι_homotopy_injective_of_surjective,KIP126.Classical.Adams.mod2Unit_homotopy_surjective,KIP126.Classical.Adams.mod2UnitFiber_homotopy_injective,KIP126.Classical.Adams.mod2UnitFiber_homotopy_subsingleton}
  \uses{def:h6-fiber-tower,prop:h6-sphere-unit-image}
  \leanok
  For any map $f:X\to Y$, surjectivity of $\pi_nf$ implies injectivity of
  $\pi_{n-1}(\operatorname{fib}(f))\to\pi_{n-1}X$.
  For the supplied mod--2 Eilenberg--Mac Lane object, the actual unit
  $S\to H$ is surjective on every homotopy group. Thus
  $\pi_n(B)\to\pi_n(S)$ is injective for every integer $n$.
  In particular any vanishing degree of the sphere is also a vanishing
  degree of $B$. This does not claim that $B\to S$ is a categorical
  monomorphism or that sphere connectivity is already supplied.
\end{proposition}
\begin{proof}
  Use the long exact sequence of the actual fiber triangle: surjectivity
  makes the connecting map zero, so the preceding kernel is zero.
  The mod--2 unit surjectivity follows from the existing sphere unit-image
  theorem after the right unitor. Apply the general result in degree $n+1$.
  No ring, tensor exactness, Milnor or Lin input is used in this argument.
  \leanok
\end{proof}

\begin{proposition}[Deriving the needed bounds from sphere homotopy]
  \label{prop:h6-unit-fiber-connectivity-from-sphere}
  \lean{KIP126.Classical.Adams.mod2UnitFiber_isLE_of_homotopy,KIP126.Classical.Adams.mod2EilenbergMacLane_isGE_of_homotopy,KIP126.Classical.Adams.mod2UnitFiber_inclusionCommutes_of_homotopy,KIP126.Classical.Adams.mod2SpherePairing_step_right_of_homotopy}
  \uses{prop:h6-unit-fiber-homotopy-injective,prop:h6-unit-fiber-commutes-from-connectivity}
  \leanok
  Suppose the t-structure detects bounds from homotopy: vanishing in all
  negative degrees implies $C^{\le0}$, and vanishing in all positive
  degrees implies $C^{\ge0}$. Suppose also $\pi_nS=0$ for $n<0$.
  Then $B\in C^{\le0}$ and $H\in C^{\ge0}$ are derived facts, not
  separate inputs. With connective tensor closure this proves the
  two-factor inclusion condition; with the preceding tensor comparison
  structures it proves the right-input transition law on the actual
  mod--2 sphere tower. All these lower hypotheses remain explicit; no
  t-structure or sphere-connectivity witness for the fixed foundation is
  constructed here.
\end{proposition}
\begin{proof}
  Apply the preceding injectivity to each negative homotopy group and then
  the connective detection hypothesis. Apply the supplied Eilenberg--Mac Lane
  vanishing to each positive degree to obtain the other bound. Invoke the
  previously proved t-structure criterion and tower transport theorem.
  \leanok
\end{proof}

\paragraph{Classical interpretation and the missing fixed witnesses.}
For spectra, $C^{\le0}$ here means connective spectra; it is the reverse
indexing of the usual homological notation. Smash products preserve this
class: see Lurie, \emph{Derived Algebraic Geometry II}, Lemma 4.3.5,
\url{https://arxiv.org/pdf/math/0702299}. In the classical sphere situation,
the long exact sequence and the surjection $\mathbb Z\to\mathbb F_2$
make $B=\operatorname{fib}(S\to H\mathbb F_2)$ connective, while
$H\mathbb F_2$ is in the heart. This explains the intended source of the
hypotheses. The long-exact-sequence deduction is now formalized above, but
the current fixed foundation still supplies neither sphere connectivity
nor a t-structure detecting its bounds from homotopy groups.
No new project axiom is introduced here.

\begin{definition}[Mixed pairings on the existing Adams layers]
  \label{def:h6-mixed-layer-pairings}
  \lean{KIP126.Classical.Adams.adamsSphereLayerPairingTriangleIso,KIP126.Classical.Adams.adamsSphereLayerPairingIso,KIP126.Classical.Adams.adamsSphereTowerLayerPairingTriangleIso,KIP126.Classical.Adams.adamsSphereTowerLayerPairingIso}
  \uses{prop:h6-existing-tower-pairing-left-transition,prop:h6-existing-tower-pairing-right-transition,def:h6-fiber-tower}
  \leanok
  Let $L_s$ denote the existing layer of the sphere tower, with triangle
  $T_{s+1}S\to T_sS\xrightarrow{q_s}L_s\xrightarrow{\delta_s}\Sigma T_{s+1}S$.
  Under right tensor exactness and monoidal additivity, complete the proved
  first-input tower square to an isomorphism of distinguished triangles.
  Its third component gives a chosen isomorphism
  $\ell_{s,t}:L_s\wedge T_tS\simeq L_{t+s}$.
  With left tensor exactness and the two-factor inclusion condition as well,
  the second-input square similarly gives
  $r_{s,t}:T_sS\wedge L_t\simeq L_{t+s}$.
  These use the existing integer-indexed layers at natural filtration
  indices, not a replacement layer object. Completion is a theorem of the
  pretriangulated interface; the choices are not new project axioms.
  Neither uniqueness nor coherence of these choices is asserted.
\end{definition}

\begin{proposition}[The mixed pairings preserve inclusions and boundaries]
  \label{prop:h6-mixed-layer-pairing-boundaries}
  \lean{KIP126.Classical.Adams.adamsSphereLayerPairingTriangleIso_hom₁,KIP126.Classical.Adams.adamsSphereLayerPairingTriangleIso_hom₂,KIP126.Classical.Adams.adamsSphereLayerPairingIso_ι,KIP126.Classical.Adams.adamsSphereLayerPairingIso_δ,KIP126.Classical.Adams.adamsSphereTowerLayerPairingTriangleIso_hom₁,KIP126.Classical.Adams.adamsSphereTowerLayerPairingTriangleIso_hom₂,KIP126.Classical.Adams.adamsSphereTowerLayerPairingIso_ι,KIP126.Classical.Adams.adamsSphereTowerLayerPairingIso_δ,KIP126.Classical.Adams.adamsSphereMixedLayerPairings_agree}
  \uses{def:h6-mixed-layer-pairings}
  \leanok
  With the hypotheses appropriate to each construction, the first two
  components are the previously constructed tower pairings, and
  \[
    (q_s\wedge\mathrm{id})\ell_{s,t}
      =\mu_{s,t}q_{t+s}
      =(\mathrm{id}\wedge q_t)r_{s,t}.
  \]
  Write $c^R$ and $c^L$ for the suspension comparisons of the right and
  left tensor functors at the relevant tower objects. The actual connecting
  maps satisfy
  \[
    \ell_{s,t}\delta_{t+s}
      =(\delta_s\wedge\mathrm{id})c^R\Sigma\mu_{s+1,t},
    \qquad
    r_{s,t}\delta_{t+s}
      =(\mathrm{id}\wedge\delta_t)c^L\Sigma\mu^+_{s,t}.
  \]
  The successor reindexing in $\mu^+$ is the one already defined above.
  The suspension comparisons are retained explicitly, rather than silently
  replacing them by identities or prescribing a sign convention.
\end{proposition}
\begin{proof}
  The first two component formulas come from the triangle-completion
  construction. Its second square proves each $q$ identity, and its third
  square proves each $\delta$ identity. The common-source equality follows
  since both maps equal $\mu_{s,t}q_{t+s}$.
  \leanok
\end{proof}

\begin{definition}[Coefficient-induced multiplication on actual layers]
  \label{def:h6-coefficient-layer-product}
  \lean{KIP126.StableHomotopy.Cohomology.mod2CoefficientPairing,KIP126.Classical.Adams.adamsSphereLayerProduct}
  \uses{def:h6-ring-cooperations,def:h6-existing-tower-pairing,prop:h6-tower-layer-comparison}
  \leanok
  Give the stable category a braiding and the specified $H$ its associative
  unital multiplication $m_H$. For $f:X\wedge Y\to Z$, define
  \[
    c_f:(H\wedge X)\wedge(H\wedge Y)
       \longrightarrow(H\wedge H)\wedge(X\wedge Y)
       \xrightarrow{m_H\wedge f}H\wedge Z,
  \]
  where the first map uses the braiding and associators to exchange the
  middle factors. Write $e_s:L_s\simeq H\wedge T_sS$ for the existing layer
  comparisons. Under the right-tensor structures used to construct $\mu$,
  define the actual layer product
  $P_{s,t}=(e_s\wedge e_t)c_{\mu_{s,t}}e_{t+s}^{-1}$.
  This is a concrete map $L_s\wedge L_t\to L_{t+s}$ using the supplied
  coefficient multiplication, not a triangle-completion choice or a new
  spectral-sequence multiplication axiom. No associativity, boundary
  compatibility or Lin comparison for this product is asserted here.
\end{definition}

\begin{proposition}[The coefficient product preserves tower images]
  \label{prop:h6-coefficient-layer-product-unit}
  \lean{KIP126.StableHomotopy.Cohomology.mod2Unit_tensor_mul,KIP126.StableHomotopy.Cohomology.mod2CoefficientPairing_unit,KIP126.Classical.Adams.adamsSphereLayerProduct_comparison,KIP126.Classical.Adams.adamsSphereLayerProduct_ι}
  \uses{def:h6-coefficient-layer-product}
  \leanok
  Writing $u_X:X\to H\wedge X$ for the actual Adams unit, one has
  $(u_X\wedge u_Y)c_f=fu_Z$. The defined layer product satisfies its
  comparison formula $P_{s,t}e_{t+s}=(e_s\wedge e_t)c_{\mu_{s,t}}$ and
  \[
    (q_s\wedge q_t)P_{s,t}=\mu_{s,t}q_{t+s}.
  \]
  In particular this product agrees with the prescribed tower pairing on
  pairs of tower-to-layer images. This does not identify its one-sided
  restrictions with the separately chosen mixed pairings above on all
  inputs, nor supply the two-term boundary formula.
\end{proposition}
\begin{proof}
  Naturality of the middle-factor exchange moves the two units onto the
  coefficient factors. The ring unit law and monoidal unitor coherence
  combine them into one unit. For the layer statement, postcompose with
  $e_{t+s}$, apply $q_se_s=u_{T_sS}$ and the coefficient-unit formula,
  then cancel the comparison isomorphism.
  \leanok
\end{proof}

\begin{proposition}[One-unit restrictions of the coefficient layer product]
  \label{prop:h6-coefficient-layer-product-restrictions}
  \lean{KIP126.StableHomotopy.Cohomology.mod2CoefficientPairing_unit_right,KIP126.StableHomotopy.Cohomology.mod2CoefficientPairing_unit_left,KIP126.Classical.Adams.adamsSphereLayerProduct_ι_right_comparison,KIP126.Classical.Adams.adamsSphereLayerProduct_ι_left_comparison}
  \uses{def:h6-coefficient-layer-product,prop:h6-coefficient-layer-product-unit}
  \leanok
  With the same explicit braiding and coefficient ring structure, the
  one-unit restrictions are
  \[
    (\mathrm{id}_{H\wedge X}\wedge u_Y)c_f
       =\alpha_{H,X,Y}(H\wedge f),
  \]
  \[
    (u_X\wedge\mathrm{id}_{H\wedge Y})c_f
       =\alpha_{X,H,Y}^{-1}(\beta_{X,H}\wedge\mathrm{id}_Y)
          \alpha_{H,X,Y}(H\wedge f).
  \]
  Here composition is written in execution order. In particular, after
  using $e_s^{-1}$ on the layer input and $e_{t+s}$ on the output,
  $(\mathrm{id}_{L_s}\wedge q_t)P_{s,t}$ has the first formula with
  $X=T_sS$, $Y=T_tS$, $f=\mu_{s,t}$; using $e_t^{-1}$ instead gives
  the second formula for $(q_s\wedge\mathrm{id}_{L_t})P_{s,t}$.
  The right-hand sides no longer involve the coefficient multiplication.
  These are formulas for the actual coefficient-induced product, not an
  assertion that its restrictions equal the independently chosen mixed
  triangle completions, and not a boundary product formula.
\end{proposition}
\begin{proof}
  Naturality of the middle-factor exchange moves the inserted unit to the
  corresponding coefficient input. Apply the left or right ring unit law
  and monoidal coherence, retaining the braiding in the left-input case.
  Transport through the layer comparisons using $q_se_s=u_{T_sS}$ and
  cancellation of each inverse comparison with its forward map.
  \leanok
\end{proof}

\begin{definition}[Right-tensor suspension compatibility]
  \label{def:h6-right-tensor-suspension-compatibility}
  \lean{KIP126.StableHomotopy.RightTensorSuspensionCompatibility}
  \uses{def:stable-context}
  \leanok
  For the supplied right-tensor shift structures, write
  $c_X(B):(\Sigma B)\wedge X\to\Sigma(B\wedge X)$ for their comparison
  at suspension one. Require naturality in the right variable and
  compatibility with reassociation:
  \[
    (\Sigma B\wedge f)c_Y(B)=c_X(B)\Sigma(B\wedge f),
  \]
  \[
    \alpha_{\Sigma B,X,Y}c_{X\wedge Y}(B)
       =(c_X(B)\wedge\mathrm{id}_Y)c_Y(B\wedge X)
         \Sigma\alpha_{B,X,Y}.
  \]
  Composition is in execution order. This predicate concerns only the
  ambient tensor and shift structures; it neither mentions Adams pages
  nor assumes a differential formula. Individual shift structures on
  the right-tensor functors do not include these family compatibilities.
  No witness for the fixed foundation is supplied here.
\end{definition}

\begin{proposition}[Tensor pairing with a suspended-output map]
  \label{prop:h6-tensor-suspended-pairing}
  \lean{KIP126.StableHomotopy.rightTensorSuspension_pairing,KIP126.Classical.Adams.adamsFiberTensorIso_pairing_δ}
  \uses{def:h6-right-tensor-suspension-compatibility,prop:h6-fiber-smash-step}
  \leanok
  For any $d:H\to\Sigma B$ and $f:X\wedge Y\to Z$, the preceding
  compatibility implies
  \[
    \alpha_{H,X,Y}(H\wedge f)(d\wedge\mathrm{id}_Z)c_Z(B)
      =((d\wedge\mathrm{id}_X)c_X(B)\wedge\mathrm{id}_Y)
        c_Y(B\wedge X)\Sigma(\alpha_{B,X,Y}(B\wedge f)).
  \]
  Put $B=\operatorname{fib}(S\to H)$ and let $d$ be its actual fiber
  connecting map. Write $F_X:\operatorname{fib}(u_X)\simeq B\wedge X$
  for the existing fiber comparison and $D_X:H\wedge X\to
  \Sigma\operatorname{fib}(u_X)$ for the actual connecting map. Then
  \[
    \alpha_{H,X,Y}(H\wedge f)D_Z
       =(D_X\wedge\mathrm{id}_Y)c_Y(\operatorname{fib}(u_X))
         \Sigma((F_X\wedge\mathrm{id}_Y)\alpha_{B,X,Y}
                    (B\wedge f)F_Z^{-1}).
  \]
\end{proposition}
\begin{proof}
  Move $d$ past $f$ and the associator by naturality, and apply the two
  suspension compatibility identities. For the second formula use
  $D_X=(d\wedge\mathrm{id}_X)c_X(B)\Sigma F_X^{-1}$, naturality of
  $c_Y$ in its left variable, and cancellation of $\Sigma F_X^{-1}$
  with $\Sigma F_X$.
  \leanok
\end{proof}

\begin{proposition}[Successor tower pairing and its fiber boundary]
  \label{prop:h6-tower-pairing-fiber-boundary}
  \lean{KIP126.Classical.Adams.adamsTowerSpherePairingIso_succ_comparison,KIP126.Classical.Adams.adamsTowerSpherePairingIso_δ_left}
  \uses{def:h6-existing-tower-pairing,prop:h6-tensor-suspended-pairing}
  \leanok
  Write $X=T_sS$, $Y=T_tS$, $Z=T_{t+s}S$. The actual tower pairing has
  successor formula
  \[
    \mu_{s+1,t}F_Z=(F_X\wedge\mathrm{id}_Y)
       \alpha_{B,X,Y}(B\wedge\mu_{s,t}).
  \]
  This identity needs only the previously supplied exact right-tensor
  structures. With right-tensor suspension compatibility it implies
  \[
    \alpha_{H,X,Y}(H\wedge\mu_{s,t})D_Z
       =(D_X\wedge\mathrm{id}_Y)c_Y(T_{s+1}S)\Sigma\mu_{s+1,t}.
  \]
\end{proposition}
\begin{proof}
  Compare with the iterated-smash tower, expand its successor pairing,
  and use associator naturality. Substitute this identity into the
  preceding fiber boundary formula. No new triangle-completion choice
  is made.
  \leanok
\end{proof}

\begin{proposition}[Coefficient-side boundary compatibility]
  \label{prop:h6-coefficient-side-boundary}
  \lean{KIP126.Classical.Adams.adamsSphereLayerProduct_ι_right_δ_comparison,KIP126.Classical.Adams.adamsSphereLayerProduct_ι_right_δ}
  \uses{prop:h6-tower-pairing-fiber-boundary,prop:h6-coefficient-layer-product-restrictions}
  \leanok
  Under the preceding suspension compatibility and biadditivity of the
  tensor product, the actual coefficient-induced product satisfies
  \[
    (\mathrm{id}_{L_s}\wedge q_t)P_{s,t}\delta_{t+s}
       =(\delta_s\wedge\mathrm{id}_{T_tS})c_{T_tS}(T_{s+1}S)
          \Sigma\mu_{s+1,t}.
  \]
  This uses the given coefficient product, rather than replacing it by
  an independently selected triangle completion. It is only a one-sided
  formula, not the full two-term Leibniz rule.
\end{proposition}
\begin{proof}
  First precompose the layer input with $e_s^{-1}$ and use its one-unit
  restriction formula. The existing layer comparison gives
  $\delta_{t+s}=-e_{t+s}D_Z$, so the preceding tower boundary identity
  computes this composite, including its minus sign. Since
  $e_s^{-1}\delta_s=-D_X$, biadditivity cancels the matching sign on the
  source side. Cancel the isomorphism $e_s^{-1}\wedge\mathrm{id}$.
  \leanok
\end{proof}

\begin{definition}[Braiding and tensor suspension comparisons]
  \label{def:h6-tensor-suspension-braiding}
  \lean{KIP126.StableHomotopy.TensorSuspensionBraidingCompatibility}
  \uses{def:h6-right-tensor-suspension-compatibility}
  \leanok
  In addition to the right-tensor shift structures, supply left-tensor
  shift structures, and write their suspension-one comparisons as
  $c_X^R(B)$ and $c_X^L(B)$. The braiding compatibility condition is
  \[
    \beta_{X,\Sigma B}c_X^R(B)=c_X^L(B)\Sigma\beta_{X,B}.
  \]
  This concerns the ambient structures. It does not assert that the
  tower pairing is symmetric, and no fixed-foundation witness is provided.
\end{definition}

\begin{proposition}[Transporting a boundary past a braiding]
  \label{prop:h6-braided-boundary-transport}
  \lean{KIP126.StableHomotopy.braidedTensorPairing_swap,KIP126.StableHomotopy.braidedSuspension_transport}
  \uses{def:h6-tensor-suspension-braiding}
  \leanok
  For $f:X\wedge Y\to Z$, the coefficient-side exchange equals
  \[
    \alpha_{X,H,Y}^{-1}(\beta_{X,H}\wedge\mathrm{id}_Y)
       \alpha_{H,X,Y}(H\wedge f)
      =\beta_{X,H\wedge Y}\alpha_{H,Y,X}
         (H\wedge(\beta_{X,Y}^{-1}f)).
  \]
  For $d:Y\to\Sigma B$ and $k:B\wedge X\to Z$, the suspension
  compatibility gives
  \[
    \beta_{X,Y}(d\wedge\mathrm{id}_X)c_X^R(B)\Sigma k
       =(X\wedge d)c_X^L(B)\Sigma(\beta_{X,B}k).
  \]
  In particular the braiding on the final fiber-level map must be retained.
\end{proposition}
\begin{proof}
  The first identity is the braiding hexagon followed by cancellation of
  the inverse braiding. For the second, apply braiding naturality, the
  displayed suspension compatibility, and functoriality of suspension.
  \leanok
\end{proof}

\begin{definition}[The braided right fiber lift]
  \label{def:h6-braided-right-fiber-lift}
  \lean{KIP126.Classical.Adams.adamsFiberTensorPairingRight,KIP126.Classical.Adams.adamsTowerSpherePairingBoundaryRight}
  \uses{prop:h6-fiber-smash-step,def:h6-existing-tower-pairing}
  \leanok
  Put $B=\operatorname{fib}(S\to H)$, $B_Y=\operatorname{fib}(u_Y)$,
  and retain the actual comparisons $F_Y:B_Y\simeq B\wedge Y$.
  For $f:X\wedge Y\to Z$ define
  \[
    g_f=\beta_{X,B_Y}(F_Y\wedge\mathrm{id}_X)\alpha_{B,Y,X}
        (B\wedge(\beta_{X,Y}^{-1}f))F_Z^{-1}.
  \]
  Thus $g_f:X\wedge B_Y\to B_Z$. For $f=\mu_{s,t}$ this defines
  $g_{s,t}:T_sS\wedge T_{t+1}S\to T_{t+s+1}S$ on the existing tower.
  It is not a replacement tower pairing and is not identified by definition
  with the ordered successor $\mu^+_{s,t}$.
\end{definition}

\begin{proposition}[The actual second-input boundary]
  \label{prop:h6-braided-right-boundary}
  \lean{KIP126.Classical.Adams.adamsFiberTensorPairingRight_δ,KIP126.Classical.Adams.adamsTowerSpherePairingBoundaryRight_δ,KIP126.Classical.Adams.adamsSphereLayerProduct_ι_left_δ_comparison,KIP126.Classical.Adams.adamsSphereLayerProduct_ι_left_δ}
  \uses{def:h6-braided-right-fiber-lift,prop:h6-braided-boundary-transport,prop:h6-tensor-suspended-pairing,prop:h6-coefficient-layer-product-restrictions}
  \leanok
  Assume right-tensor suspension compatibility and braiding compatibility
  of the two tensor suspension structures. With $D_Y:H\wedge Y\to
  \Sigma B_Y$ as above, the second-input coefficient boundary is
  \[
    \alpha_{X,H,Y}^{-1}(\beta_{X,H}\wedge\mathrm{id}_Y)
       \alpha_{H,X,Y}(H\wedge f)D_Z
       =(X\wedge D_Y)c_X^L(B_Y)\Sigma g_f.
  \]
  For the specified coefficient-induced layer multiplication, and
  assuming biadditivity of the tensor product, this becomes
  \[
    (q_s\wedge\mathrm{id}_{L_t})P_{s,t}\delta_{t+s}
       =(T_sS\wedge\delta_t)c_{T_sS}^L(T_{t+1}S)\Sigma g_{s,t}.
  \]
  There is also the coefficient-coordinate formula, precomposed with
  $\mathrm{id}_{T_sS}\wedge e_t^{-1}$, with right side
  $-(T_sS\wedge D_{T_tS})c_{T_sS}^L(T_{t+1}S)\Sigma g_{s,t}$.
\end{proposition}
\begin{proof}
  Exchange the inputs using the first transport identity, apply the
  first-input fiber boundary calculation to $\beta_{X,Y}^{-1}f$, and
  use the second transport identity. For the layer statement use its
  existing comparison and its negative connecting-map formula, then
  cancel the layer-input isomorphism. The product $P$ is unchanged.
  \leanok
\end{proof}

\begin{proposition}[The ordered-successor obstruction]
  \label{prop:h6-ordered-right-boundary-obstruction}
  \lean{KIP126.Classical.Adams.adamsTowerSpherePairingBoundaryRight_obstruction,KIP126.Classical.Adams.adamsSphereLayerProduct_ι_left_δ_ordered_iff}
  \uses{prop:h6-braided-right-boundary,prop:h6-existing-tower-pairing-right-transition}
  \leanok
  Write $a_{s,t}=(T_sS\wedge\delta_t)c_{T_sS}^L(T_{t+1}S)$.
  Under the preceding hypotheses, the morphism-level boundary formula with the
  ordered successor holds exactly when the following concrete composite
  vanishes:
  \[
    (q_s\wedge\mathrm{id})P_{s,t}\delta_{t+s}
         =a_{s,t}\Sigma\mu^+_{s,t}
    \quad\Longleftrightarrow\quad
    a_{s,t}\Sigma(g_{s,t}-\mu^+_{s,t})=0.
  \]
  The same difference criterion compares the two coefficient-side
  composites with $D_{T_tS}$ in place of $\delta_t$.
  Neither vanishing nor equality of the two lifts is asserted. Vanishing
  of this composite is the precise condition for this morphism identity;
  equality of lifts would be a stronger sufficient condition. Necessity
  of this global condition for page-level multiplicativity is not asserted.
\end{proposition}
\begin{proof}
  Substitute the proved boundary formula with $g_{s,t}$. Additivity of
  suspension and composition identifies the difference of the resulting
  two composites with the displayed obstruction.
  \leanok
\end{proof}

\begin{proposition}[Connecting maps on representatives and products of tower images]
  \label{prop:h6-product-tower-image-cycles}
  \lean{KIP126.Classical.Adams.adamsK_eq_zero_iff_comp_δ_eq_zero,KIP126.Classical.Adams.adamsCycles_mem_of_comp_δ_eq_zero,KIP126.Classical.Adams.adamsCycleSubmodule_top_mem_of_comp_δ_eq_zero,KIP126.Classical.Adams.adamsSphereLayerProduct_comp_δ_eq_zero,KIP126.Classical.Adams.adamsSphereLayerProduct_mem_cycles}
  \uses{prop:h6-coefficient-layer-product-unit}
  \leanok
  For a homotopy representative $x$ in an actual layer, $Kx=0$ if and
  only if $x\delta=0$. Such a representative lies in every finite cycle
  submodule and hence, as a $Z_2$ representative, in the actual $Z_\infty$
  intersection used by internal SSData. For maps $x:A\to L_s$ and $y:B\to L_t$ with
  $x\delta_s=y\delta_t=0$, the specified product satisfies
  $(x\wedge y)P_{s,t}\delta_{t+s}=0$. After any specified sphere
  precomposition this is an element of every actual $Z_r$.
  The product assertion assumes only the ambient structures used to
  construct $P$, not the additional suspension compatibility predicates.
  No nonvanishing or graded sphere-product construction is asserted.
\end{proposition}
\begin{proof}
  Exactness identifies zero connecting image with a lift through $q$.
  Lift both inputs, use $(q_s\wedge q_t)P=\mu q_{t+s}$, and use
  $q_{t+s}\delta_{t+s}=0$. The zero connecting image lifts through every
  tower transition by the zero element, giving membership in every $Z_r$.
  \leanok
\end{proof}

\begin{proposition}[The second-input formula evaluated on representatives]
  \label{prop:h6-right-boundary-representatives}
  \lean{KIP126.Classical.Adams.adamsSphereLayerProduct_ι_left_δ_representatives,KIP126.Classical.Adams.adamsSphereLayerProduct_ι_left_δ_representatives_ordered_iff}
  \uses{prop:h6-braided-right-boundary,prop:h6-ordered-right-boundary-obstruction}
  \leanok
  Under the hypotheses of the second-input boundary formula, for
  $x:A\to T_sS$ and $y:B\to L_t$ put
  $a(x,y)=(x\wedge(y\delta_t))c^L_{T_sS}(T_{t+1}S)$. Then
  \[
    ((xq_s)\wedge y)P_{s,t}\delta_{t+s}=a(x,y)\Sigma g_{s,t}.
  \]
  Replacing $g_{s,t}$ here by $\mu^+_{s,t}$ is equivalent to
  $a(x,y)\Sigma(g_{s,t}-\mu^+_{s,t})=0$. This tests the particular
  decomposable representatives; it does not require the global morphism
  obstruction to vanish. Taking $A,B$ to be spheres supplies the inputs
  relevant to homotopy products, still before page quotients.
\end{proposition}
\begin{proof}
  Precompose the layer formula with $x\wedge y$, use tensor functoriality,
  and subtract the ordered-successor expression.
  \leanok
\end{proof}

\paragraph{Remaining multiplicativity boundary.}
These equalities are in the homotopy-category interface. They do not supply
coherent tower-level data or the full two-term boundary product formula
needed for Leibniz. The mixed pairings above supply individual one-sided
boundary formulas, while the coefficient construction supplies a concrete
layer-layer product with explicit one-unit restrictions. Under the stated
ambient suspension compatibility, that specified product now satisfies
the first-input boundary formula after restriction of the second input to
the tower image. The other side now has a formula with the explicit braided
lift. Global comparison with the ordered successor would require the
morphism-level vanishing condition above, but this is only a stronger
possible route, not a proved prerequisite for page products. On represented
inputs the exact condition is the corresponding precomposed obstruction;
passing to page quotients may weaken it further. Products of two tower
images are already cycles on every finite page by exactness. The full
two-term boundary formula for general cycles, coherent
gluing, passage to spectral-sequence pages, and identification with the
Lin product remain unproved. Equality
with the independently chosen mixed triangle completions is not claimed.
Dugger, \emph{Multiplicative structures on homotopy spectral sequences, Part I},
\url{https://arxiv.org/pdf/math/0305173}, Sections 1 and 7, explains why pairings
commuting only up to homotopy need not induce spectral-sequence pairings and
gives a counterexample. This is a design constraint, not an additional Lean
axiom or a proved project theorem. The fixed two-factor witness, tensor
structures and their suspension compatibilities, the full boundary product
formula, and identification with Lin's product remain separate obligations.

\begin{proposition}[The tower is killed by coefficient homology]
  \label{prop:h6-tower-homology-zero}
  \lean{KIP126.StableHomotopy.Cohomology.mod2FreeAction,KIP126.StableHomotopy.Cohomology.mod2FreeAction_unit,KIP126.StableHomotopy.Cohomology.mod2_adamsTowerStep_homology_eq_zero,KIP126.StableHomotopy.Cohomology.mod2_adamsTowerMap_eq_zero,KIP126.StableHomotopy.Cohomology.mod2_homology_eq_zero_of_tower_lift}
  \uses{def:h6-ring-cooperations,def:h6-fiber-tower}
  \leanok
  Suppose $H$ has the specified associative unital multiplication and
  $H\wedge-$ preserves zero morphisms. Multiplication on the first two
  factors retracts $H\wedge u_X$. Consequently $H\wedge(T_{s+1}\to T_s)$
  is the zero map; the same holds for every strictly descending composite
  of natural-indexed tower maps. If a map $f:X\to Y$ is supplied with a
  factorization through any positive stage of the actual tower of $Y$,
  then $H_nf=0$ in every degree $n$.
\end{proposition}
\begin{proof}
  Use the right unit identity for the multiplication, associator naturality,
  and the triangle identity to construct the retraction. The fiber inclusion
  composed with the unit is zero. Smash this equality with $H$, then compose
  with the retraction. Factor a positive-length tower map through one such
  step, and postcompose representatives of homotopy classes to obtain the
  homology statement.
  \leanok
\end{proof}

These are generic theorems with explicit structural arguments. The selected
\texttt{standardFoundation} does not yet supply the multiplication and
tensor-zero-preservation witnesses used in the second proposition. They are
not added as new global axioms. Neither proposition provides the Milnor
cooperation coordinates, a concrete model of classical spectra, or permanence
of the specified square.

\begin{definition}[Quotient pages from the Adams tower]
  \label{def:h6-tower-pages}
  \lean{KIP126.Classical.Adams.adamsE1}
  \lean{KIP126.Classical.Adams.adamsI}
  \lean{KIP126.Classical.Adams.adamsJ}
  \lean{KIP126.Classical.Adams.adamsK}
  \lean{KIP126.Classical.Adams.adamsCycles}
  \lean{KIP126.Classical.Adams.adamsBoundaries}
  \lean{KIP126.Classical.Adams.adamsCycleBoundaries}
  \lean{KIP126.Classical.Adams.adamsPage}
  \uses{def:h6-fiber-tower,thm:stable-homotopy-long-exact}
  \notready
  Write $E_1^{s,t}=\pi_{t-s}Q_s$.  Let $j$ be induced by $T_s\to Q_s$,
  let $k:E_1^{s,t}\to\pi_{t-s-1}T_{s+1}$ be the connecting map, and
  write $i_{a,b}$ for the map on a fixed homotopy group induced by
  $T_b\to T_a$.  For $r\ge1$ define
  \[
    Z_r^{s,t}=k^{-1}(\operatorname{im}i_{s+1,s+r}),\qquad
    B_r^{s,t}=j(\ker i_{s-r+1,s}),\qquad
    E_r^{s,t}=Z_r^{s,t}/(Z_r^{s,t}\cap B_r^{s,t}).
  \]
  These are actual subgroups and quotients of the homotopy groups of the
  constructed tower.  The exact-couple identities imply $B_r\subseteq Z_r$.
\end{definition}

\begin{proposition}[The first quotient page is coefficient homology]
  \label{prop:h6-first-page-homology}
  \lean{KIP126.Classical.Adams.adamsCycles_one,KIP126.Classical.Adams.adamsBoundaries_one,KIP126.Classical.Adams.adamsPageOneHomologyEquiv,KIP126.Classical.Adams.adamsE1HomologyEquiv_J,KIP126.Classical.Adams.adamsPageOneHomologyEquiv_mkQ}
  \uses{def:h6-tower-pages,prop:h6-tower-layer-comparison}
  \leanok
  At page one, $Z_1=E_1$ and $B_1=0$. For $s\ge0$ the constructed quotient
  page is therefore integer-linearly isomorphic to
  $\pi_{t-s}(H\wedge T_s)$, which is $H_{t-s}(T_s;\mathbb F_2)$ when
  $H=H\mathbb F_2$. The isomorphism sends a quotient representative to
  its postcomposition with the layer comparison and sends $j(x)$ to the
  homotopy class induced by the Adams unit. This is the tower's actual
  first quotient page, not the out-of-range first page of the internal
  spectral sequence whose starting page is two.
\end{proposition}
\begin{proof}
  The ordered tower map from a stage to itself is the identity. Its image
  is the full group and its kernel is zero. Remove the zero quotient and
  compose with the homotopy-group isomorphism induced by the layer comparison.
  \leanok
\end{proof}

\begin{definition}[Page differential from a tower lift]
  \label{def:h6-tower-differential}
  \lean{KIP126.Classical.Adams.adamsCycleLift}
  \lean{KIP126.Classical.Adams.adamsDifferentialValue}
  \lean{KIP126.Classical.Adams.adamsDifferential}
  \lean{KIP126.Classical.Adams.adamsPageComplex}
  \uses{def:h6-tower-pages}
  \notready
  For $x\in Z_r^{s,t}$, choose $y\in\pi_{t-s-1}T_{s+r}$ with
  $i_{s+1,s+r}(y)=k(x)$.  Define
  \[
    d_r([x])=[j(y)]\in E_r^{s+r,t+r-1}.
  \]
  Thus the map is specified on representatives before descending to the
  quotient.  Its additivity, independence of the lift and representative,
  and square-zero law are established in Proposition~\ref{prop:h6-tower-differential-laws}; they are not inputs selecting a map.
\end{definition}

\begin{proposition}[Well-definedness and square-zero of the tower differential]
  \label{prop:h6-tower-differential-laws}
  \lean{KIP126.Classical.Adams.adamsDifferentialValue_eq_of_lift}
  \lean{KIP126.Classical.Adams.adamsDifferentialValue_add}
  \lean{KIP126.Classical.Adams.adamsDifferentialValue_smul}
  \lean{KIP126.Classical.Adams.adamsBoundaries_le_differential_ker}
  \lean{KIP126.Classical.Adams.adamsPageD_comp}
  \uses{def:h6-tower-differential}
  \leanok
  The differential is independent of the chosen lift, is additive and
  integer-linear, vanishes on cycle boundaries, and satisfies $d_r^2=0$.
\end{proposition}
\begin{proof}
  Two lifts differ by an element in the kernel of the tower composite,
  so their layer images differ by a boundary. Use sums and scalar multiples
  of lifts to establish linearity. Exactness gives $kj=0$; hence every
  layer image has zero next differential, proving both boundary vanishing
  and square-zero.
  \leanok
\end{proof}

\begin{definition}[Long-layer representatives of the existing tower]
  \label{def:h6-tower-long-layer}
  \lean{KIP126.StableHomotopy.cofiberFactorizationMap,KIP126.Classical.Adams.adamsLongLayer,KIP126.Classical.Adams.adamsLongLayerProjection,KIP126.Classical.Adams.adamsLongLayerToE1,KIP126.Classical.Adams.adamsLongLayerK}
  \uses{def:h6-tower-pages}
  \leanok
  For $r\ge1$ and integer $s$, put
  $Q_s^{(r)}=\operatorname{cofib}(T_{s+r}\to T_s)$.
  The factorization through $T_{s+1}$ and the identity on $T_s$ give a
  specified map $p_r:Q_s^{(r)}\to Q_s$. Let $k_r$ be the connecting
  homomorphism of the long cofiber triangle. These use the actual tower
  and its existing cofiber choices, without any assumed composition law
  for those choices or replacement of its spectral sequence.
\end{definition}

\begin{proposition}[All cycles and pages have long-layer representatives]
  \label{prop:h6-tower-long-layer-cycles}
  \lean{KIP126.StableHomotopy.cofiberFactorizationMap_ι,KIP126.StableHomotopy.cofiberFactorizationMap_δ,KIP126.StableHomotopy.cofiberFactorizationMap_connecting,KIP126.StableHomotopy.cofiberFactorizationMap_image_iff,KIP126.Classical.Adams.adamsLongLayerK_lift,KIP126.Classical.Adams.adamsCycles_mem_iff_longLayer,KIP126.Classical.Adams.adamsLongLayerToE1_mem_cycles,KIP126.Classical.Adams.adamsCycles_eq_range_longLayer,KIP126.Classical.Adams.adamsLongLayerToCycles,KIP126.Classical.Adams.adamsLongLayerToPage,KIP126.Classical.Adams.adamsLongLayerToCycles_surjective,KIP126.Classical.Adams.adamsLongLayerToPage_surjective}
  \uses{def:h6-tower-long-layer,thm:stable-homotopy-long-exact}
  \leanok
  With $n=t-s$, one has an equality of the existing submodules
  \[
    Z_r^{s,t}=\operatorname{im}\bigl(\pi_nQ_s^{(r)}
        \xrightarrow{(p_r)_*}\pi_nQ_s\bigr).
  \]
  Hence the induced linear maps to $Z_r^{s,t}$ and to the already
  constructed $E_r^{s,t}$ are surjective. Naturality gives
  $k(p_rz)=i_{s+1,s+r}(k_rz)$, so these statements cover general cycles,
  not merely representatives with zero connecting image.
\end{proposition}
\begin{proof}
  More generally, for a factorization $A\xrightarrow{a}B\xrightarrow{g}D$
  with composite $f$, a class $x\in\pi_n\operatorname{cofib}(g)$ is in
  the image of $\pi_n\operatorname{cofib}(f)$ if and only if its
  connecting image is in the image of $\pi_{n-1}a$.
  For the reverse direction choose a lift $y$ of that connecting image.
  Exactness gives $f_*y=0$, hence a long-cofiber class $z$ with boundary
  $y$. Then $x-p_*z$ has zero boundary and comes from $\pi_nD$; add
  the corresponding long-cofiber image to $z$. The forward direction
  is boundary naturality. Apply this to the tower factorization and
  then the surjective quotient map. No octahedral or tensor hypothesis
  is needed.
  \leanok
\end{proof}

\begin{proposition}[The existing differential from the long-layer boundary]
  \label{prop:h6-tower-long-layer-differential}
  \lean{KIP126.Classical.Adams.adamsDifferential_of_longLayer,KIP126.Classical.Adams.adamsDifferential_longLayerToPage}
  \uses{prop:h6-tower-long-layer-cycles,prop:h6-tower-differential-laws}
  \leanok
  For any $z\in\pi_{t-s}Q_s^{(r)}$, the actual quotient differential is
  \[
    d_r([p_rz])=[j(k_rz)]\in E_r^{s+r,t+r-1}.
  \]
  In particular the long-cofiber formula computes $d_r$ on every page
  class, not a newly chosen differential.
\end{proposition}
\begin{proof}
  Boundary naturality makes $k_rz$ a valid lift in the existing
  definition of $d_r$. Apply its proved independence-of-lift theorem;
  the homotopy degrees agree since $(t+r-1)-(s+r)=t-s-1$.
  \leanok
\end{proof}

\paragraph{Boundary of this representative construction.}
These results supply all-page relative representatives needed for a
possible multiplicativity proof. They do not construct pairings between
the long layers or prove a relative product boundary formula. The conditional
quotient-descent construction below handles the subsequent algebraic step.
Neither Leibniz nor compatibility with the Lin multiplication follows from
the surjectivity statement alone.

\begin{proposition}[Long-layer representatives on the internal SSData page]
  \label{prop:h6-tower-long-layer-internal}
  \lean{KIP126.Classical.Adams.adamsLongLayerToInternalPage,KIP126.Classical.Adams.adamsLongLayerToInternalPage_comparison,KIP126.Classical.Adams.adamsLongLayerToInternalPage_surjective,KIP126.Classical.Adams.adamsTowerInternalD_longLayer}
  \uses{prop:h6-tower-long-layer-differential,prop:h6-tower-internal-laws}
  \leanok
  At internal stage $n$ (classical page $n+2$), the long-layer map also
  surjects onto the existing SSData page, through the inverse of its
  representative-preserving quotient comparison. Under the target
  comparison, its actual internal differential is again $[j(k_{n+2}z)]$.
  This uses KIP126's internal SSData/PreSS, not a Mathlib spectral sequence.
\end{proposition}
\begin{proof}
  Compose the long-layer surjection with the inverse page comparison.
  Cancel that isomorphism for surjectivity and use the already proved
  differential comparison and long-layer formula.
  \leanok
\end{proof}

\begin{definition}[Inputs for long-layer pairing descent]
  \label{def:h6-long-pairing-inputs}
  \lean{KIP126.Classical.Adams.AdamsLongLayerPairing,KIP126.Classical.Adams.AdamsLongLayerPairing.ProjectionCompatible,KIP126.Classical.Adams.AdamsLongLayerPairing.BoundaryCompatible}
  \uses{def:h6-tower-long-layer,def:h6-tower-pages}
  \leanok
  Fix $r\ge1$ and bidegrees $p,q$. A candidate consists of bilinear maps
  $m_1:E_1^p\times E_1^q\to E_1^{p+q}$ and
  $m_r:\pi Q_p^{(r)}\times\pi Q_q^{(r)}\to\pi Q_{p+q}^{(r)}$,
  in the corresponding homotopy degrees. Projection compatibility means
  $p_rm_r(a,b)=m_1(p_ra,p_rb)$.
  Boundary compatibility asks that $m_1(j(u),p_rb)$ lie in $B_r^{p+q}$
  whenever $u$ is killed by the tower map defining $B_r^p$, and analogously
  for $m_1(p_ra,j(v))$. These are predicates on supplied maps, not axioms
  asserting their existence. A sphere application must still identify $m_1$
  with the intended coefficient-induced product.
\end{definition}

\begin{definition}[Specified sphere representative product]
  \label{def:h6-specified-sphere-pairing}
  \lean{KIP126.StableHomotopy.sphereTensorIsoFromRightShift,KIP126.StableHomotopy.tensorHomPairing,KIP126.StableHomotopy.homotopyTensorPairing,KIP126.Classical.Adams.adamsSphereLayerProductOrdered,KIP126.Classical.Adams.adamsSphereE1Product,KIP126.Classical.Adams.adamsSphereLongLayerPairing}
  \uses{def:h6-coefficient-layer-product,def:h6-long-pairing-inputs}
  \leanok
  Assume additive tensor and the given right-tensor shift and exactness
  structures. The right shift comparison, left unitor and shift-addition
  comparison specify $S^a\wedge S^b\simeq S^{a+b}$. For an actual map
  $f:X\wedge Y\to Z$, precompose $(x\wedge y)f$ with its inverse to get
  an integer-bilinear pairing of homotopy groups. Apply this to the
  coefficient-induced layer product, reindexed from $t+s$ to $s+t$, to
  fix $m_1$. Given an actual long-layer map $\mu$, the same construction
  fixes $m_r$ and hence a candidate with this specified $m_1$.
  No graded symmetry, sign coherence, or Lin comparison is asserted.
\end{definition}

\begin{proposition}[Spectrum projection compatibility implies represented compatibility]
  \label{prop:h6-sphere-pairing-projection}
  \lean{KIP126.StableHomotopy.tensorHomPairing_apply,KIP126.StableHomotopy.tensorHomPairing_naturality,KIP126.StableHomotopy.homotopyTensorPairing_naturality,KIP126.Classical.Adams.adamsSphereLongLayerPairing_projection}
  \uses{def:h6-specified-sphere-pairing}
  \leanok
  If $\mu p_r=(p_r\wedge p_r)P$ as spectrum maps, the specified
  representative pairing is projection compatible. This is conditional on
  the spectrum map and its square, not a construction of them for $r>1$.
\end{proposition}
\begin{proof}
  Tensor respects composition; precompose the assumed square with the
  tensor of the two representatives and the specified sphere comparison.
  \leanok
\end{proof}

\begin{proposition}[Conditional pairing on the actual quotient pages]
  \label{prop:h6-long-pairing-descent}
  \lean{KIP126.Classical.Adams.AdamsLongLayerPairing.mem_cycles,KIP126.Classical.Adams.AdamsLongLayerPairing.boundary_left,KIP126.Classical.Adams.AdamsLongLayerPairing.boundary_right,KIP126.Classical.Adams.AdamsLongLayerPairing.onCycles,KIP126.Classical.Adams.AdamsLongLayerPairing.cyclesToPage,KIP126.Classical.Adams.AdamsLongLayerPairing.boundaries_le_left_ker,KIP126.Classical.Adams.AdamsLongLayerPairing.boundaries_le_right_ker,KIP126.Classical.Adams.AdamsLongLayerPairing.onPage,KIP126.Classical.Adams.AdamsLongLayerPairing.onPage_mk,KIP126.Classical.Adams.AdamsLongLayerPairing.onPage_long,KIP126.Classical.Adams.AdamsLongLayerPairing.onPage_unique,KIP126.Classical.Adams.AdamsLongLayerPairing.onInternalPage,KIP126.Classical.Adams.AdamsLongLayerPairing.onInternalPage_comparison,KIP126.Classical.Adams.AdamsLongLayerPairing.onInternalPage_long}
  \uses{def:h6-long-pairing-inputs,prop:h6-tower-long-layer-cycles,prop:h6-tower-long-layer-internal}
  \leanok
  Assume both compatibility predicates. The specified $m_1$ preserves
  $Z_r\times Z_r$ and sends $B_r\times Z_r$ and $Z_r\times B_r$ into
  $B_r$. It therefore induces a bilinear map on the existing quotient
  pages, with $[p_ra]\,[p_rb]=[p_rm_r(a,b)]$.
  This formula determines that map uniquely. For $r=n+2$ transport through
  the existing quotient comparisons gives the same pairing on internal
  SSData pages, preserving the long-layer formula.
  The statement fixes a page and two bidegrees; it does not assert
  associativity, a unit, or multiplicative compatibility between pages.
\end{proposition}
\begin{proof}
  Represent both cycles by long-layer classes and apply projection
  compatibility. Represent a boundary as $j$ of an element killed by the
  appropriate tower composite, and the other cycle by a long-layer class;
  apply the corresponding boundary condition. After projecting the output,
  these give the two kernel inclusions required for bilinear quotient descent.
  Surjectivity of both long-layer-to-page maps proves uniqueness.
  \leanok
\end{proof}

\begin{proposition}[Constructed one-step sphere pairing]
  \label{prop:h6-sphere-pairing-one}
  \lean{KIP126.StableHomotopy.cofiberFactorizationMap_isIso,KIP126.Classical.Adams.adamsLongLayerProjection_one_isIso,KIP126.Classical.Adams.adamsSphereLongLayerOneProduct,KIP126.Classical.Adams.adamsSphereLongLayerOnePairing,KIP126.Classical.Adams.AdamsLongLayerPairing.boundaryCompatible_one,KIP126.Classical.Adams.adamsSphereLongLayerOneProduct_projection,KIP126.Classical.Adams.adamsSphereLongLayerOnePairing_projection,KIP126.Classical.Adams.adamsSpherePageOneProduct,KIP126.Classical.Adams.adamsSpherePageOneProduct_long}
  \uses{prop:h6-sphere-pairing-projection,prop:h6-long-pairing-descent,prop:h6-first-page-homology}
  \leanok
  Under the same ambient assumptions, the one-step projection $p_1$ is
  an isomorphism. Define $\mu_1=(p_1\wedge p_1)Pp_1^{-1}$. Its projection
  square holds, and its represented pairing is boundary compatible because
  $B_1=0$. Thus the existing first quotient page has a constructed bilinear
  product, agreeing with $\mu_1$ on long-layer representatives, without a
  free pairing or compatibility input. This does not prove a derivation law
  for $d_1$, descent to $E_2$, or a witness for the fixed standard foundation.
\end{proposition}
\begin{proof}
  The cofiber comparison has isomorphisms on its first two triangle terms;
  the triangulated five lemma makes the third map an isomorphism as well.
  This does not require that a chosen cofiber map of identities be identity.
  Inverse cancellation proves the projection square. A tower-kernel image
  defining a first-page boundary is zero, so bilinearity proves both boundary
  conditions. Apply quotient descent and its represented formula.
  \leanok
\end{proof}

\begin{proposition}[Exact relative-boundary criterion for the page derivation formula]
  \label{prop:h6-long-pairing-leibniz-criterion}
  \lean{KIP126.Classical.Adams.adamsLongLayerBoundaryToPage,KIP126.Classical.Adams.adamsDifferential_longLayerToPage_eq_boundary,KIP126.Classical.Adams.AdamsLongLayerPairing.RelativeBoundaryFormula,KIP126.Classical.Adams.AdamsLongLayerPairing.leibniz_of_relativeBoundary,KIP126.Classical.Adams.AdamsLongLayerPairing.relativeBoundary_iff_leibniz,KIP126.Classical.Adams.AdamsLongLayerPairing.internalD_of_relativeBoundary}
  \uses{prop:h6-long-pairing-descent,prop:h6-tower-long-layer-differential}
  \leanok
  Put $b_r(z)=[j(k_rz)]$. Fix an integer coefficient $e$ and bilinear maps
  $L:E_r^{dp}\times E_r^q\to E_r^{d(p+q)}$ and
  $R:E_r^p\times E_r^{dq}\to E_r^{d(p+q)}$, where
  $d(s,t)=(s+r,t+r-1)$. Then the representative identity
  \[
    b_r(m_r(a,b))=L(b_ra,[p_rb])+eR([p_ra],b_rb)
  \]
  is equivalent to the page formula
  $d_r(xy)=L(d_rx,y)+eR(x,d_ry)$ for all $x,y$.
  For $r=n+2$, the same identity computes the existing internal differential
  of the descended internal product: compare the output and both input
  differentials with their existing quotient-page isomorphisms.
  This is a reduction of the remaining proof obligation, not its discharge:
  the boundary identity is unproved. One must also construct the intended
  adjacent-degree pairings $L,R$, supply their target identifications and
  correct sign, and prove compatibility with Lin multiplication.
\end{proposition}
\begin{proof}
  Lift both page inputs to long-layer representatives. Substitute the proved
  pairing formula and $d_r[p_rz]=b_r(z)$ in each term. The converse evaluates
  the page identity on those same representatives.
  Naturality of the existing internal differential comparison transports
  the formula to internal pages without replacing their differentials.
  \leanok
\end{proof}

\begin{definition}[The three geometric products for the square's second differential]
  \label{def:h6-three-long-layer-products}
  \lean{KIP126.Classical.Adams.SphereH6LongLayerMaps,KIP126.Classical.Adams.SphereH6LongLayerMaps.squarePairing,KIP126.Classical.Adams.SphereH6LongLayerMaps.leftPairing,KIP126.Classical.Adams.SphereH6LongLayerMaps.rightPairing,KIP126.Classical.Adams.SphereH6LongLayerMaps.Compatible,KIP126.Classical.Adams.SphereH6LongLayerMaps.Compatible.square,KIP126.Classical.Adams.SphereH6LongLayerMaps.Compatible.left,KIP126.Classical.Adams.SphereH6LongLayerMaps.Compatible.right,KIP126.Classical.Adams.SphereH6LongLayerMaps.RelativeBoundary}
  \uses{def:h6-specified-sphere-pairing,prop:h6-sphere-pairing-projection,prop:h6-long-pairing-leibniz-criterion}
  \leanok
  Write $Q_s=Q_s^{(2)}$. Supply actual spectrum maps
  $Q_1\wedge Q_1\to Q_2$, $Q_3\wedge Q_1\to Q_4$, and
  $Q_1\wedge Q_3\to Q_4$. Their representative pairings have input
  bidegrees respectively $((1,64),(1,64))$, $((3,65),(1,64))$, and
  $((1,64),(3,65))$. Each first-page product is fixed to the specified
  coefficient-induced product. Compatibility consists of the three
  spectrum projection squares and the two-sided tower-kernel boundary
  conditions for each product. The squares imply represented projection
  compatibility by tensor naturality.
  The relative-boundary predicate uses the latter two descended products
  as the two terms in the first product's boundary identity, with a plus
  sign on the target page. Thus the adjacent products are no longer free
  bilinear maps. The plus form is intended for characteristic two, not
  asserted as a spectrum-level graded-sign rule.
  These are defined data and predicates; no existence or compatibility
  witness for the specified fixed sphere is supplied here.
\end{definition}

\begin{proposition}[Simultaneous cofiber and connecting-map lifting]
  \label{prop:h6-simultaneous-cofiber-lift}
  \lean{KIP126.StableHomotopy.cofiberFactorizationMap_lift_of_boundary,KIP126.StableHomotopy.cofiberFactorizationMap_lift_iff,KIP126.Classical.Adams.adamsLongLayerProjection_lift_of_boundary,KIP126.Classical.Adams.adamsLongLayerProjection_lift_iff}
  \uses{def:h6-tower-long-layer,thm:stable-homotopy-long-exact}
  \leanok
  Suppose $f=ga:A\to D$, with $a:A\to B$ and $g:B\to D$,
  and let $c:\operatorname{cofib}(f)\to\operatorname{cofib}(g)$ be
  the chosen comparison. For any spectrum $W$, maps
  $x:W\to\operatorname{cofib}(g)$ and $y:W\to\Sigma A$ satisfying
  $(\Sigma a)y=\delta_gx$ admit a simultaneous lift
  $z:W\to\operatorname{cofib}(f)$ with $cz=x$ and $\delta_fz=y$.
  In particular, $x$ lifts through $c$ if and only if its connecting map
  lifts through $\Sigma a$. This applies to arbitrary source spectra,
  not just sphere representatives, and specializes to $Q_s^{(r)}\to L_s$.
\end{proposition}
\begin{proof}
  Since $(\Sigma f)y=0$, exactness of the represented distinguished
  triangle gives $z_0$ with $\delta_fz_0=y$. The difference $x-cz_0$
  has zero connecting map and therefore equals $\iota_gb$ for some
  $b:W\to D$. Then $z=z_0+\iota_fb$ has both required properties.
  Conversely, a lift $z$ supplies $y=\delta_fz$.
  \leanok
\end{proof}

\begin{proposition}[Constructed products retain a prescribed lower boundary]
  \label{prop:h6-product-from-boundary-lift}
  \lean{KIP126.Classical.Adams.adamsSphereLongLayerProjectedProduct,KIP126.Classical.Adams.adamsSphereLongLayerProductBoundary,KIP126.Classical.Adams.adamsSphereLongLayerProduct_exists_of_boundaryLift,KIP126.Classical.Adams.adamsSphereLongLayerProduct_exists_iff_boundaryLift,KIP126.Classical.Adams.adamsSphereLongLayerProductOfBoundaryLift,KIP126.Classical.Adams.adamsSphereLongLayerProductOfBoundaryLift_projection,KIP126.Classical.Adams.adamsSphereLongLayerProductOfBoundaryLift_boundary,KIP126.Classical.Adams.adamsSphereLongLayerPairingOfBoundaryLift_projection}
  \uses{prop:h6-simultaneous-cofiber-lift,def:h6-specified-sphere-pairing,prop:h6-sphere-pairing-projection}
  \leanok
  Fix the coefficient-induced product $m_1:L_s\wedge L_t\to L_{s+t}$
  and put $P=m_1(\pi_s\wedge\pi_t)$. A product
  $\mu:Q_s^{(r)}\wedge Q_t^{(r)}\to Q_{s+t}^{(r)}$ with
  $\pi_{s+t}\mu=P$ exists if and only if there is a map
  $y:Q_s^{(r)}\wedge Q_t^{(r)}\to\Sigma T_{s+t+r}$ satisfying
  \[
    (\Sigma I_{s+t+1,s+t+r})y=\delta_{s+t}^{(1)}P.
  \]
  Given such a $y$, the specified choice of $\mu$ satisfies both
  $\pi_{s+t}\mu=P$ and $\delta_{s+t}^{(r)}\mu=y$.
  Its homotopy pairings consequently satisfy projection compatibility
  in every input bidegree. No existence of $y$ is asserted.
\end{proposition}
\begin{proof}
  Apply simultaneous lifting with source $Q_s^{(r)}\wedge Q_t^{(r)}$
  and choose a witness retaining both equalities. Tensor naturality
  gives the representative projection formula.
  \leanok
\end{proof}

\begin{proposition}[The supplied lift computes the actual differential]
  \label{prop:h6-product-lift-differential}
  \lean{KIP126.StableHomotopy.connectingHomomorphism_eq_desuspend,KIP126.Classical.Adams.adamsSphereLongLayerProductOfBoundaryLift_K,KIP126.Classical.Adams.adamsSphereLongLayerProductOfBoundaryLift_differential}
  \uses{prop:h6-product-from-boundary-lift,prop:h6-tower-long-layer-differential}
  \leanok
  For homotopy representatives $a,b$ in these two long layers, the
  connecting class of $\mu_*(a,b)$ is the desuspension of $y_*(a,b)$.
  Thus, with the prescribed degree identifications,
  \[
    d_r[\mu_*(a,b)]=j_r\bigl(\operatorname{desusp}(y_*(a,b))\bigr).
  \]
  This computes the existing Adams differential, without assuming a
  Leibniz rule or that the product has descended to quotient pages.
  The formula does not establish a two-term boundary identity for $y$.
\end{proposition}
\begin{proof}
  Express the connecting homomorphism as postcomposition by the triangle's
  third map followed by desuspension. Substitute $\delta\mu=y$ and the
  already proved long-layer formula for the differential.
  \leanok
\end{proof}

\begin{definition}[Lower lifts for the three square products]
  \label{def:h6-three-boundary-lifts}
  \lean{KIP126.Classical.Adams.SphereH6BoundaryLifts,KIP126.Classical.Adams.SphereH6BoundaryLifts.FitsProducts,KIP126.Classical.Adams.SphereH6BoundaryLifts.toLongLayerMaps,KIP126.Classical.Adams.SphereH6BoundaryLifts.toLongLayerMaps_boundaries,KIP126.Classical.Adams.SphereH6BoundaryLifts.toLongLayerMaps_compatible}
  \uses{prop:h6-product-from-boundary-lift,def:h6-three-long-layer-products}
  \leanok
  At $r=2$, supply three maps
  $Q_1\wedge Q_1\to\Sigma T_4$,
  $Q_3\wedge Q_1\to\Sigma T_6$, and
  $Q_1\wedge Q_3\to\Sigma T_6$. The predicate
  \texttt{FitsProducts} requires their composites through
  $\Sigma I_{3,4}$ and $\Sigma I_{5,6}$ to equal the three prescribed
  first-layer product boundaries. Simultaneous lifting constructs the
  three long-layer products, proves their projection squares, and retains
  these exact connecting maps. The two-sided tower-kernel conditions
  are inputs to the general descent constructor; four of the six are
  discharged by the next proposition. Neither a value of this structure for the
  fixed foundation nor the relative-boundary or Lin comparisons is supplied.
\end{definition}

\begin{proposition}[Filtration-one boundaries and redundant descent conditions]
  \label{prop:h6-filtration-one-boundaries}
  \lean{KIP126.Classical.Adams.adamsBoundaries_succ_eq_of_differential_zero,KIP126.Classical.Adams.sphereAdamsBoundaries_filtration_one_succ,KIP126.Classical.Adams.sphereAdamsBoundaries_filtration_one,KIP126.Classical.Adams.adamsJ_eq_zero_of_boundaries_eq_bot,KIP126.Classical.Adams.AdamsLongLayerPairing.boundaryCompatible_of_eq_bot,KIP126.Classical.Adams.SphereH6LongLayerMaps.CrossBoundaryCompatible,KIP126.Classical.Adams.SphereH6LongLayerMaps.square_boundary,KIP126.Classical.Adams.SphereH6LongLayerMaps.CrossBoundaryCompatible.left_boundary,KIP126.Classical.Adams.SphereH6LongLayerMaps.CrossBoundaryCompatible.right_boundary,KIP126.Classical.Adams.SphereH6BoundaryLifts.toLongLayerMaps_compatible_of_cross}
  \uses{prop:h6-sphere-initial-survival,def:h6-three-long-layer-products,def:h6-three-boundary-lifts}
  \leanok
  In the actual sphere tower, $B_r^{1,t}=0$ for every $r\ge1$ and
  every integer $t$. The first incoming differential is zero by the
  initial-column theorem, and later incoming sources have negative
  filtration. Thus any bilinear representative pairing kills a boundary
  input of filtration one. For the three products used for $h_6^2$,
  this proves both boundary conditions of the square product and the
  filtration-one condition of each cross product. Only the two conditions
  with boundary input in bidegree $(3,65)$ remain as descent obligations.
  They are recorded in \texttt{CrossBoundaryCompatible}; no Lin comparison,
  fixed-foundation axiom, or relative-boundary formula enters this reduction.
\end{proposition}
\begin{proof}
  A zero incoming differential makes consecutive boundary submodules
  equal: every new boundary is represented by an image of that differential.
  Apply this at filtration zero for $r=1$, and use the zero negative
  source for $r>1$. Induction starts from $B_1=0$.
  If $I(a)=0$ in filtration one, then $J(a)\in B_r^{1,t}=0$.
  Bilinearity therefore kills every such boundary input, proving four
  of the six conditions. Combine these proofs with the two remaining
  cross conditions and the constructed projection squares.
  \leanok
\end{proof}

\begin{proposition}[The remaining descent conditions from local first-differential rules]
  \label{prop:h6-cross-boundary-from-first-differential}
  \lean{KIP126.Classical.Adams.adamsPageOneEquiv_d_mem_boundaries,KIP126.Classical.Adams.adamsBoundary_two_exists_first_primitive,KIP126.Classical.Adams.SphereH6FirstCycleProductRule,KIP126.Classical.Adams.SphereH6FirstCycleProductRule.crossBoundaryCompatible}
  \uses{prop:h6-filtration-one-boundaries,prop:h6-long-pairing-descent}
  \leanok
  Fix the already constructed coefficient product on the first quotient
  page. Assume its two local $d_1$ product identities on an arbitrary
  input in bidegree $(2,65)$ and a $d_1$ cycle in $(1,64)$, in both
  orders. Writing these inputs as $x$ and $z$, the identities are
  \[
    d_1m_{2,1}(x,z)=m_{3,1}(d_1x,z),\qquad
    d_1m_{1,2}(z,x)=m_{1,3}(z,d_1x),
  \]
  with the specified first-page-to-layer comparisons and the mod-two
  unsigned convention. Then the two remaining cross boundary conditions follow for
  every choice of the three long-layer maps. The proof represents a
  boundary in $(3,65)$ as $d_1$ of an input in $(2,65)$ and uses the
  identities to exhibit the product as a boundary in $(4,129)$.
  The local identities themselves remain geometric obligations; neither
  a global Leibniz theorem nor a new axiom is asserted. This is a sufficient
  criterion for descent, not a proof of those identities or of the
  second-differential relative-boundary formula.
\end{proposition}
\begin{proof}
  Use exactness and the first-page quotient comparison to choose a
  primitive $x$ whose actual $d_1$ representative is the given boundary.
  Map the other long-layer representative into the actual second-cycle
  submodule. Apply the appropriate identity: the product is the first-page
  representative of $d_1m(x,z)$ or $d_1m(z,x)$, hence is in $B_2^{4,129}$.
  The prescribed first-page component of every long-layer pairing is the
  same coefficient product, so the conclusion does not depend on its
  chosen long-layer spectrum map.
  \leanok
\end{proof}

\begin{proposition}[Geometric square vanishing from cross-term cancellation]
  \label{prop:h6-long-layer-cross-cancellation}
  \lean{KIP126.Classical.Adams.SphereH6LongLayerMaps.internalD_square_eq_zero_of_cross_sum}
  \uses{def:h6-three-long-layer-products,prop:h6-long-pairing-leibniz-criterion}
  \leanok
  Let the three maps satisfy the descent and relative-boundary predicates,
  and let $x$ be an internal second-page class in bidegree $(1,64)$.
  If the two descended cross products $m_L(d_2x,x)+m_R(x,d_2x)$ sum to
  zero, then the existing internal differential of $m_0(x,x)$ is zero.
  This implication uses no fixed-foundation or Lin axiom.
\end{proposition}
\begin{proof}
  Compare the internal relative-boundary formula with the two internal
  cross products, apply additivity to their zero sum, and use injectivity
  of the target page isomorphism.
  \leanok
\end{proof}

\begin{definition}[Coefficient-homology resolution differential]
  \label{def:h6-resolution-differential}
  \lean{KIP126.Classical.Adams.adamsResolutionConnecting,KIP126.Classical.Adams.adamsResolutionBoundary,KIP126.Classical.Adams.adamsResolutionDifferential}
  \uses{def:h6-fiber-tower,prop:h6-tower-layer-comparison}
  \leanok
  For $s\ge0$, let $\partial_s:H\wedge T_s\to\Sigma T_{s+1}$ be
  minus the third morphism of the inverse-rotated unit cofiber triangle.
  The sign is required by the chosen positive fiber inclusion.
  Its induced connecting homomorphism is
  $\beta_s:\pi_n(H\wedge T_s)\to\pi_{n-1}T_{s+1}$.
  Define the integer-linear map
  \[
    D_s:\pi_{t-s}(H\wedge T_s)\longrightarrow
      \pi_{t-s-1}(H\wedge T_{s+1}),\qquad
    D_s=(\eta_{T_{s+1}})_*\circ\beta_s.
  \]
  This definition uses only the unit and its cofiber, not page coordinates
  or an independently supplied spectral-sequence differential.
\end{definition}

\begin{proposition}[The actual first differential is the resolution differential]
  \label{prop:h6-first-differential-homology}
  \lean{KIP126.Classical.Adams.adamsK_eq_resolutionBoundary,KIP126.Classical.Adams.adamsCycleLift_one,KIP126.Classical.Adams.adamsPageD_one_homology,KIP126.Classical.Adams.adamsResolutionDifferential_comp,KIP126.Classical.Adams.sphereFirstDifferential_homology}
  \uses{def:h6-resolution-differential,prop:h6-first-page-homology,prop:h6-tower-differential-laws}
  \leanok
  Write $e_s:E_1^{s,t}\cong\pi_{t-s}(H\wedge T_s)$ for the constructed
  first-page comparison. Then $e_{s+1}d_1=D_se_s$ and $D_{s+1}D_s=0$.
  This applies in particular to the precise \texttt{sphereFirstDifferential}
  used by the existing Milnor-coordinate interface, without requiring a
  value of that interface. The comparison with the Milnor coproduct
  differential is a further, still unproved step.
\end{proposition}
\begin{proof}
  The connecting-map comparison of the layer triangles identifies $k$
  with $\beta_s$. At page one, the map through which the cycle is lifted
  is the identity, so the chosen lift equals $k(x)$. The quotient
  differential is consequently $jk$. Under $e_{s+1}$ the map $j$ is
  induced by the next unit. Explicit degree transports identify
  $t-s-1$ with $t-(s+1)$. Finally transport the proved square-zero
  identity of the actual page differential through these isomorphisms.
  \leanok
\end{proof}

\begin{definition}[The smashed unit triangle and its homology maps]
  \label{def:h6-homology-resolution}
  \lean{KIP126.Classical.Adams.adamsResolutionTriangle,KIP126.Classical.Adams.adamsResolutionTriangle_distinguished,KIP126.Classical.Adams.adamsHomologySequence,KIP126.Classical.Adams.adamsHomologyUnit,KIP126.Classical.Adams.adamsHomologyAction,KIP126.Classical.Adams.adamsHomologyBoundary}
  \uses{def:h6-resolution-differential,def:h6-ring-cooperations}
  \leanok
  In addition to the specified ring structure on $H$, suppose that
  $H\wedge-$ commutes coherently with suspension and is a triangulated
  functor. These are explicit structural parameters, not global instances.
  For $F(X)=\operatorname{fib}(\eta_X)$, apply this functor to the
  distinguished triangle
  $F(X)\to X\to H\wedge X\to\Sigma F(X)$ with the positive fiber
  inclusion. The sign-corrected triangle is proved distinguished by an
  isomorphism with the inverse rotation of the unit cofiber triangle.
  In each degree $n$, write
  \[
    U_X:H_n(X)\to H_n(H\wedge X),\qquad
    A_X:H_n(H\wedge X)\to H_n(X),\qquad
    b_X:H_n(H\wedge X)\to H_{n-1}(F(X)).
  \]
  Here $U_X$ is induced by the unit, $A_X$ by multiplication on the first
  two factors, and $b_X$ is the connecting map of the smashed triangle.
  The notation $H_n(Y)$ means $\pi_n(H\wedge Y)$.
\end{definition}

\begin{proposition}[Derived normalized homology of the next tower stage]
  \label{prop:h6-homology-normalization}
  \lean{KIP126.Classical.Adams.adamsHomologyAction_unit,KIP126.Classical.Adams.adamsHomologyUnit_injective,KIP126.Classical.Adams.adamsHomologyBoundary_exact,KIP126.Classical.Adams.adamsHomologyBoundary_surjective,KIP126.Classical.Adams.adamsHomologyKernelEquiv,KIP126.Classical.Adams.adamsHomologyQuotientEquiv,KIP126.Classical.Adams.adamsHomologyQuotientEquiv_mkQ,KIP126.Classical.Adams.adamsHomologyKernelEquiv_symm_boundary}
  \uses{def:h6-homology-resolution,prop:h6-tower-homology-zero}
  \leanok
  The maps satisfy $A_XU_X=1$, $\ker b_X=\operatorname{im}U_X$, and
  $b_X$ is surjective. Thus there are constructed integer-linear isomorphisms
  \[
    H_n(H\wedge X)/\operatorname{im}U_X
      \ \cong\ H_{n-1}(F(X))
      \ \cong\ \ker A_X.
  \]
  The quotient map sends $[z]$ to $b_X(z)$; the inverse of the restricted
  boundary sends $b_X(z)$ to $z-U_XA_X(z)$. In particular, setting $X=T_s$
  gives a recursive description of $H_{n-1}(T_{s+1})$ from $T_s$.
  No coordinates on Adams pages, K\"unneth isomorphism, or Milnor polynomial
  algebra are assumed in this construction.
\end{proposition}
\begin{proof}
  The unit law gives the retraction. Apply the homotopy long exact sequence
  to the smashed triangle. Its first map is zero by
  Proposition~\ref{prop:h6-tower-homology-zero}, giving surjectivity of $b_X$.
  Exactness identifies its kernel. Restricting to $\ker A_X$ is injective
  because $U_X(y)\in\ker A_X$ forces $y=0$. It is surjective because any
  lift $z$ can be replaced by $z-U_XA_X(z)$ without changing its boundary.
  These maps, not abstract dimension comparisons, construct the isomorphisms.
  \leanok
\end{proof}

\begin{proposition}[Filtration one of the actual first sphere page]
  \label{prop:h6-first-filtration-cooperations}
  \lean{KIP126.Classical.Adams.sphereCoefficientHomologyEquiv_map_ker,KIP126.Classical.Adams.sphereFirstTowerHomologyEquiv,KIP126.Classical.Adams.sphereFirstTowerHomologyEquiv_boundary,KIP126.Classical.Adams.sphereFirstPageFiltrationOneEquiv}
  \uses{prop:h6-homology-normalization,prop:h6-first-page-homology,def:h6-ring-cooperations}
  \leanok
  Under the same explicit structural parameters, the right unitor
  identifies $A_S$ with the cooperation counit
  $\epsilon_t:\pi_t(H\wedge H)\to\pi_tH$. Therefore
  \[
    E_1^{1,t}\cong H_{t-1}(T_1S)\cong\ker\epsilon_t.
  \]
  The boundary-class formula agrees with normalization by $z-U_SA_S(z)$.
  This constructs the first nontrivial filtration of the first-page
  coordinate comparison, before identifying cooperations with the Milnor
  algebra. It is not the complete first-page Milnor comparison.
\end{proposition}
\begin{proof}
  Right-unitor naturality and its tensor identity identify multiplication on
  $H\wedge(H\wedge S)$ with multiplication on $H\wedge H$. Restrict the
  induced homotopy equivalence to the two kernels and compose with the
  normalized tower-homology and actual first-page equivalences.
  \leanok
\end{proof}

The fixed foundation does not yet supply the ring structure and the coherent
triangulated structure on $H\wedge-$ used above. These results do not add
such structures as new fixed axioms. K\"unneth compatibility, the Milnor
cooperation identification, the higher-filtration coordinate construction,
and its differential compatibility remain to be derived from the agreed
lower-level foundations.

\begin{proposition}[Coaction identities before K\"unneth]
  \label{prop:h6-coaction-identities}
  \lean{KIP126.StableHomotopy.Cohomology.mod2Coaction,KIP126.StableHomotopy.Cohomology.mod2CoactionMap,KIP126.StableHomotopy.Cohomology.mod2Coaction_naturality,KIP126.StableHomotopy.Cohomology.mod2Coaction_coface,KIP126.StableHomotopy.Cohomology.mod2Coaction_counit,KIP126.StableHomotopy.Cohomology.mod2FreeAction_naturality,KIP126.StableHomotopy.Cohomology.cooperationDiagonal_counit_left,KIP126.StableHomotopy.Cohomology.mod2CoactionMap_coface}
  \uses{def:h6-ring-cooperations,def:h6-fiber-tower}
  \leanok
  Write $\eta_X:X\to H\wedge X$ for the actual Adams unit and
  $\delta_X=H\wedge\eta_X$. Then $\delta$ is natural in $X$ and satisfies
  \[
    \delta_{H\wedge X}\circ\delta_X
      =(H\wedge\delta_X)\circ\delta_X.
  \]
  With the specified ring structure, the free action $a_X$ is natural and
  $a_X\circ\delta_X=\operatorname{id}$. For $X=H$, this is the left
  counit identity for the concrete cooperation diagonal; the right counit
  identity uses $H\wedge\mu$. These identities also hold on represented
  homology. They do not assert a K\"unneth isomorphism or polynomial coordinates.
\end{proposition}
\begin{proof}
  Unitor naturality and the interchange law prove naturality of $\eta$.
  Apply it to $\eta_X$, then tensor with $H$, to obtain the coface identity.
  Associator naturality proves naturality of the free action; the ring unit
  laws give both counit identities. Postcomposition induces the homology laws.
  \leanok
\end{proof}

\begin{proposition}[Natural normalization and actual tower splitting]
  \label{prop:h6-natural-tower-splitting}
  \lean{KIP126.Classical.Adams.adamsHomologyUnit_eq_coaction,KIP126.Classical.Adams.adamsHomologyNormalize,KIP126.Classical.Adams.adamsHomologyNormalize_naturality,KIP126.Classical.Adams.adamsHomologyNormalize_idempotent,KIP126.Classical.Adams.adamsHomologyNormalize_range,KIP126.Classical.Adams.adamsHomologyNormalize_ker,KIP126.Classical.Adams.adamsHomologyBoundary_normalize,KIP126.Classical.Adams.adamsHomologySplitEquiv,KIP126.Classical.Adams.adamsHomologySplitEquiv_apply,KIP126.Classical.Adams.adamsHomologySplitEquiv_symm_apply}
  \uses{prop:h6-coaction-identities,prop:h6-homology-normalization}
  \leanok
  The homology unit $U_X$ is the induced coaction. The operator
  $P_X=\operatorname{id}-U_XA_X$ is natural and idempotent, with image
  $\ker A_X$ and kernel $\operatorname{im}U_X$.
  Under the tensor exactness and shift compatibility used above,
  the actual connecting map $b_X$ satisfies $b_XP_X=b_X$, and
  \[
    H_n(H\wedge X;\mathbb F_2)
      \xrightarrow[\cong]{(A_X,b_X)}
    H_n(X;\mathbb F_2)\oplus H_{n-1}(\operatorname{fib}\eta_X;\mathbb F_2).
  \]
  Here $H_n(Y;\mathbb F_2)=\pi_n(H\wedge Y)$; initially this is an
  isomorphism of integer modules. It applies in particular to every
  $X=T_s$, with fiber $T_{s+1}$, without a Milnor-coordinate input.
\end{proposition}
\begin{proof}
  Naturality follows from that of $U$ and $A$, and $AU=1$ proves the
  projector identities. The inverse of $(A,b)$ sends $(x,y)$ to
  $U(x)+k^{-1}(y)$, where $k:\ker A\cong H_{n-1}(\operatorname{fib}\eta_X;\mathbb F_2)$
  is the proved restriction of the actual boundary. The normalization
  formula for $k^{-1}b$ verifies both inverse identities.
  \leanok
\end{proof}

\begin{proposition}[Mod--two scalars derived from the ring object]
  \label{prop:h6-derived-mod2-scalars}
  \lean{KIP126.StableHomotopy.Cohomology.mod2Unit_add_self,KIP126.StableHomotopy.Cohomology.mod2_id_add_self,KIP126.StableHomotopy.Cohomology.mod2Homology_two_nsmul_zero,KIP126.StableHomotopy.Cohomology.mod2Cohomology_two_nsmul_zero,KIP126.StableHomotopy.Cohomology.mod2HomologyModule,KIP126.StableHomotopy.Cohomology.mod2CohomologyModule,KIP126.StableHomotopy.Cohomology.mod2HomologyModule_eq,KIP126.StableHomotopy.Cohomology.mod2HomologyF2Functor,KIP126.StableHomotopy.Cohomology.cooperationCounitF2,KIP126.StableHomotopy.Cohomology.cooperationDiagonalF2}
  \uses{def:h6-ring-cooperations,def:h6-em-object}
  \leanok
  Suppose the tensor product is additive in both variables
  (\texttt{MonoidalPreadditive}) and the chosen $H$ has the specified
  unital multiplication. Then $2u=0$, $2\operatorname{id}_H=0$, and
  $2\operatorname{id}_{H\wedge X}=0$. Consequently the existing represented
  homology and cohomology groups have canonical $\mathbb F_2$-module
  structures. Homology pushforward, the cooperation counit and diagonal,
  and the specified $\pi_0H\cong\mathbb F_2$ coordinate are linear for these
  structures. The underlying groups and maps are unchanged. The homology
  scalar action is independent of the compatible multiplication used to
  prove two-torsion.
\end{proposition}
\begin{proof}
  The unit represents $1\in\mathbb F_2$, hence doubles to zero. Tensor
  additivity and the unit law transfer this to the identity of $H$, and
  then to the identity of $H\wedge X$. Postcomposition proves two-torsion
  for homology and cohomology classes. Apply the canonical
  $\mathbb Z/2$-module construction for an abelian group killed by two.
  Additive homomorphisms of such modules are automatically linear; the
  scalar structure is unique on the given additive group.
  \leanok
\end{proof}

\begin{proposition}[The constructed finite pages and comparisons are mod--two linear]
  \label{prop:h6-derived-mod2-pages}
  \lean{KIP126.Classical.Adams.adamsE1_two_nsmul_zero,KIP126.Classical.Adams.adamsPage_two_nsmul_zero,KIP126.Classical.Adams.adamsInternalPage_two_nsmul_zero,KIP126.Classical.Adams.adamsPageF2Module,KIP126.Classical.Adams.adamsInternalPageF2Module,KIP126.Classical.Adams.adamsDifferentialF2_apply,KIP126.Classical.Adams.adamsInternalDifferentialF2_apply,KIP126.Classical.Adams.adamsInternalDifferentialF2_comp,KIP126.Classical.Adams.adamsPageOneHomologyF2Equiv,KIP126.Classical.Adams.adamsHomologyKernelF2Equiv,KIP126.Classical.Adams.adamsHomologyKernelF2Equiv_apply,KIP126.Classical.Adams.adamsHomologySplitF2Equiv,KIP126.Classical.Adams.adamsHomologySplitF2Equiv_apply}
  \uses{prop:h6-derived-mod2-scalars,prop:h6-first-page-homology,def:h6-tower-internal-sequence,prop:h6-homology-normalization,prop:h6-natural-tower-splitting}
  \leanok
  Under the same ring and additive-tensor parameters, every finite tower
  quotient page and every finite internal SSData page is killed by two.
  They carry constructed $\mathbb F_2$-module structures on their existing
  additive groups. Their differentials are the original maps, now expressed
  as $\mathbb F_2$-linear maps, and retain the proved square-zero identity.
  The first-page homology comparison is $\mathbb F_2$-linear; with the
  additional tensor exactness and suspension compatibility of
  Proposition~\ref{prop:h6-homology-normalization}, so is the normalized
  next-stage kernel comparison and the actual tower splitting $(A_X,b_X)$.
\end{proposition}
\begin{proof}
  In nonnegative filtration transport two-torsion through the layer
  homology equivalence; in negative filtration the layer is zero.
  Subgroups and quotient groups inherit two-torsion. Transport this
  through the proved internal finite-page isomorphism. Reinterpret the
  existing additive maps and equivalences using their automatic
  $\mathbb F_2$ linearity; the representative maps do not change.
  \leanok
\end{proof}

These are explicit local scalar structures. The existing integer-module
SSData object is retained, and no multiplication or tensor-additivity
witness has been silently added to the fixed foundation. This proves a
prerequisite for the graded K\"unneth comparison, not that comparison itself.

\begin{proposition}[Reduced cooperation tensors and their actual counit kernel]
  \label{prop:h6-graded-tensor-kernel}
  \lean{KIP126.Core.Algebra.gradedTensor,KIP126.Core.Algebra.gradedTensorKernelEquiv,KIP126.Core.Algebra.gradedTensorKernelEquiv_coe,KIP126.StableHomotopy.Cohomology.cooperationTensorCounit,KIP126.StableHomotopy.Cohomology.reducedCooperationTensorEquiv,KIP126.StableHomotopy.Cohomology.coefficientTensorEquiv,KIP126.StableHomotopy.Cohomology.coefficientTensorEquiv_zero_tmul,KIP126.StableHomotopy.Cohomology.reducedCooperationAugmentationEquiv}
  \uses{def:h6-ring-cooperations,prop:h6-derived-mod2-scalars}
  \leanok
  Put $\mathcal A_i=\pi_i(H\wedge H)$ and
  $\overline{\mathcal A}_i=\ker(\varepsilon_i:\mathcal A_i\to\pi_iH)$,
  using the actual multiplication-induced counit. For every graded
  $\mathbb F_2$-module $V$, define
  $(\mathcal A\otimes V)_n=\bigoplus_{i\in\mathbb Z}\mathcal A_i\otimes V_{n-i}$.
  Then
  \[
    (\overline{\mathcal A}\otimes V)_n
      \cong\ker\bigl((\mathcal A\otimes V)_n\longrightarrow V_n\bigr).
  \]
  The arrow is $\varepsilon\otimes 1$ followed by the constructed
  identification $(\pi_*H\otimes V)_n\cong V_n$. The kernel isomorphism
  preserves the ordinary tensor representatives. No K\"unneth or Milnor
  identification is required for this algebraic result.
\end{proposition}
\begin{proof}
  Tensoring over a field preserves kernels, and direct sums preserve the
  componentwise kernel. The coefficient tensor has only its $i=0$ summand:
  all other $\pi_iH$ vanish. On that summand apply the specified
  $\pi_0H\cong\mathbb F_2$ and the tensor unitor. This last map is an
  isomorphism and hence does not change the kernel.
  \leanok
\end{proof}

\begin{definition}[Homology-level K\"unneth input]
  \label{def:h6-homology-kunneth-input}
  \lean{KIP126.StableHomotopy.Cohomology.Mod2CooperationKunneth}
  \uses{prop:h6-graded-tensor-kernel,prop:h6-coaction-identities}
  \leanok
  A \texttt{Mod2CooperationKunneth} value consists of natural linear
  isomorphisms, for every spectrum $X$ and integer $n$,
  \[
    \kappa_{X,n}:H_n(H\wedge X;\mathbb F_2)
       \cong(\mathcal A\otimes H_*(X;\mathbb F_2))_n,
  \]
  compatible with the specified multiplication: tensor augmentation of
  $\kappa_{X,n}(z)$ equals the actual free-action image $A_X(z)$.
  This record contains no Adams pages, coordinates on those pages,
  differential formula, or permanence assertion. The record is defined,
  but no value for the fixed foundation is postulated or constructed here.
  In particular, its existence and the additional diagonal compatibilities
  needed for a full Milnor differential comparison remain obligations.
\end{definition}

\begin{definition}[Below-page suspension compatibility of K\"unneth]
  \label{def:h6-kunneth-suspension-input}
  \lean{KIP126.StableHomotopy.Cohomology.Mod2KunnethSuspensionCompatible}
  \uses{def:h6-homology-kunneth-input}
  \leanok
  The predicate \texttt{Mod2KunnethSuspensionCompatible} requires that
  $\kappa$ commute with the actual degree-minus-one desuspension maps:
  double coefficient homology desuspension on the source corresponds to
  $1\otimes\sigma_X$ on the target. The degree-$i$ cooperation factor is
  unchanged. This condition concerns spectra and suspension only; it has
  no distinguished triangle, Adams boundary, page, or differential field.
  The predicate is defined, but its satisfaction on the fixed foundation
  has not been proved or postulated.
\end{definition}

\begin{proposition}[K\"unneth and the actual coaction commute with boundaries]
  \label{prop:h6-kunneth-coaction-boundary}
  \lean{KIP126.StableHomotopy.connectingHomomorphism_map_naturality,KIP126.StableHomotopy.Cohomology.mod2Kunneth_connecting,KIP126.StableHomotopy.Cohomology.mod2TensorCoaction,KIP126.StableHomotopy.Cohomology.mod2TensorCoaction_counit,KIP126.StableHomotopy.Cohomology.mod2TensorCoaction_injective,KIP126.StableHomotopy.Cohomology.mod2TensorCoaction_naturality,KIP126.StableHomotopy.Cohomology.mod2TensorCoaction_connecting}
  \uses{def:h6-kunneth-suspension-input,prop:h6-coaction-identities}
  \leanok
  Assume the explicit ring and additive tensor structures, exactness and
  suspension compatibility of $H\wedge-$, and the preceding K\"unneth
  suspension condition. For every distinguished triangle $T$, its actual
  coefficient-homology connecting map $\partial_T$ satisfies
  \[
    \kappa\partial_{H\wedge T}
       =(1\otimes\partial_T)\kappa.
  \]
  Define $\rho_X=\kappa_X U_X$. This is natural, has tensor augmentation
  as a left inverse, and hence is injective. If the unit natural
  transformation commutes with suspension, then also
  \[
    \rho_{T.X}\partial_T=(1\otimes\partial_T)\rho_{T.Z}.
  \]
  These are conditional theorems, not new boundary comparison axioms.
  A tensor coproduct formula or Milnor identification is not asserted.
\end{proposition}
\begin{proof}
  Factor the actual connecting map into the last arrow of the triangle
  followed by desuspension. Apply naturality of $\kappa$ to that arrow
  and the explicit suspension condition to desuspension. Naturality of
  connecting maps for shift-compatible natural transformations proves
  that inner-unit insertion commutes with the boundary. Compose the two
  equalities. The counit identity follows from multiplication compatibility
  of $\kappa$ and the already proved unit/action identity.
  \leanok
\end{proof}

\begin{definition}[Actual tensor diagonal and below-page K\"unneth coherence]
  \label{def:h6-kunneth-diagonal-input}
  \lean{KIP126.Core.Algebra.gradedTensorAssoc,KIP126.StableHomotopy.Cohomology.cooperationTensorDiagonal,KIP126.StableHomotopy.Cohomology.cooperationTensorComultiply,KIP126.StableHomotopy.Cohomology.Mod2KunnethDiagonalCompatible}
  \uses{def:h6-homology-kunneth-input,prop:h6-kunneth-coaction-boundary}
  \leanok
  Define the actual tensor-coordinate diagonal
  $\Delta=\rho_H:\mathcal A\to\mathcal A\otimes\mathcal A$.
  On $\mathcal A\otimes V$, define $\Delta\otimes1$ followed by the
  explicit graded tensor reassociation. Its value on an elementary
  tensor is the ordinary reassociation, with the necessary integer-degree
  cast; this is constructed using finite-support direct sums.
  The predicate \texttt{Mod2KunnethDiagonalCompatible} requires, for every
  spectrum $X$,
  \[
    (1\otimes\kappa_X)\rho_{H\wedge X}
        =(\Delta\otimes1)\kappa_X.
  \]
  This is an explicit coherence premise on the homology-level K\"unneth
  comparison, not an Adams page or differential comparison. No fixed
  witness is postulated, and no Milnor polynomial formula is part of it.
\end{definition}

\begin{proposition}[Derived coassociativity of the actual tensor coaction]
  \label{prop:h6-tensor-coaction-coassoc}
  \lean{KIP126.Core.Algebra.gradedTensorAssoc_lof_tmul,KIP126.StableHomotopy.Cohomology.cooperationTensorDiagonal_apply,KIP126.StableHomotopy.Cohomology.cooperationTensorDiagonal_counit,KIP126.StableHomotopy.Cohomology.mod2TensorCoaction_coassoc,KIP126.StableHomotopy.Cohomology.cooperationTensorDiagonal_coassoc}
  \uses{def:h6-kunneth-diagonal-input}
  \leanok
  Under that explicit diagonal coherence premise, every actual coaction
  satisfies $(\Delta\otimes1)\rho_X=(1\otimes\rho_X)\rho_X$.
  Taking $X=H$ gives coassociativity of $\Delta$. Its left tensor
  augmentation is the identity, independently of diagonal coherence.
  The diagonal is exactly the original middle-unit insertion followed
  by K\"unneth; it is not defined by the desired polynomial formula.
\end{proposition}
\begin{proof}
  Apply naturality of the actual coaction to $\eta_X$. Use diagonal
  coherence to replace the free-object coaction on one side, then compose
  the maps on the right tensor factor. The counit is the previously
  proved tensor-coaction counit specialized to $H$.
  \leanok
\end{proof}

\begin{proposition}[Derived reduced-tensor recurrence for the actual tower]
  \label{prop:h6-kunneth-tower-recurrence}
  \lean{KIP126.Classical.Adams.adamsHomologyKunneth_map_ker,KIP126.Classical.Adams.adamsNextHomologyTensorEquiv,KIP126.Classical.Adams.adamsTowerHomologyTensorEquiv,KIP126.Classical.Adams.adamsNextHomologyTensorEquiv_boundary}
  \uses{def:h6-homology-kunneth-input,prop:h6-graded-tensor-kernel,prop:h6-derived-mod2-pages}
  \leanok
  Given that homology-level input, the explicit ring structure, additive
  tensor, and the tensor exactness and shift compatibility used in the
  homology sequence, the actual next fiber satisfies
  \[
    H_{n-1}(\operatorname{fib}\eta_X;\mathbb F_2)
      \cong(\overline{\mathcal A}\otimes H_*(X;\mathbb F_2))_n.
  \]
  For an actual boundary $b_X(z)$, its reduced tensor coordinate, included
  in the unreduced tensor, is $\kappa_{X,n}(z-U_XA_Xz)$.
  The result applies to every $X=T_s$, so it is a derived recurrence for
  the existing tower, not extra tower data.
\end{proposition}
\begin{proof}
  The previously constructed boundary isomorphism identifies next-stage
  homology with $\ker A_X$. Multiplication compatibility restricts
  $\kappa$ to the tensor augmentation kernel, which is the reduced tensor
  by Proposition~\ref{prop:h6-graded-tensor-kernel}. The representative
  formula is the proved normalization formula for the boundary inverse.
  \leanok
\end{proof}

\begin{proposition}[Actual next-stage coaction from a normalized lift]
  \label{prop:h6-next-stage-coaction}
  \lean{KIP126.Classical.Adams.adamsHomologyBoundary_eq_connecting,KIP126.Classical.Adams.adamsTensorCoaction_boundary,KIP126.Classical.Adams.adamsTensorCoaction_next}
  \uses{prop:h6-kunneth-coaction-boundary,prop:h6-homology-normalization,prop:h6-kunneth-tower-recurrence}
  \leanok
  Under the preceding explicit hypotheses, let $b_X$ be the actual Adams
  homology boundary. For any $y\in H_{n-1}(\operatorname{fib}\eta_X)$,
  let $z\in\ker A_X$ be its unique normalized lift, so $b_X(z)=y$.
  Then
  \[
    \rho_{\operatorname{fib}\eta_X}(y)
       =(1\otimes b_X)\rho_{H\wedge X}(z).
  \]
  This holds for every spectrum $X$, in particular every actual tower
  stage $T_s(S^0)$. The coaction on $H\wedge X$ still needs a compatible
  tensor coproduct description before this yields the Milnor cobar formula.
\end{proposition}
\begin{proof}
  The previously defined Adams boundary is definitionally the connecting
  map of its coefficient-homology triangle. Specialize the preceding
  proposition and use the proved inverse of the normalized boundary
  isomorphism. No new tower recurrence is assumed.
  \leanok
\end{proof}

\begin{proposition}[Next-stage coaction expressed by the actual coproduct]
  \label{prop:h6-next-stage-coproduct}
  \lean{KIP126.Classical.Adams.adamsTensorBoundary,KIP126.Classical.Adams.adamsTensorBoundary_comparison,KIP126.Classical.Adams.adamsNextHomologyTensorEquiv_normalized,KIP126.Classical.Adams.adamsTensorBoundary_normalization,KIP126.Classical.Adams.adamsTensorCoaction_next_coproduct}
  \uses{prop:h6-next-stage-coaction,prop:h6-kunneth-tower-recurrence,def:h6-kunneth-diagonal-input}
  \leanok
  Let $e_X$ be the constructed reduced-tensor coordinate on next-stage
  homology and $\iota$ its inclusion into $\mathcal A\otimes H_*X$.
  Put $q_X=b_X\kappa_X^{-1}$, using the actual Adams boundary.
  Without diagonal coherence one already has the normalization formula
  \[
    \iota e_Xq_X(w)=w-\rho_X\operatorname{aug}(w).
  \]
  Under both K\"unneth suspension and diagonal compatibility, as well as
  the explicit unit/suspension and tensor exactness hypotheses used above,
  \[
    \rho_{\operatorname{fib}\eta_X}(y)
       =(1\otimes q_X)(\Delta\otimes1)\iota e_X(y).
  \]
  The map $1\otimes q_X$ lowers the total degree by one. This applies to
  every actual tower stage, without a supplied tower coaction recurrence.
  A compatible Milnor basis and the unit-coordinate comparisons are still
  needed to expand this into the specified polynomial cobar differential.
\end{proposition}
\begin{proof}
  The earlier normalization theorem gives the first formula after
  applying K\"unneth; multiplication compatibility identifies the action
  with tensor augmentation. For the second, use the normalized lift of
  $y$ in the preceding next-stage coaction theorem. Diagonal coherence
  replaces its free-object coaction. The equality $q_X\kappa_X=b_X$
  and the composition law for degree-lowering tensor maps finish the proof.
  \leanok
\end{proof}

\begin{proposition}[All-filtration coordinates in genuine reduced cooperations]
  \label{prop:h6-iterated-cooperation-coordinates}
  \lean{KIP126.StableHomotopy.Cohomology.iteratedReducedCooperations,KIP126.Classical.Adams.adamsTowerHomologyIteratedEquiv,KIP126.Classical.Adams.adamsTowerHomologyIteratedEquiv_succ,KIP126.Classical.Adams.adamsFirstPageIteratedEquiv,KIP126.Classical.Adams.adamsFirstPageIteratedEquiv_apply}
  \uses{prop:h6-kunneth-tower-recurrence,prop:h6-first-page-homology}
  \leanok
  Under the same explicit hypotheses, set $N_0(n)=H_n(X;\mathbb F_2)$ and
  \[
    N_{s+1}(n)=\bigoplus_i\overline{\mathcal A}_i\otimes N_s(n+1-i).
  \]
  There are constructed linear isomorphisms
  $H_n(T_s;\mathbb F_2)\cong N_s(n)$ and
  $E_1^{s,t}(X)\cong N_s(t-s)$ for every $s\geq0$ and integer $t$.
  These concern the actual first quotient page, not the out-of-range
  first page of the internal sequence starting at $r_0=2$.
\end{proposition}
\begin{proof}
  Induct on $s$, using the identity at zero and the actual fiber recurrence
  at the successor. Transport each remaining tensor factor through the
  induction isomorphism. Finally compose with the proved first-page layer
  homology isomorphism, retaining the Adams degree $n=t-s$.
  \leanok
\end{proof}

\begin{definition}[Independent reduced cooperation step]
  \label{def:h6-reduced-cobar-step}
  \lean{KIP126.StableHomotopy.Cohomology.mod2OuterUnitMap,KIP126.StableHomotopy.Cohomology.mod2CobarDifference,KIP126.StableHomotopy.Cohomology.mod2CobarStep}
  \uses{def:h6-homology-kunneth-input,prop:h6-graded-tensor-kernel,prop:h6-coaction-identities}
  \leanok
  Write $V_X=(\eta_{H\wedge X})_*$ for outer-unit insertion and
  $U_X=(H\wedge\eta_X)_*$ for inner-unit insertion. Define
  \[
    c_{X,n}:H_n(X;\mathbb F_2)
       \longrightarrow(\overline{\mathcal A}\otimes H_*(X;\mathbb F_2))_n
  \]
  as the unique reduced tensor whose unreduced image is
  $\kappa_{X,n}(V_Xx-U_Xx)$.
  Both unit insertions split the free action $A_X$, so this difference
  has zero tensor augmentation. This construction uses the specified
  ring and homology K\"unneth input, but no Adams page, differential,
  cofiber exactness, or unit/shift compatibility.
  It is not yet a formula in Milnor polynomial generators.
\end{definition}

\begin{proposition}[Actual first differential equals the reduced step]
  \label{prop:h6-first-differential-cobar-step}
  \lean{KIP126.StableHomotopy.Cohomology.mod2UnitNatTrans,KIP126.StableHomotopy.Cohomology.mod2Unit_shift,KIP126.StableHomotopy.connectingHomomorphism_naturality,KIP126.Classical.Adams.adamsHomologyD1_eq_boundary_outerUnit,KIP126.Classical.Adams.adamsHomologyD1_eq_cobarStep,KIP126.Classical.Adams.adamsPageD_one_cobarStep}
  \uses{def:h6-reduced-cobar-step,prop:h6-kunneth-tower-recurrence,prop:h6-first-differential-homology}
  \leanok
  In addition to the recurrence hypotheses, assume explicitly that the
  natural transformation $\eta:\mathrm{id}\to H\wedge-$ commutes with
  integer suspension. This is extra coherence data, not a consequence
  here of the functor $H\wedge-$ commuting with suspension.
  Under the proved next-stage homology isomorphism,
  \[
    H_{n-1}(\operatorname{fib}\eta_X;\mathbb F_2)
       \cong(\overline{\mathcal A}\otimes H_*(X;\mathbb F_2))_n,
  \]
  the actual resolution differential $\eta_{\operatorname{fib}\eta_X,*}
  \partial_X(x)$ maps to $c_{X,n}(x)$.
  For $X=T_s$ and $n=t-s$, this gives the actual quotient-page
  $d_1:E_1^{s,t}\to E_1^{s+1,t}$ after the existing first-page
  homology comparisons and the reindexing $n-1=t-(s+1)$.
\end{proposition}
\begin{proof}
  Naturality of the unit and its suspension coherence identify the last
  square from the unit-fiber triangle to its smashed triangle.
  Naturality of connecting homomorphisms therefore gives
  $\eta_{\operatorname{fib}\eta_X,*}\partial_X(x)=b_X(V_Xx)$.
  Apply the already proved boundary coordinate formula
  $\kappa(z-U_XA_Xz)$ and the outer-unit retraction $A_XV_X=1$.
  Injectivity of reduced-tensor inclusion gives equality with the
  independently defined step. No first-differential formula is supplied
  as an input field.
  \leanok
\end{proof}

\begin{definition}[Actual cooperation unit and K\"unneth unit coherence]
  \label{def:h6-kunneth-unit-input}
  \lean{KIP126.StableHomotopy.Cohomology.cooperationUnit,KIP126.StableHomotopy.Cohomology.cooperationTensorUnit,KIP126.StableHomotopy.Cohomology.Mod2KunnethUnitCompatible}
  \uses{def:h6-homology-kunneth-input,def:h6-reduced-cobar-step}
  \leanok
  The degree-zero element $1_{\mathcal A}$ is the specified class
  $1\in\pi_0H$ followed by $\eta_H:H\to H\wedge H$.
  For a graded module $V$, the map $u_V:V_n\to(\mathcal A\otimes V)_n$
  sends $x$ to the elementary tensor $1_{\mathcal A}\otimes x$ in the
  zero cooperation-degree summand. The degree cast is explicit.
  The predicate \texttt{Mod2KunnethUnitCompatible} requires
  $\kappa_X V_X=u_{H_*X}$ for every spectrum $X$, where $V_X$ is the actual
  outer-unit insertion. It mentions no Adams page, differential, or Milnor
  basis. No witness for the fixed foundation is supplied or postulated.
\end{definition}

\begin{proposition}[Counit and group-like properties of the actual unit]
  \label{prop:h6-cooperation-unit-laws}
  \lean{KIP126.StableHomotopy.Cohomology.cooperationCounit_unit,KIP126.StableHomotopy.Cohomology.cooperationTensorUnit_augmentation,KIP126.StableHomotopy.Cohomology.cooperationTensorUnit_injective,KIP126.StableHomotopy.Cohomology.cooperationTensorUnit_naturality,KIP126.StableHomotopy.Cohomology.cooperationTensorDiagonal_unit}
  \uses{def:h6-kunneth-unit-input,def:h6-kunneth-diagonal-input}
  \leanok
  Tensor augmentation retracts $u_V$, so $u_V$ is injective, and it is
  natural in graded linear maps. The actual cooperation counit sends
  $1_{\mathcal A}$ to the specified class $1\in\pi_0H$.
  Under K\"unneth unit coherence,
  $\Delta(1_{\mathcal A})=1_{\mathcal A}\otimes1_{\mathcal A}$.
  This last statement does not require diagonal or suspension coherence.
\end{proposition}
\begin{proof}
  The ring unit law proves the cooperation counit calculation. Apply it
  to the zero summand of the graded tensor augmentation. Naturality of
  the actual unit identifies inner and outer insertion on this class;
  K\"unneth unit coherence gives the displayed group-like formula.
  \leanok
\end{proof}

\begin{proposition}[Unit-minus-coaction formula and the first-cycle criterion]
  \label{prop:h6-first-cycle-coaction}
  \lean{KIP126.StableHomotopy.Cohomology.mod2CobarStep_unit_sub_coaction,KIP126.StableHomotopy.Cohomology.mod2CobarStep_eq_zero_iff,KIP126.Classical.Adams.adamsHomologyD1_unit_sub_coaction,KIP126.Classical.Adams.adamsHomologyD1_eq_zero_iff,KIP126.Classical.Adams.adamsPageD_one_eq_zero_iff}
  \uses{def:h6-kunneth-unit-input,prop:h6-first-differential-cobar-step}
  \leanok
  Under the explicit K\"unneth unit coherence, the independently
  constructed reduced step has unreduced representative
  $1_{\mathcal A}\otimes x-\rho_X(x)$. Thus the step vanishes exactly when
  $\rho_X(x)=1_{\mathcal A}\otimes x$.
  With the exactness and unit/suspension hypotheses of the actual
  first-differential comparison, this is the representative of the actual
  $d_1$ in next-stage reduced-tensor coordinates. For the actual quotient
  page, writing $y$ for the existing first-page homology coordinate of $x$,
  \[
    d_1(x)=0\quad\Longleftrightarrow\quad
    \rho_{T_s}(y)=1_{\mathcal A}\otimes y.
  \]
  No Milnor coordinates, Lin table, K\"unneth diagonal coherence, or
  K\"unneth suspension coherence are required for this first-cycle
  criterion. It is not a criterion for higher-page or permanent survival.
\end{proposition}
\begin{proof}
  Apply K\"unneth unit coherence to the outer-unit term of the already
  proved two-unit difference. Reduced-tensor inclusion is injective;
  the next-stage and first-page homology comparisons are isomorphisms.
  They therefore detect zero, including after the original degree casts.
  \leanok
\end{proof}

\begin{proposition}[Sphere coefficient homology comes from the actual unit]
  \label{prop:h6-sphere-unit-image}
  \lean{KIP126.StableHomotopy.Cohomology.adamsUnit_sphere_unitor,KIP126.StableHomotopy.Cohomology.mod2SphereHomology_unit_surjective,KIP126.StableHomotopy.Cohomology.mod2CoactionMap_unit_image,KIP126.StableHomotopy.Cohomology.mod2CoactionMap_sphere}
  \uses{def:h6-em-object,prop:h6-coaction-identities}
  \leanok
  For every integer $n$, the actual unit induces a surjection
  $\pi_nS^0\to H_n(S^0;\mathbb F_2)$. Consequently inner and outer unit
  insertions agree on all sphere coefficient classes. Neither a ring
  structure nor K\"unneth is needed for these statements.
\end{proposition}
\begin{proof}
  The right tensor unitor identifies the target with $\pi_nH$. For
  $n\ne0$ it vanishes. For $n=0$, its two classes are zero and the image
  of the identity of the sphere (with the shift-zero identification).
  Naturality of the actual unit identifies inner and outer insertions
  on every unit image, hence on every sphere coefficient class.
  \leanok
\end{proof}

\begin{proposition}[The sphere is the derived initial coaction for the recurrence]
  \label{prop:h6-sphere-initial-coaction}
  \lean{KIP126.StableHomotopy.Cohomology.mod2TensorCoaction_sphere,KIP126.StableHomotopy.Cohomology.mod2CobarStep_sphere_zero}
  \uses{prop:h6-sphere-unit-image,def:h6-kunneth-unit-input,def:h6-reduced-cobar-step}
  \leanok
  Under the homology-level K\"unneth input and unit coherence, the actual
  initial coaction is $\rho_{S^0}(x)=1_{\mathcal A}\otimes x$.
  The independent reduced step on sphere homology vanishes even without
  unit coherence. Thus the initial coaction for the actual tower
  recurrence is derived, not a separately supplied sphere assumption.
\end{proposition}
\begin{proof}
  Replace inner insertion by outer insertion using the preceding
  proposition, then apply K\"unneth unit coherence. For the reduced step,
  their difference is already zero before applying K\"unneth.
  \leanok
\end{proof}

These conditional constructions do not remove \texttt{standardMilnor}:
the fixed ring, exactness, unit/shift coherence and K\"unneth witnesses,
the identification of genuine reduced cooperations with the Milnor
polynomial algebra, and the diagonal/coaction comparisons needed to
express this step as the Milnor cobar differential are still missing.
No new fixed axiom is introduced, and the permanent-cycle Solution remains open.

\begin{proposition}[Page passage for the constructed tower]
  \label{prop:h6-tower-page-passage}
  \lean{KIP126.Classical.Adams.adamsCycles_succ_le}
  \lean{KIP126.Classical.Adams.adamsNextPageToHomology}
  \lean{KIP126.Classical.Adams.adamsNextPageToHomology_bijective}
  \lean{KIP126.Classical.Adams.adamsPageHomologyIso}
  \uses{def:h6-tower-differential}
  \notready
  Sending an $(r+1)$-cycle to the homology class of its image in $E_r$
  induces an isomorphism $E_{r+1}^{s,t}\cong H^{s,t}(E_r,d_r)$.
\end{proposition}
\begin{proof}
  A lift one stage farther gives an $r$-cycle with zero differential.
  The next-page boundaries become the current-page differential boundaries.
  Conversely, exactness adjusts the lift of an $r$-cycle with zero
  differential so that it lifts one stage farther.  The same exactness
  calculation identifies the kernel with the next-page boundaries.
  Invert this specified quotient map to obtain page passage in the
  direction $H(E_r)\to E_{r+1}$.
\end{proof}

\begin{definition}[Internal nested subobjects of the constructed tower]
  \label{def:h6-tower-internal-data}
  \lean{KIP126.Classical.Adams.adamsTowerSSData,KIP126.Classical.Adams.adamsTowerSSDataPageIso}
  \uses{def:h6-tower-pages}
  \leanok
  For any chosen unit and object in the stated stable context, take the
  ambient module to be $Z_2^{s,t}$ of the constructed tower. Internal stage
  $n$ consists of the submodules $Z_{n+2}^{s,t}$ and $B_{n+2}^{s,t}$ inside
  that ambient module. At infinity take the intersection of the cycles
  and the increasing union of the boundaries. These define an internal
  SSData object, and its finite quotient page is canonically isomorphic
  to the tower quotient $E_{n+2}^{s,t}$. No independent spectral-sequence
  model or Lin-table assumption enters this generic construction.
\end{definition}

\begin{proposition}[Representative and successor laws for the internal tower]
  \label{prop:h6-tower-internal-laws}
  \lean{KIP126.Classical.Adams.adamsTowerSSDataPageIso_π,KIP126.Classical.Adams.adamsTowerInternalD_comparison,KIP126.Classical.Adams.adamsTowerInternalD_comp,KIP126.Classical.Adams.adamsTowerInternalD_kernel,KIP126.Classical.Adams.adamsTowerInternalD_image}
  \uses{def:h6-tower-internal-data,prop:h6-tower-differential-laws}
  \leanok
  The quotient comparison preserves actual representatives. Transporting
  the tower differential through it gives a square-zero internal map.
  Its kernel is the image of the next cycle subobject, and its image is
  the image of the next boundary subobject on the current quotient page.
\end{proposition}
\begin{proof}
  Compare the categorical cokernel projection with the ordinary module
  quotient. A representative has zero differential exactly when its
  connecting image lifts one stage farther in the tower. The exact-couple
  identities likewise characterize differential images as next boundaries.
  Transport these elementwise facts through the specified quotient map.
  \leanok
\end{proof}

\begin{definition}[Internal spectral sequence assembled from the tower]
  \label{def:h6-tower-internal-sequence}
  \lean{KIP126.Classical.Adams.adamsTowerInternalSpectralSequence}
  \uses{prop:h6-tower-internal-laws}
  \leanok
  Display the internal tower quotients from page two, with differential
  degree $(r,r-1)$. The representative, square-zero, kernel, and image laws
  above supply all fields of KIP126's internal spectral-sequence structure.
  This construction does not import the Mathlib spectral-sequence type.
\end{definition}

\begin{definition}[Actual first cycles as internal second-page classes]
  \label{def:h6-first-cycle-internal-class}
  \lean{KIP126.Classical.Adams.adamsTowerE2OfFirstCycle}
  \uses{def:h6-tower-internal-sequence,prop:h6-tower-page-passage}
  \leanok
  An actual first-page cycle determines an element of the existing internal
  $E_2$ page: choose a $Z_2$ lift, take its actual tower quotient class,
  and use the inverse representative-preserving page comparison.
  This is not a new spectral sequence or an assumed page identification.
\end{definition}

\begin{proposition}[Independence and exact vanishing of the internal class]
  \label{prop:h6-first-cycle-internal-class}
  \lean{KIP126.Classical.Adams.adamsNextCycle_exists_of_differential_eq_zero,KIP126.Classical.Adams.adamsNextCycle_quotient_eq_of_page_eq,KIP126.Classical.Adams.adamsTowerE2OfFirstCycle_comparison,KIP126.Classical.Adams.adamsTowerE2OfFirstCycle_eq_zero_iff}
  \uses{def:h6-first-cycle-internal-class,prop:h6-tower-internal-laws}
  \leanok
  At every finite page, a cycle has a next-cycle representative, and two
  such representatives of the same current class have the same next-page
  quotient. In particular, the internal $E_2$ construction agrees with
  any actual $Z_2$ lift and vanishes exactly when that lift belongs to
  the actual $B_2$ submodule.
\end{proposition}
\begin{proof}
  Lift through the quotient map and use the proved differential-kernel
  criterion. Two lifts differ by a current boundary, hence by a next
  boundary. Transport the quotient equality through the existing page
  isomorphism; its injectivity gives the exact vanishing criterion.
  \leanok
\end{proof}

\begin{proposition}[Incoming-image test for internal second-page vanishing]
  \label{prop:h6-internal-first-cycle-image}
  \lean{KIP126.Classical.Adams.adamsNextBoundary_iff_is_differential,KIP126.Classical.Adams.adamsTowerE2OfFirstCycle_eq_zero_iff_is_differential}
  \uses{prop:h6-first-cycle-internal-class,prop:h6-tower-internal-laws}
  \leanok
  For every $r\geq1$, a next-page cycle belongs to $B_{r+1}$ if and only
  if its current-page class is an incoming $d_r$ boundary. Consequently
  the internal $E_2$ class of an $E_1$ cycle $z$ is zero exactly when
  $z$ lies in the incoming $d_1$ image. The latter criterion has no
  chosen $Z_2$ representative in its statement.
\end{proposition}
\begin{proof}
  Use the actual next-cycle-to-homology map: its kernel is precisely
  the next boundaries, and zero in current homology means membership
  in the incoming image. Combine with the internal quotient vanishing
  criterion for any next-cycle lift.
  \leanok
\end{proof}

\begin{proposition}[All-page comparison of the two tower presentations]
  \label{prop:h6-tower-internal-comparison}
  \lean{KIP126.Classical.Adams.adamsTowerPageComparison,KIP126.Classical.Adams.adamsTowerPageComparison_differential,KIP126.Classical.Adams.adamsTowerPageComparison_passage}
  \uses{def:h6-tower-internal-sequence,prop:h6-tower-page-passage}
  \leanok
  On every page $r\ge2$, the representative-preserving quotient isomorphism
  identifies the internal tower page with the Mathlib tower page. It commutes
  with the tower differential. Two internal classes satisfy the specified
  next-page representative relation if and only if their images correspond
  under the Mathlib homology-to-next-page map. The comparison is constructed
  for every unit and object in the stated stable context; it does not use
  a Lin presentation, Milnor coordinates, or a permanence assumption.
\end{proposition}
\begin{proof}
  Reindex internal stage $n$ as classical page $n+2$ and use the proved
  differential comparison. A next-cycle representative gives both the
  current quotient class and the next quotient class. The explicitly
  constructed homology isomorphism sends its homology class to that same
  next quotient class. Conversely, surjectivity of the quotient projection
  and the kernel calculation recover such a representative.
  \leanok
\end{proof}

\begin{proposition}[A common representative for compatible tower classes]
  \label{prop:h6-tower-common-representative}
  \lean{KIP126.Classical.Adams.adamsNextPage_representative_independence,KIP126.Classical.Adams.adamsPages_common_representative,KIP126.Classical.Adams.adamsTower_nonzeroSurvival_iff,KIP126.Classical.Adams.adamsTower_nonzeroSurvival_iff_compatible}
  \uses{def:h6-tower-internal-sequence,prop:h6-tower-internal-laws}
  \leanok
  A compatible family of classes in the actual finite tower quotients
  has a single representative in $Z_2$ which belongs to every later cycle
  submodule. Consequently internal nonzero survival is equivalent to a
  compatible family of nonzero classes in all these quotients, with the
  prescribed initial class. No convergence or finiteness hypothesis is needed.
\end{proposition}
\begin{proof}
  Choose a representative of the initial class. If it agrees with a
  next-cycle representative on the current page, their difference is a
  current boundary. Every boundary lies in every later cycle submodule,
  and current boundaries are next boundaries. Thus the original representative
  remains a next cycle and gives the prescribed next class. Induct on the
  finite page. Membership in $Z_\infty$ is membership in all finite cycles;
  membership in $B_\infty$ is membership in some finite boundary submodule.
  These two facts identify nonzero classes at infinity with the stated
  nonzero compatible family. The entire argument is internal to the tower's
  SSData and does not import the Mathlib spectral-sequence type.
  \leanok
\end{proof}

\begin{proposition}[Comparison of the two nonzero survival predicates]
  \label{prop:h6-tower-survival-comparison}
  \lean{KIP126.Classical.Adams.adamsTower_survival_comparison}
  \uses{prop:h6-tower-common-representative,prop:h6-tower-internal-comparison,def:h6-pagewise-permanence}
  \leanok
  For every E$_2$ class of the constructed tower, internal nonzero survival
  is equivalent to standard pagewise permanence of its image under the same
  tower comparison. This is a generic theorem, independent of the Lin table,
  Milnor coordinates, and any assertion of permanence for a specified class.
\end{proposition}
\begin{proof}
  Reindex the standard trajectory's iterated pages as $n+2$. The actual
  homology-to-next-page isomorphism identifies its cycle and passage conditions
  with the next-cycle representative relation. Apply the common-representative
  result and check that its initial quotient map is exactly the E$_2$ component
  of the fixed tower comparison.
  \leanok
\end{proof}

\begin{definition}[Constructed sphere Adams sequence]
  \label{def:h6-constructed-sphere-sequence}
  \lean{KIP126.Classical.Adams.adamsTowerSpectralSequence}
  \lean{KIP126.Classical.Adams.mod2SphereAdams}
  \uses{prop:h6-tower-page-passage,def:h6-em-object}
  \notready
  Specialize the tower to $X=S^0$ and display its quotient-page spectral
  sequence from $E_2$, with differential bidegree $(r,r-1)$.
  Its underlying abelian groups may be regarded as integer modules.
  Given the ring and additive-tensor structures, the mod--$2$ scalars are
  derived in Proposition~\ref{prop:h6-derived-mod2-pages}. The Milnor
  coordinates below specify the further cobar comparison; no
  convergence-to-an-abutment assertion is required to state pagewise
  nonzero permanence.
\end{definition}

% SOURCE: J. M. Boardman, Stable operations in generalized cohomology,
% p. 208, (7.4), verified at
% https://www.sas.rochester.edu/mth/sites/doug-ravenel/otherpapers/boardman-8.pdf
\begin{definition}[Normalized Milnor cobar cochains]
  \label{def:h6-milnor-cobar}
  \lean{KIP126.Steenrod.Milnor.TensorPower}
  \lean{KIP126.Steenrod.Milnor.weight}
  \lean{KIP126.Steenrod.Milnor.coproductGenerator}
  \lean{KIP126.Steenrod.Milnor.cochains}
  \lean{KIP126.Steenrod.Milnor.differential}
  \lean{KIP126.Steenrod.Milnor.cup}
  \uses{def:mathlib-f2}
  \leanok
  Set $\mathcal A_*=\mathbb F_2[\xi_1,\xi_2,\ldots]$, with
  $|\xi_i|=2^i-1$, $\xi_0=1$, and
  \[
    \Delta\xi_n=\sum_{i=0}^n\xi_{n-i}^{2^i}\otimes\xi_i.
  \]
  Define $C^{s,t}$ as the internal-degree-$t$ polynomials in $s$ tensor
  slots whose augmentation in every slot is zero.  The cobar differential
  is the sum of the two end coaugmentations and the coproduct in each slot.
  The product $C^{s,t}\otimes C^{s',t'}\to C^{s+s',t+t'}$ concatenates slots.
\end{definition}

\begin{proposition}[Monomial basis and word tensor decomposition]
  \label{prop:h6-cobar-monomial-basis}
  \lean{KIP126.Steenrod.Milnor.mem_cochains_iff_monomials,KIP126.Steenrod.Milnor.cochainsMonomialBasis,KIP126.Steenrod.Milnor.cochainsMonomialBasis_val,KIP126.Steenrod.Milnor.cochainsWordEquiv,KIP126.Steenrod.Milnor.wordTensorEquiv,KIP126.Steenrod.Milnor.wordTensorEquiv_single}
  \uses{def:h6-milnor-cobar}
  \leanok
  The normalized homogeneous cochains have the specified monomial basis:
  an exponent belongs to the basis exactly when it has weight $t$ and
  uses every tensor slot. Let $P_i$ be the nonconstant Milnor monomials of
  integer degree $i$, and let $W_s(t)$ be the length-$s$ words in these
  monomials with total degree $t$. Then there are constructed isomorphisms
  \[
    C^{s,t}\cong\mathbb F_2^{(W_s(t))},\qquad
    \bigoplus_{i\in\mathbb Z}
       \mathbb F_2^{(P_i)}\otimes\mathbb F_2^{(W_s(t-i))}
       \cong\mathbb F_2^{(W_{s+1}(t))}.
  \]
  Superscripts in parentheses denote finite support. The first map takes
  actual polynomial coefficients. The second sends the tensor of two
  basis vectors to the prepended word with multiplied coefficient.
  Integer degrees are retained, so negative-degree word sets are empty.
\end{proposition}
\begin{proof}
  Augmenting a slot deletes precisely the monomials using that slot; all
  remaining monomials retain their coefficients, so there is no hidden
  cancellation. Combine this with weighted homogeneity and restrict the
  polynomial coefficient equivalence. Curry exponents into individual
  slots and split a nonempty word into its head and tail. Finally apply
  the finite-support tensor and direct-sum equivalences, checking their
  values on basis vectors.
  \leanok
\end{proof}

\begin{definition}[Below-page reduced Milnor basis input]
  \label{def:h6-reduced-milnor-basis-input}
  \lean{KIP126.StableHomotopy.Cohomology.Mod2ReducedMilnorBasis}
  \uses{prop:h6-graded-tensor-kernel,prop:h6-cobar-monomial-basis}
  \leanok
  A \texttt{Mod2ReducedMilnorBasis} specifies for each integer $i$ a basis
  of the actual reduced cooperation group
  $\ker(\pi_i(H\wedge H)\xrightarrow{\mu_*}\pi_iH)$ indexed by $P_i$.
  It has no Adams page, page coordinate, differential, or permanence field.
  The input structure is defined, but its existence on the fixed foundation
  is not proved or postulated here. A graded basis alone does not supply
  the Milnor coproduct compatibility.
\end{definition}

\begin{definition}[Full single-slot monomials]
  \label{def:h6-full-milnor-monomials}
  \lean{KIP126.Steenrod.Milnor.MilnorMonomial,KIP126.Steenrod.Milnor.zeroMilnorMonomial,KIP126.Steenrod.Milnor.milnorMonomialEquivPositive,KIP126.Steenrod.Milnor.milnorZeroCoefficientsEquiv}
  \uses{prop:h6-cobar-monomial-basis}
  \leanok
  Let $M_n$ consist of finite exponent vectors $d$ of integer weight $n$,
  including the zero exponent vector. The unique degree-zero monomial
  is the constant monomial. For $n\ne0$, forgetting or adding the
  nonzero-exponent condition gives $M_n\cong P_n$.
  The degree-zero coefficient space is identified with $\mathbb F_2$
  by evaluation at the constant monomial. Negative degrees are retained,
  not replaced with their nonnegative truncation.
\end{definition}

\begin{proposition}[Complete cooperation coordinates with the actual unit]
  \label{prop:h6-full-milnor-coordinates}
  \lean{KIP126.Steenrod.Milnor.slotWeight_eq_zero_iff,KIP126.Steenrod.Milnor.positiveMonomial_degree_pos,KIP126.StableHomotopy.Cohomology.cooperationCounit_surjective,KIP126.StableHomotopy.Cohomology.reducedCooperations_zero_subsingleton,KIP126.StableHomotopy.Cohomology.cooperationCounitF2_zero_bijective,KIP126.StableHomotopy.Cohomology.cooperationMilnorEquiv,KIP126.StableHomotopy.Cohomology.cooperationMilnorEquiv_zero_coefficient,KIP126.StableHomotopy.Cohomology.cooperationMilnorEquiv_unit,KIP126.StableHomotopy.Cohomology.cooperationMilnorEquiv_of_ne,KIP126.StableHomotopy.Cohomology.cooperationMilnorEquiv_reduced_basis}
  \uses{def:h6-full-milnor-monomials,def:h6-reduced-milnor-basis-input,prop:h6-cooperation-unit-laws}
  \leanok
  Given only the existing reduced Milnor basis, the ring structure and
  additive tensor, construct for every integer $n$ a linear isomorphism
  \[
    \Phi_n:\pi_n(H\wedge H)\cong\mathbb F_2^{(M_n)}.
  \]
  It sends the actual cooperation unit to the constant monomial with
  coefficient one, and every supplied reduced basis vector to its original
  nonconstant monomial. In degree zero its sole coefficient is exactly
  the actual counit followed by the specified $\pi_0H\cong\mathbb F_2$.
  In other degrees the actual counit is zero and the comparison uses
  precisely the original reduced basis. No full-basis existence input,
  K\"unneth comparison, or Adams page is used.
  This does not prove preservation of multiplication or the Milnor
  coproduct; it supplies the unit- and counit-compatible coordinates
  needed to formulate that remaining coproduct comparison.
\end{proposition}
\begin{proof}
  All single-variable weights are positive, so weight zero forces the
  zero exponent vector. Hence $P_0$ is empty and its basis forces the
  actual counit kernel in degree zero to vanish. The actual unit splits
  the counit, which is therefore a linear isomorphism in that degree.
  Use its coefficient to fill the unique constant monomial. Away from
  zero, $\pi_nH=0$, so the full cooperation group equals its counit kernel;
  apply the original basis coordinates and the identification $P_n\cong M_n$.
  The ring unit law proves the asserted constant coefficient, while
  basis-coordinate evaluation proves the reduced-vector formula.
  \leanok
\end{proof}

\begin{definition}[Below-page Milnor coproduct comparison]
  \label{def:h6-milnor-coproduct-input}
  \lean{KIP126.StableHomotopy.Cohomology.milnorMonomialPolynomial,KIP126.StableHomotopy.Cohomology.cooperationMilnorPolynomial,KIP126.StableHomotopy.Cohomology.cooperationTensorMilnorPolynomial,KIP126.StableHomotopy.Cohomology.Mod2MilnorCoproductCompatible}
  \uses{prop:h6-full-milnor-coordinates,def:h6-kunneth-diagonal-input}
  \leanok
  Realize the full coordinates as a linear map $P_n$ into the existing
  one-slot polynomial algebra. On the tensor square, let $Q_n$ concatenate
  the two polynomial coordinates in adjacent slots. The remaining Milnor
  coproduct condition is, for each supplied reduced basis vector $b_d$,
  \[
    Q_n\Delta(b_d)=\operatorname{splitSlot}_0(m_d).
  \]
  Here $\Delta$ is actual middle-unit insertion followed by K\"unneth,
  and $m_d$ is the corresponding monomial. Thus the two sides are
  independently defined. The quantifier runs over all reduced basis
  monomials, not just algebra generators. This definition neither asserts
  existence of a compatible basis on the fixed foundation nor prescribes
  an Adams page differential.
\end{definition}

\begin{proposition}[Extend the coproduct comparison from the basis]
  \label{prop:h6-milnor-coproduct-linear-extension}
  \lean{KIP126.StableHomotopy.Cohomology.cooperation_linearMap_ext,KIP126.StableHomotopy.Cohomology.cooperationMilnorPolynomial_unit,KIP126.StableHomotopy.Cohomology.cooperationMilnorPolynomial_reduced_basis,KIP126.StableHomotopy.Cohomology.cooperationTensorMilnorPolynomial_lof_tmul,KIP126.StableHomotopy.Cohomology.cooperationTensorMilnorPolynomial_unit,KIP126.StableHomotopy.Cohomology.cooperationMilnorCoproduct_unit,KIP126.StableHomotopy.Cohomology.cooperationMilnorCoproduct,KIP126.StableHomotopy.Cohomology.mod2MilnorCoproductCompatible_iff}
  \uses{def:h6-milnor-coproduct-input,def:h6-kunneth-unit-input,prop:h6-cooperation-unit-laws}
  \leanok
  Two families of linear maps on actual cooperations agree if they agree
  on the actual unit and the reduced basis. Under K\"unneth unit coherence,
  the constant case of the coproduct formula follows from the actual unit
  law. Consequently the reduced-basis condition is equivalent to
  \[
    Q_n\Delta(a)=\operatorname{splitSlot}_0(P_n(a))
    \qquad\text{for every integer }n\text{ and }a\in\pi_n(H\wedge H).
  \]
  This is a conditional linear-extension theorem, not a proof of the
  coproduct input or of the full cobar differential comparison.
\end{proposition}
\begin{proof}
  Pull the two linear maps back along the full coordinate isomorphism
  $\Phi_n$. Finite-support linear maps are determined by their singletons.
  The degree-zero singleton corresponds to the actual unit, and all other
  singletons correspond to the original reduced basis. The coordinate
  map sends the unit to $1$, and the tensor coordinate map sends the unit
  tensor to $1$. Apply the previously proved group-like unit identity and
  then the singleton criterion. The reverse implication is specialization
  to the reduced basis.
  \leanok
\end{proof}

\begin{definition}[Derived basis of the cooperation tensor square]
  \label{def:h6-milnor-tensor-basis}
  \lean{KIP126.StableHomotopy.Cohomology.MilnorPair,KIP126.StableHomotopy.Cohomology.milnorPairExponents,KIP126.StableHomotopy.Cohomology.cooperationMilnorBasis,KIP126.StableHomotopy.Cohomology.cooperationTensorMilnorBasis}
  \uses{prop:h6-full-milnor-coordinates,def:h6-milnor-coproduct-input}
  \leanok
  In total degree $n$, index tensor basis vectors by
  $\coprod_{i\in\mathbb Z} M_i\times M_{n-i}$.
  Derive the basis of each genuine cooperation group from its existing
  full coordinates, tensor these bases, and take their direct-sum basis.
  Encode a pair of exponents in slots zero and one. There is no additional
  basis-existence assumption.
\end{definition}

\begin{proposition}[Faithful polynomial coordinates and coproduct tests]
  \label{prop:h6-milnor-polynomial-faithful}
  \lean{KIP126.StableHomotopy.Cohomology.milnorPairExponents_injective,KIP126.StableHomotopy.Cohomology.cup_milnorMonomialPolynomial,KIP126.StableHomotopy.Cohomology.cooperationMilnorPolynomial_basis,KIP126.StableHomotopy.Cohomology.cooperationMilnorPolynomial_injective,KIP126.StableHomotopy.Cohomology.cooperationTensorMilnorBasis_apply,KIP126.StableHomotopy.Cohomology.cooperationTensorMilnorPolynomial_basis,KIP126.StableHomotopy.Cohomology.cooperationTensorMilnorPolynomial_injective,KIP126.StableHomotopy.Cohomology.cooperationTensorDiagonal_eq_iff}
  \uses{def:h6-milnor-tensor-basis,prop:h6-milnor-coproduct-linear-extension}
  \leanok
  Both the one-slot map $P_n$ and the tensor map $Q_n$ are injective.
  Assuming the stated unit and Milnor coproduct compatibilities, any
  proposed tensor value $z$ can therefore be tested exactly:
  \[
    \Delta(a)=z\quad\Longleftrightarrow\quad
    \operatorname{splitSlot}_0(P_n(a))=Q_n(z).
  \]
  The polynomial comparison does not silently discard a tensor kernel.
  Surjectivity onto the entire polynomial algebra is neither claimed nor
  needed; $n$ remains a fixed total degree.
\end{proposition}
\begin{proof}
  The one-slot monomial embedding has distinct exponents. In two slots,
  restriction to each slot recovers the two exponent vectors; the weight
  of the first vector also recovers its direct-sum degree $i$.
  Thus all tensor basis vectors map to distinct polynomial monomials.
  Linear independence of polynomial monomials gives both injectivity
  assertions. Apply tensor injectivity and the previously proved full
  coproduct formula to obtain the equivalence.
  \leanok
\end{proof}

\begin{definition}[Right unit tensor]
  \label{def:h6-right-unit-tensor}
  \lean{KIP126.StableHomotopy.Cohomology.cooperationTensorRightUnit}
  \uses{prop:h6-cooperation-unit-laws}
  \leanok
  Define the linear map $a\mapsto a\otimes1$ by inserting the actual
  cooperation unit in the second factor of the degree-$n$ summand, with
  the canonical degree cast $0=n-n$. This construction uses neither
  Milnor coordinates nor K\"unneth.
\end{definition}

\begin{proposition}[The actual primitive cooperation representing the h-six polynomial]
  \label{prop:h6-actual-primitive-cooperation}
  \lean{KIP126.StableHomotopy.Cohomology.cooperationTensorMilnorPolynomial_leftUnit,KIP126.StableHomotopy.Cohomology.cooperationTensorMilnorPolynomial_rightUnit,KIP126.StableHomotopy.Cohomology.cooperationMilnorPolynomial_h6_existsUnique,KIP126.StableHomotopy.Cohomology.cooperationTensorDiagonal_primitive_iff,KIP126.StableHomotopy.Cohomology.cooperationTensorDiagonal_primitive_of_h6Polynomial,KIP126.StableHomotopy.Cohomology.cooperationTensorDiagonal_h6_existsUnique}
  \uses{prop:h6-milnor-polynomial-faithful,def:h6-right-unit-tensor}
  \leanok
  The two actual unit tensors have coordinates given by the left and right
  polynomial insertions. Under the explicit unit and Milnor coproduct
  compatibilities, $\Delta(a)=1\otimes a+a\otimes1$ is equivalent to the
  same primitive identity for $P_n(a)$.
  There is a unique actual cooperation $a\in\pi_{64}(H\wedge H)$ with
  $P_{64}(a)=\xi_1^{64}$; under those compatibilities it is primitive.
  Existence is constructed from the derived full basis, not supplied as
  another input. This is a statement about a cooperation underlying the
  standard one-bar cocycle, not an identification of a permanent Adams
  class or a proof of survival of its square.
\end{proposition}
\begin{proof}
  The exponent vector with exponent $64$ at $\xi_1$ has weight $64$.
  Its basis vector supplies the required cooperation, uniquely by
  injectivity of $P_{64}$. The unit tensor formulas and injectivity of
  $Q_{64}$ reduce the primitive identity to the existing polynomial
  coproduct. In characteristic two, the $64$th power of a sum is the sum
  of the $64$th powers, giving the desired formula.
  \leanok
\end{proof}

\begin{proposition}[The actual tensor boundary is a comodule map]
  \label{prop:h6-tensor-boundary-comodule}
  \lean{KIP126.StableHomotopy.Cohomology.cooperationTensorLowerMap_unit,KIP126.Classical.Adams.adamsTensorCoaction_tensorBoundary,KIP126.Classical.Adams.adamsTensorBoundary_unit}
  \uses{prop:h6-next-stage-coproduct,prop:h6-first-cycle-coaction}
  \leanok
  Write $q_X$ for the actual homology connecting map after inverse
  K\"unneth, lowering degree by one. Under the explicit suspension,
  diagonal, and unit compatibilities, for every tensor $w$,
  \[
    \rho_{\operatorname{fib}(\eta_X)}q_X(w)
      =(1\otimes q_X)(\Delta\otimes1)(w),
    \qquad q_X(1\otimes x)=D_1(x).
  \]
  The first identity is valid before normalization; the second uses the
  original tower first differential. Unit tensors also commute with
  arbitrary degree-minus-one linear maps, with all grading casts retained.
\end{proposition}
\begin{proof}
  Choose the inverse K\"unneth representative of $w$. Apply naturality of
  the actual coaction with the connecting map, then diagonal coherence
  and composition of the right tensor maps. For the unit tensor, unit
  coherence converts it to outer-unit insertion; the previously proved
  boundary formula gives the actual $D_1$. The degree-cast identities
  for the elementary unit tensor follow by transporting equal indices.
  \leanok
\end{proof}

\begin{proposition}[Primitive tensors give cycles at the next tower stage]
  \label{prop:h6-primitive-next-stage-cycle}
  \lean{KIP126.StableHomotopy.Cohomology.cooperationTensorComultiply_primitive,KIP126.Classical.Adams.adamsHomologyD1_tensorBoundary_primitive,KIP126.Classical.Adams.sphereAdamsHomologyD1_tensorBoundary_primitive,KIP126.Classical.Adams.sphereAdamsPageD_one_tensorBoundary_primitive}
  \uses{prop:h6-tensor-boundary-comodule,def:h6-right-unit-tensor,prop:h6-first-cycle-coaction,prop:h6-sphere-unit-image}
  \leanok
  Under the same low-level compatibilities, let $a\in\pi_k(H\wedge H)$
  satisfy $\Delta a=1\otimes a+a\otimes1$ and let
  $x\in H_{n-k}(X;\mathbb F_2)$ satisfy $D_1x=0$.
  Then the actual next-stage class $q_X(a\otimes x)$ has $D_1=0$.
  For $X=S^0$ every coefficient class satisfies the required condition.
  Transferring by the existing first-page homology isomorphism gives
  a class in the actual $E_1^{1,n}$ whose actual quotient-page
  differential to $E_1^{2,n}$ is zero. No Milnor coordinates are used
  in this general statement, and nonzero classes are not asserted here.
\end{proposition}
\begin{proof}
  Reassociation and primitivity expand $(\Delta\otimes1)(a\otimes x)$
  into $1\otimes(a\otimes x)+a\otimes(1\otimes x)$.
  The boundary-comodule identity and $q_X(1\otimes x)=D_1x=0$
  give $\rho(q_X(a\otimes x))=1\otimes q_X(a\otimes x)$.
  The actual first-cycle coaction criterion proves the claim.
  Sphere coefficients are cycles by the already proved unit-surjectivity
  argument. The original page/homology comparison transports this result
  to the quotient-page differential.
  \leanok
\end{proof}

\begin{proposition}[The primitive h-six polynomial gives an actual first-page cycle]
  \label{prop:h6-polynomial-actual-first-cycle}
  \lean{KIP126.Classical.Adams.sphereAdamsPageD_one_h6_tensorBoundary,KIP126.Classical.Adams.sphereAdamsPageD_one_h6_tensorBoundary_existsUnique}
  \uses{prop:h6-primitive-next-stage-cycle,prop:h6-actual-primitive-cooperation}
  \leanok
  Assume the explicit low-level K\"unneth and Milnor compatibilities.
  There is a unique cooperation $a$ with $P_{64}(a)=\xi_1^{64}$.
  For every degree-zero sphere coefficient $x$, the actual first-page
  class obtained from $q_{S^0}(a\otimes x)$ lies in the kernel of
  $d_1:E_1^{1,64}\to E_1^{2,64}$.
  This uses the actual tower differential, not a differential supplied
  by the fixed Milnor page-comparison input. It does not yet identify the
  class with the fixed standard $h_6$ or prove any higher-page permanence.
  Nonvanishing for a nonzero coefficient is established below.
\end{proposition}
\begin{proof}
  The preceding polynomial calculation makes the unique cooperation
  primitive. Apply the general primitive-cycle construction in total
  degree $64$ with first-factor degree $64$. Uniqueness remains exactly
  uniqueness of the single-slot cooperation coordinate.
  \leanok
\end{proof}

\begin{proposition}[Nonzero primitive-polynomial first cycle]
  \label{prop:h6-nonzero-first-cycle}
  \lean{KIP126.StableHomotopy.Cohomology.mod2SphereHomology_exists_ne_zero,KIP126.StableHomotopy.Cohomology.cooperation_ne_zero_of_h6Polynomial,KIP126.StableHomotopy.Cohomology.cooperationTensor_lof_tmul_ne_zero,KIP126.Classical.Adams.adamsTensorBoundary_eq_zero_iff_of_augmentation_eq_zero,KIP126.Classical.Adams.sphereAdamsPageOne_h6_tensorBoundary_ne_zero,KIP126.Classical.Adams.sphereAdamsPageOne_h6_nonzero_cycle_exists}
  \uses{prop:h6-polynomial-actual-first-cycle,prop:h6-homology-normalization,prop:h6-sphere-unit-image}
  \leanok
  Under the same explicit structural hypotheses, choose a nonzero
  degree-zero sphere coefficient $x$. The first-page cycle obtained from
  the cooperation with polynomial $\xi_1^{64}$ is nonzero.
\end{proposition}
\begin{proof}
  The sphere unitor and the specified degree-zero coefficient equivalence
  provide $x\ne0$. A cooperation with nonzero polynomial is nonzero.
  A pure tensor of two nonzero vectors over $\mathbb F_2$ is nonzero,
  as witnessed by a linear functional on the first factor. Its positive
  cooperation degree makes its augmentation zero. The proved normalized
  boundary identity then prevents its image under $q$ from being zero.
  Transport along the actual first-page homology equivalence.
  \leanok
\end{proof}

\begin{proposition}[Nonzero second-page class with polynomial provenance]
  \label{prop:h6-nonzero-second-page}
  \lean{KIP126.Classical.Adams.sphereAdamsPageTwo_nonzero_of_first_cycle,KIP126.Classical.Adams.sphereAdamsPageTwo_h6_nonzero_exists}
  \uses{prop:h6-nonzero-first-cycle,prop:h6-sphere-unit-image}
  \leanok
  Every nonzero first-page sphere cycle in filtration one has a
  $Z_2$ representative whose $E_1$ image is that cycle and whose $E_2$
  quotient class is nonzero. In particular the preceding polynomial
  construction gives such a representative in bidegree $(1,64)$.
  All structural hypotheses remain explicit. This does not identify it
  with the fixed standard $h_6$, prove a multiplicative comparison, or
  establish the permanence of its square.
\end{proposition}
\begin{proof}
  Lift the first-page class to $Z_1$; the zero differential makes the
  same element lie in $Z_2$. If its second-page class vanished, the
  exact-couple boundary theorem would express the original class as an
  incoming $d_1$ from filtration zero. That differential is zero by the
  sphere unit-image theorem, contradicting the original nonvanishing.
  \leanok
\end{proof}

\begin{proposition}[The primitive-polynomial class on the internal page]
  \label{prop:h6-nonzero-internal-second-page}
  \lean{KIP126.Classical.Adams.sphereAdamsInternalE2_h6_nonzero_exists}
  \uses{prop:h6-nonzero-second-page,prop:h6-first-cycle-internal-class}
  \leanok
  Under the explicit structural hypotheses above, the polynomial
  $\xi_1^{64}$ gives a nonzero class on KIP126's internal $E_2^{1,64}$.
  The statement retains its actual $Z_2$ representative, its original
  first-page tensor-boundary class, and the specified cooperation polynomial.
\end{proposition}
\begin{proof}
  Apply the internal first-cycle construction to the proved nonzero
  second-page representative. The exact comparison preserves nonvanishing.
  No fixed Milnor or Lin comparison is invoked.
  \leanok
\end{proof}

\begin{definition}[Twice-applied tensor-boundary representative]
  \label{def:h6-double-tensor-representative}
  \lean{KIP126.Classical.Adams.sphereH6TensorRepresentative,KIP126.Classical.Adams.sphereH6DoubleTensorRepresentative}
  \uses{prop:h6-tensor-boundary-comodule,prop:h6-actual-primitive-cooperation}
  \leanok
  For a degree-$64$ cooperation $a$ and $x\in H_0(S^0;\mathbb F_2)$,
  set $y=q_{S^0}(a\otimes x)\in H_{63}(T_1)$ and
  $w=q_{T_1}(a\otimes y)\in H_{126}(T_2)$.
  These are actual connecting maps, not a prescribed page multiplication.
\end{definition}

\begin{proposition}[The double representative is a nonzero first cycle]
  \label{prop:h6-double-first-cycle}
  \lean{KIP126.Classical.Adams.sphereH6DoubleTensorRepresentative_ne_zero,KIP126.Classical.Adams.sphereAdamsPageOne_h6_double_ne_zero,KIP126.Classical.Adams.sphereH6DoubleTensorRepresentative_d1_zero,KIP126.Classical.Adams.sphereAdamsPageD_one_h6_double_zero,KIP126.Classical.Adams.sphereAdamsDifferential_one_h6_double_zero}
  \uses{def:h6-double-tensor-representative,prop:h6-nonzero-first-cycle,prop:h6-primitive-next-stage-cycle}
  \leanok
  Suppose $P_{64}(a)=\xi_1^{64}$ and $x\ne0$. Under the explicit structural
  compatibilities, the resulting actual $E_1^{2,128}$ class is nonzero
  and has zero $d_1$. This does not assert that it is not an incoming boundary.
\end{proposition}
\begin{proof}
  Apply the primitive-cycle theorem twice. Each tensor has nonzero first
  factor in positive cooperation degree, so the normalized-boundary
  identity proves nonvanishing at each step. Transfer to the original
  quotient page through its homology equivalence.
  \leanok
\end{proof}

\begin{definition}[The double construction on the internal target page]
  \label{def:h6-double-internal-second-page}
  \lean{KIP126.Classical.Adams.sphereH6DoubleInternalE2}
  \uses{prop:h6-double-first-cycle,def:h6-first-cycle-internal-class}
  \leanok
  Send the double first cycle to the existing internal $E_2^{2,128}$
  by the general first-cycle construction. The definition takes the
  low-level structures and compatibilities explicitly.
\end{definition}

\begin{proposition}[Provenance and remaining boundary test for the double class]
  \label{prop:h6-double-internal-representative}
  \lean{KIP126.Classical.Adams.sphereH6DoubleInternalE2_representative,KIP126.Classical.Adams.sphereH6DoubleInternalE2_eq_zero_iff}
  \uses{def:h6-double-internal-second-page,prop:h6-first-cycle-internal-class}
  \leanok
  The double internal class has a genuine $Z_2$ lift of the specified
  double first cycle. It vanishes exactly when such a lift lies in $B_2$.
  Exclusion of that boundary, identification with the fixed standard or
  computed square, and higher-page survival remain open for this construction.
\end{proposition}
\begin{proof}
  Specialize the general first-cycle lifting, comparison, and vanishing
  theorems. No non-boundary or permanence premise is inserted.
  \leanok
\end{proof}

\begin{proposition}[Derived all-filtration Milnor vector-space coordinates]
  \label{prop:h6-derived-milnor-coordinates}
  \lean{KIP126.StableHomotopy.Cohomology.reducedTensorMilnorWordEquiv,KIP126.Classical.Adams.sphereHomologyEmptyWordEquiv,KIP126.Classical.Adams.sphereTowerHomologyWordEquiv,KIP126.Classical.Adams.sphereFirstPageMilnorEquiv,KIP126.Classical.Adams.sphereTowerHomologyWordEquiv_d1}
  \uses{def:h6-reduced-milnor-basis-input,prop:h6-cobar-monomial-basis,prop:h6-kunneth-tower-recurrence,prop:h6-first-differential-cobar-step}
  \leanok
  Given the reduced-cooperation basis, the homology-level K\"unneth input,
  the ring structure, and the tensor exactness and suspension structure
  used in the recurrence, one obtains for every $s\geq0$ and integer $n$
  \[
    H_n(T_s(S^0);\mathbb F_2)\cong\mathbb F_2^{(W_s(n+s))}.
  \]
  Composing with actual first-page homology and the proved polynomial
  coefficient equivalence gives $E_1^{s,t}(S^0)\cong C^{s,t}$ for all
  $s,t\geq0$. No page-coordinate field is used.
  If the unit also commutes with suspension, the resolution differential
  in these word coordinates equals the tensor-word conversion of
  $c_{T_s,n}$ from Definition~\ref{def:h6-reduced-cobar-step}.
  The equality with the explicitly specified Milnor polynomial
  differential is not yet proved: it still requires the compatible
  coproduct/coaction identification. The fixed input witnesses are also
  still missing, so \texttt{standardMilnorCooperations} remains an axiom.
\end{proposition}
\begin{proof}
  In length zero, the tensor unitor and the Eilenberg--Mac Lane property
  identify sphere homology with the empty word in degree zero and zero
  in all other degrees. Induct using the actual next-fiber K\"unneth
  isomorphism, replace each reduced factor using its basis, and prepend
  words by the preceding tensor decomposition. The degree becomes
  $n+s$, hence $t$ on $E_1^{s,t}$. The differential statement follows
  from the already proved boundary/coaction formula, not an additional
  differential input.
  \leanok
\end{proof}

\begin{proposition}[Reduced tensors are the exact coordinates of actual boundaries]
  \label{prop:h6-reduced-boundary-coordinates}
  \lean{KIP126.StableHomotopy.Cohomology.reducedCooperationTensorInclusion_lof_tmul,KIP126.StableHomotopy.Cohomology.cooperationTensorAugmentation_reduced_inclusion,KIP126.Classical.Adams.adamsNextHomologyTensorEquiv_boundary_reduced,KIP126.Classical.Adams.adamsNextHomologyTensorEquiv_boundary_lof_tmul,KIP126.StableHomotopy.Cohomology.reducedTensorMilnorWordEquiv_basis_single}
  \uses{prop:h6-tensor-boundary-comodule,prop:h6-derived-milnor-coordinates}
  \leanok
  For every reduced tensor $w$, the derived next-stage coordinates send
  the actual boundary of its inclusion back to $w$. An elementary reduced
  basis tensor followed by a single word has the concatenated word as
  its coordinates, with the same coefficient. No fixed page coordinates
  or coproduct compatibility are required for these assertions.
\end{proposition}
\begin{proof}
  A reduced tensor has augmentation zero, so the normalization identity
  reduces to the inclusion of $w$. Cancel the injective reduced inclusion.
  The direct-sum tensor coordinate map sends a reduced basis vector to its
  single coefficient and then applies the already constructed word-cons
  equivalence.
  \leanok
\end{proof}

\begin{definition}[Specified h-six words and normalized sphere coefficient]
  \label{def:h6-normalized-word-representatives}
  \lean{KIP126.Steenrod.Milnor.emptyMilnorWord,KIP126.Steenrod.Milnor.h6PositiveMonomial,KIP126.Steenrod.Milnor.h6MilnorWord,KIP126.Steenrod.Milnor.h6SquareMilnorWord,KIP126.Classical.Adams.sphereMilnorUnitCoefficient}
  \uses{prop:h6-derived-milnor-coordinates}
  \leanok
  Specify the empty word, the cooperation monomial $\xi_1^{64}$, its
  length-one word, and the length-two word with two such factors. Normalize
  the sphere coefficient by requiring its degree-zero empty-word coordinate
  to be one, using the existing sphere unitor and specified $\pi_0$ coordinate.
\end{definition}

\begin{proposition}[Exact polynomial coordinates of the constructed representatives]
  \label{prop:h6-exact-polynomial-representatives}
  \lean{KIP126.Steenrod.Milnor.cochainsWordEquiv_h6,KIP126.Steenrod.Milnor.cochainsWordEquiv_h6Square,KIP126.Classical.Adams.sphereMilnorUnitCoefficient_ne_zero,KIP126.Classical.Adams.cooperation_h6_eq_reduced_basis,KIP126.Classical.Adams.sphereTowerHomologyWordEquiv_unitCoefficient,KIP126.Classical.Adams.sphereTowerHomologyWordEquiv_h6,KIP126.Classical.Adams.sphereTowerHomologyWordEquiv_h6_double,KIP126.Classical.Adams.sphereFirstPageMilnorEquiv_h6,KIP126.Classical.Adams.sphereFirstPageMilnorEquiv_h6_double}
  \uses{def:h6-normalized-word-representatives,prop:h6-reduced-boundary-coordinates,def:h6-double-tensor-representative,prop:h6-first-boundary-word-coordinate,prop:h6-word-polynomial-realization}
  \leanok
  The normalized coefficient is nonzero. With this coefficient, the
  derived first-page coordinates send the actual single and double tensor
  boundaries exactly to the existing cochains $[\xi_1^{64}]$ and
  $[\xi_1^{64}\mid\xi_1^{64}]$. This is equality of specified elements,
  not merely agreement of their bidegrees.
\end{proposition}
\begin{proof}
  Polynomial injectivity identifies the cooperation with its reduced
  monomial basis vector. Apply the reduced-boundary coordinate formula
  once and twice, starting with the empty word of coefficient one. The
  inverse word map is the polynomial monomial with exactly those slot
  exponents; for two slots it is the existing concatenation product.
  Compose with the original first-page homology equivalence.
  \leanok
\end{proof}

\begin{proposition}[The internal double class has the specified square cochain]
  \label{prop:h6-internal-square-polynomial-provenance}
  \lean{KIP126.Classical.Adams.sphereH6DoubleInternalE2_polynomial_representative}
  \uses{prop:h6-exact-polynomial-representatives,prop:h6-double-internal-representative}
  \leanok
  Under the explicit low-level cycle-compatibility hypotheses, the internal
  double class with normalized sphere coefficient has an actual $Z_2$ lift
  whose first-page coordinate under the derived equivalence is precisely
  the existing square cochain. This does not compare the derived coordinates
  with the independently assumed fixed Milnor coordinates or with Lin data.
  Non-boundary, full differential compatibility, and permanence obligations
  remain open.
\end{proposition}
\begin{proof}
  Use the already proved actual $Z_2$ representative and substitute its
  exact polynomial coordinate, retaining its image in the internal page.
  \leanok
\end{proof}

\begin{definition}[Unit-corrected actual cooperation split]
  \label{def:h6-actual-cobar-split}
  \lean{KIP126.StableHomotopy.Cohomology.cooperationCobarDiagonal,KIP126.StableHomotopy.Cohomology.cooperationTensorCobarSplit}
  \uses{prop:h6-tensor-boundary-comodule,def:h6-right-unit-tensor}
  \leanok
  Set $c(a)=1\otimes a+a\otimes1-\Delta(a)$ for actual cooperations.
  On an arbitrary cooperation tensor $w$, define
  $C(w)=U(w)+(1\otimes U)(w)-(\Delta\otimes1)(w)$, with reassociation
  understood. These are operations on the existing tensors; reducedness
  is not assumed for arbitrary inputs.
\end{definition}

\begin{proposition}[Actual splitting and its single-slot polynomial formula]
  \label{prop:h6-actual-cobar-split}
  \lean{KIP126.StableHomotopy.Cohomology.cooperationTensor_unitParts,KIP126.StableHomotopy.Cohomology.cooperationTensorCobarSplit_lof_tmul,KIP126.StableHomotopy.Cohomology.cooperationCobarDiagonal_polynomial}
  \uses{def:h6-actual-cobar-split,prop:h6-actual-primitive-cooperation}
  \leanok
  The operation $C$ applies $c$ to the first tensor factor and reassociates.
  Under the explicit low-level unit and Milnor coproduct comparisons,
  $Q(c(a))=d_{\mathrm{poly}}P(a)$, where the right side is the existing
  length-one polynomial differential, not a new prescribed differential.
\end{proposition}
\begin{proof}
  Expand both unit terms and track the graded associator casts. Apply
  the previously proved coproduct and unit-coordinate formulas. In
  characteristic two subtraction equals addition, giving exactly the
  two coaugmentation terms and the single split-slot term.
  \leanok
\end{proof}

\begin{proposition}[First-differential recurrence for arbitrary tensor boundaries]
  \label{prop:h6-actual-cobar-recurrence}
  \lean{KIP126.Classical.Adams.adamsHomologyD1_tensorBoundary_cobar_recurrence,KIP126.Classical.Adams.sphereAdamsHomologyD1_tensorBoundary_cobar}
  \uses{prop:h6-actual-cobar-split,prop:h6-first-cycle-coaction,prop:h6-tensor-boundary-comodule,prop:h6-sphere-unit-image}
  \leanok
  Write $L_q$ for tensoring the identity with the actual degree-minus-one
  boundary, and $e$ for the derived next-stage tensor coordinates.
  Under explicit K\"unneth compatibilities, for every tensor $w$ one has
  \[
    \iota e D_1 q(w)=L_q C(w)-L_{qU}(w).
  \]
  Here $qU=D_1$ on the remaining factor by the proved unit formula.
  For sphere coefficients this remaining-factor term is zero. No
  primitivity assumption on the first factor is made.
\end{proposition}
\begin{proof}
  Substitute the boundary coaction formula into the actual
  unit-minus-coaction differential formula. Expand $C$, commute the
  unit insertion with the lowering map, and cancel the remaining-factor
  term. For the sphere, its first differential already vanishes.
  \leanok
\end{proof}

\begin{proposition}[Exact test for the incoming first differential]
  \label{prop:h6-first-incoming-cobar-image}
  \lean{KIP126.Classical.Adams.adamsTensorBoundary_surjective,KIP126.Classical.Adams.sphereAdamsHomologyD1_image_cobar_iff,KIP126.Classical.Adams.sphereAdamsPageD_one_128_image_cobar_iff}
  \uses{prop:h6-actual-cobar-recurrence,prop:h6-reduced-boundary-coordinates,prop:h6-first-page-homology}
  \leanok
  Every next-stage homology element is an actual tensor boundary. Thus an
  element $y$ at the second sphere tower stage lies in the image of $D_1$
  if and only if there exists a tensor $w$ with $L_qC(w)=\iota e(y)$.
  In particular this gives an exact test for the image of
  $d_1:E_1^{1,128}\to E_1^{2,128}$ on the original quotient pages.
  The assertion does not prove that the double representative fails the
  test by itself. The polynomial image criterion and non-boundary
  calculation below subsequently exclude it for the normalized double
  representative; later survival remains open.
\end{proposition}
\begin{proof}
  For surjectivity use the reduced inclusion of $e(y)$ and the proved
  inverse identity for the actual boundary. Apply the recurrence to
  such a preimage; conversely use the injective inclusion and equivalence
  to recover equality of homology classes. Transfer the image criterion
  through the existing first-page homology equivalences.
  \leanok
\end{proof}

\begin{proposition}[Tensor equation for vanishing of the internal double class]
  \label{prop:h6-internal-square-cobar-image}
  \lean{KIP126.Classical.Adams.sphereH6DoubleTensorRepresentative_normalized,KIP126.Classical.Adams.sphereH6DoubleInternalE2_eq_zero_iff_is_differential,KIP126.Classical.Adams.sphereH6DoubleInternalE2_eq_zero_iff_cobar}
  \uses{prop:h6-internal-first-cycle-image,prop:h6-first-incoming-cobar-image,def:h6-double-internal-second-page}
  \leanok
  Assume the explicit low-level compatibility conditions used to construct
  the double class, and let $P(a)=[\xi_1^{64}]$ and $x\in H_0(S)$.
  Write $b=q_S(a\otimes x)$. The constructed internal double class in
  $E_2^{2,128}$ is zero if and only if some degree-128 cooperation tensor
  $w$ satisfies
  \[
    L_{q_S}C(w)=a\otimes b.
  \]
  Direct-sum inclusions at cooperation degree 64 are understood.
  Thus proving this equation has no solution would prove nonvanishing
  of the actual internal class. No assertion of insolubility, fixed Lin
  comparison, or higher-page survival is made here.
\end{proposition}
\begin{proof}
  Combine the internal first-cycle image test with the actual tower
  image criterion. Normalize the double tensor boundary: its
  augmentation correction vanishes because the cooperation degree is
  64, leaving exactly $a\otimes b$. This last normalization needs no
  Milnor or coproduct comparison assumptions.
  \leanok
\end{proof}

\begin{definition}[Removal of the sphere coefficient factor]
  \label{def:h6-sphere-tensor-coefficient}
  \lean{KIP126.Classical.Adams.sphereHomologyScalarEquiv,KIP126.Classical.Adams.sphereCooperationTensorEquiv}
  \uses{def:h6-normalized-word-representatives}
  \leanok
  The actual sphere coefficient homology is concentrated in degree zero,
  where the existing empty-word coordinate identifies it with $\mathbb F_2$.
  Thus in every integer degree $n$ there is a constructed linear equivalence
  $(A\otimes H_*S)_n\simeq A_n$: discard the zero summands, apply the scalar
  coordinate on the remaining factor, and use the tensor right unitor.
  No K\"unneth map, Milnor basis, or Adams-page comparison is required.
\end{definition}

\begin{proposition}[Unique cooperation coefficient of every sphere tensor]
  \label{prop:h6-sphere-tensor-coefficient}
  \lean{KIP126.Classical.Adams.sphereHomologyScalarEquiv_symm_one,KIP126.Classical.Adams.sphereCooperationTensorEquiv_symm_apply,KIP126.Classical.Adams.sphereCooperationTensor_existsUnique}
  \uses{def:h6-sphere-tensor-coefficient}
  \leanok
  The inverse sends $b\in A_n$ to the degree-$n$ summand $b\otimes1$,
  with the explicit degree transport $n-n=0$. Its sphere coefficient is
  exactly the previously specified normalized sphere unit. Consequently
  every tensor in $(A\otimes H_*S)_n$ has a unique such coefficient $b$.
\end{proposition}
\begin{proof}
  Expand the inverse direct-sum concentration equivalence, the tensor
  product of equivalences, and the right unitor. Existence and uniqueness
  then follow from bijectivity.
  \leanok
\end{proof}

\begin{proposition}[The incoming search reduces to degree-128 cooperations]
  \label{prop:h6-internal-square-cooperation-source}
  \lean{KIP126.Classical.Adams.sphereH6DoubleInternalE2_eq_zero_iff_cooperation_cobar}
  \uses{prop:h6-sphere-tensor-coefficient,prop:h6-internal-square-cobar-image,prop:h6-actual-cobar-split}
  \leanok
  Under the same explicit low-level hypotheses, the constructed internal
  double class is zero if and only if some $b\in A_{128}$ satisfies
  \[
    L_{q_S}\operatorname{assoc}(c(b)\otimes1)
      =a\otimes q_S(a\otimes x),\qquad
    c(b)=1\otimes b+b\otimes1-\Delta b.
  \]
  All direct-sum inclusions and degree transports are retained in Lean.
  The polynomial formula for $c(b)$ is already proved. The comparison of
  the remaining map $L_{q_S}\operatorname{assoc}(-\otimes1)$ and the
  non-boundary calculation for the normalized sphere coefficient are
  supplied by the later nodes below.
\end{proposition}
\begin{proof}
  Replace the quantified tensor in the exact zero criterion by its unique
  cooperation coefficient and apply the elementary-tensor split formula.
  Conversely the coefficient supplies a candidate tensor by the inverse
  equivalence. No candidates are lost and no differential is postulated.
  \leanok
\end{proof}

\begin{proposition}[Word and first-page coordinates of every reduced boundary]
  \label{prop:h6-boundary-word-coordinates}
  \lean{KIP126.Classical.Adams.sphereTowerHomologyWordEquiv_boundary_reduced,KIP126.Classical.Adams.sphereTowerHomologyWordEquiv_boundary_basis_single,KIP126.Classical.Adams.sphereFirstPageMilnorEquiv_boundary_reduced}
  \uses{prop:h6-reduced-boundary-coordinates,prop:h6-derived-milnor-coordinates}
  \leanok
  For every tower stage and every reduced tensor $w$, the word coordinate
  of the actual boundary $q(\iota w)$ is precisely the already constructed
  tensor-word coordinate of $w$. A reduced basis factor followed by a
  specified word gives their concatenation, with its coefficient unchanged.
  On the original first quotient page, the corresponding formula takes
  values in the existing normalized polynomial cochains.
  No coproduct, unit, or suspension compatibility predicate is required
  beyond the structures used to construct the tensor coordinates themselves.
\end{proposition}
\begin{proof}
  Expand the successor word coordinate, cancel the degree-transport
  equivalence, and apply the reduced-boundary inverse identity.
  For the basis formula use the proved tensor-word basis calculation.
  For the first page additionally cancel its homology-coordinate inverse.
  \leanok
\end{proof}

\begin{definition}[Word realization in the existing polynomial tensor powers]
  \label{def:h6-word-polynomial-realization}
  \lean{KIP126.Steenrod.Milnor.singleMilnorWordEquiv,KIP126.Steenrod.Milnor.milnorWordPolynomial}
  \uses{prop:h6-cobar-monomial-basis}
  \leanok
  Positive monomials of degree $t$ are equivalent to length-one words.
  More generally realize word coefficients by sending a word to the
  polynomial monomial with exactly its slot exponents and extending linearly.
  Integer degrees are allowed; the target remains the existing polynomial
  tensor power, not a replacement cochain model.
\end{definition}

\begin{proposition}[Faithfulness and agreement of polynomial realization]
  \label{prop:h6-word-polynomial-realization}
  \lean{KIP126.Steenrod.Milnor.singleMilnorWord_exponents,KIP126.Steenrod.Milnor.wordConsEquiv_zero,KIP126.Steenrod.Milnor.cochainsWordEquiv_symm_single_val,KIP126.Steenrod.Milnor.milnorWordPolynomial_single,KIP126.Steenrod.Milnor.milnorWordPolynomial_injective,KIP126.Steenrod.Milnor.milnorWordPolynomial_eq_cochainsWordEquiv_symm}
  \uses{def:h6-word-polynomial-realization,prop:h6-cobar-monomial-basis}
  \leanok
  Realization is injective at each fixed integer degree. In nonnegative
  degree it equals the inverse of the existing cochain word equivalence
  followed by inclusion of the normalized cochains into polynomials.
  The single-word realization places the monomial in slot zero.
\end{proposition}
\begin{proof}
  Distinct words have distinct slot exponent vectors, hence linearly
  independent polynomial monomials. Check agreement with the existing
  cochain map on single word coefficients and use linearity.
  \leanok
\end{proof}

\begin{proposition}[The full first boundary on reduced cooperation coefficients]
  \label{prop:h6-first-boundary-word-coordinate}
  \lean{KIP126.Classical.Adams.sphereTowerHomologyWordEquiv_firstBoundary_basis,KIP126.Classical.Adams.sphereTowerHomologyWordEquiv_firstBoundary_reduced}
  \uses{prop:h6-boundary-word-coordinates,prop:h6-sphere-tensor-coefficient,def:h6-word-polynomial-realization}
  \leanok
  For any reduced cooperation $a$ of degree $n+1$, insert the normalized
  sphere coefficient and apply the actual boundary to $H_n(T_1 S)$.
  Its word coefficients are exactly the reduced Milnor basis coefficients
  of $a$, relabeled as length-one words. This includes every reduced
  element, not only $\xi_1^{64}$. The existing $h_6$ calculation is now
  a specialization of this theorem.
\end{proposition}
\begin{proof}
  The sphere scalar coordinate is the empty word with coefficient one.
  The preceding reduced-basis formula prepends the chosen monomial.
  Extend the equality of the two linear maps from the reduced basis.
  \leanok
\end{proof}

\begin{proposition}[Polynomial value of the actual first boundary]
  \label{prop:h6-first-boundary-polynomial}
  \lean{KIP126.Classical.Adams.sphereFirstBoundary_polynomial_reduced}
  \uses{prop:h6-first-boundary-word-coordinate,prop:h6-word-polynomial-realization,prop:h6-full-milnor-coordinates}
  \leanok
  The polynomial coordinate of $q_S(a\otimes1)$ is exactly $P(a)$ for
  every reduced cooperation $a$, with the degree shifts explicit.
  This computes the actual boundary; no differential comparison is
  assumed. Double reduction of the corrected coproduct is proved below;
  the remaining application requires the graded tensor compatibility calculations.
  The full incoming differential comparison and non-boundary assertion
  are not claimed here.
\end{proposition}
\begin{proof}
  Apply word realization to the first-boundary coordinate formula.
  On each reduced basis vector its single-word monomial equals the
  already defined cooperation polynomial. Extend by linearity.
  \leanok
\end{proof}

\begin{definition}[Actual double-reduced cooperation tensors]
  \label{def:h6-double-reduced-cooperations}
  \lean{KIP126.StableHomotopy.Cohomology.reducedCooperationSquare,KIP126.StableHomotopy.Cohomology.reducedCooperationSquareInclusion}
  \uses{prop:h6-graded-tensor-kernel}
  \leanok
  Use the existing reduced tensor construction with $\ker\epsilon$ also
  as the remaining factor. Its inclusion into $A\otimes A$ is the
  composite of the two actual factor inclusions. This introduces no
  new quotient, cochain model, or spectral-sequence page.
\end{definition}

\begin{proposition}[Double inclusion and reduced-basis support]
  \label{prop:h6-double-reduced-support}
  \lean{KIP126.StableHomotopy.Cohomology.reducedCooperationSquareInclusion_lof_tmul,KIP126.StableHomotopy.Cohomology.reducedCooperationSquareInclusion_injective,KIP126.StableHomotopy.Cohomology.cooperationMilnorBasis_reduced,KIP126.StableHomotopy.Cohomology.cooperationTensorMilnorPolynomial_coeff,KIP126.StableHomotopy.Cohomology.cooperationTensorMilnorBasis_support_reduced,KIP126.StableHomotopy.Cohomology.cooperationTensor_mem_span_reduced_of_normalized,KIP126.StableHomotopy.Cohomology.cooperationTensorMilnorBasis_positive}
  \uses{def:h6-double-reduced-cooperations,prop:h6-milnor-polynomial-faithful}
  \leanok
  The double inclusion is injective and preserves elementary represented
  factors. Coefficients of the actual tensor basis are precisely the
  polynomial coefficients at its distinct pair exponents. If the
  polynomial of a tensor is killed by both slot augmentations, every
  supported basis pair has two nonconstant factors. Each such basis
  vector is exactly an included tensor of the original reduced bases.
\end{proposition}
\begin{proof}
  Tensor the injective factor maps over the field and take the direct sum.
  Compare coefficient functionals on the full tensor basis. Normalization
  forces every monomial with nonzero coefficient to use both slots;
  a constant factor would contradict this. Identify the positive full
  basis vectors with the previously specified reduced basis vectors.
  \leanok
\end{proof}

\begin{proposition}[Polynomial normalization yields an actual reduced tensor]
  \label{prop:h6-double-reduced-lift}
  \lean{KIP126.StableHomotopy.Cohomology.cooperationTensor_reduced_span_le_range,KIP126.StableHomotopy.Cohomology.cooperationTensor_existsUnique_reduced_of_normalized,KIP126.StableHomotopy.Cohomology.cooperationMilnorPolynomial_reduced_normalized}
  \uses{prop:h6-double-reduced-support}
  \leanok
  Every polynomially normalized tensor has a unique preimage under the
  actual double-reduced inclusion. Also every reduced cooperation has
  a normalized single-slot polynomial. Neither assertion assumes an
  Adams differential or a coproduct formula.
\end{proposition}
\begin{proof}
  Every positive pair basis vector has an explicit reduced tensor lift,
  so its span lies in the inclusion's range. Apply the support result,
  and use injectivity for uniqueness. Check single-slot normalization
  on the reduced basis, whose exponent vectors are nonzero, and extend
  linearly.
  \leanok
\end{proof}

\begin{proposition}[Double reduction of the actual corrected coproduct]
  \label{prop:h6-corrected-coproduct-double-reduced}
  \lean{KIP126.StableHomotopy.Cohomology.cooperationCobarDiagonal_reduced_normalized,KIP126.StableHomotopy.Cohomology.cooperationCobarDiagonal_mem_span_reduced,KIP126.StableHomotopy.Cohomology.cooperationCobarDiagonal_mem_span_reduced_of_ne,KIP126.StableHomotopy.Cohomology.cooperationCobarDiagonal_existsUnique_reduced,KIP126.StableHomotopy.Cohomology.cooperationCobarDiagonal_existsUnique_reduced_of_ne}
  \uses{prop:h6-double-reduced-lift,prop:h6-actual-cobar-split,prop:h6-polynomial-cocycles}
  \leanok
  Under the explicit unit and Milnor coproduct compatibilities, the
  corrected coproduct $c(a)=1\otimes a+a\otimes1-\Delta a$ of a reduced
  cooperation has a unique actual double-reduced tensor lift. All
  nonzero-degree cooperations, including every degree-128 candidate,
  satisfy the reduced-input condition automatically.
\end{proposition}
\begin{proof}
  The proved polynomial formula is $Q(c(a))=dP(a)$. The single-slot input
  is normalized, and the polynomial differential preserves normalization
  in both slots. Apply the actual reduced-lift theorem. In nonzero degree
  the coefficient counit vanishes by the Eilenberg--Mac Lane property.
  \leanok
\end{proof}

\begin{definition}[Derived reduced cooperation differential]
  \label{def:h6-reduced-coproduct-map}
  \lean{KIP126.StableHomotopy.Cohomology.cooperationReducedCobarDiagonal}
  \uses{prop:h6-corrected-coproduct-double-reduced}
  \leanok
  Restrict the codomain of the actual corrected coproduct to the
  double-reduced tensor via the injective inclusion and its proved
  range property. The result is a linear map, constructed rather than
  supplied as additional differential data.
\end{definition}

\begin{proposition}[The reduced map retains the actual and polynomial values]
  \label{prop:h6-reduced-coproduct-map}
  \lean{KIP126.StableHomotopy.Cohomology.cooperationReducedCobarDiagonal_inclusion,KIP126.StableHomotopy.Cohomology.cooperationReducedCobarDiagonal_polynomial}
  \uses{def:h6-reduced-coproduct-map,prop:h6-actual-cobar-split}
  \leanok
  Inclusion of the constructed reduced map gives exactly $c(a)$, and
  its polynomial is exactly the existing one-slot cobar differential
  of $P(a)$. Reducedness of the output is proved, not assumed.
\end{proposition}
\begin{proof}
  Use the defining property of codomain restriction along an injection,
  then the previously proved corrected-coproduct polynomial formula.
  \leanok
\end{proof}

\begin{proposition}[Actual first differential through the reduced coproduct]
  \label{prop:h6-actual-first-differential-reduced-coproduct}
  \lean{KIP126.Classical.Adams.sphereAdamsHomologyD1_firstBoundary_reduced}
  \uses{prop:h6-reduced-coproduct-map,prop:h6-actual-cobar-recurrence,prop:h6-sphere-tensor-coefficient}
  \leanok
  For every reduced cooperation $a$ in degree $n$, and under the stated
  low-level K\"unneth compatibilities, the actual first differential
  of $q_S(a\otimes1)$ has included next-stage tensor coordinate
  \[
    L_{q_S}\operatorname{assoc}(\iota\overline c(a)\otimes1).
  \]
  Here $\overline c$ is the constructed double-reduced map. This retains
  the original tower representative. Composing the whole graded tensor
  expression with faithful two-slot coordinates and the non-boundary
  calculation are completed below; higher-page survival remains open.
\end{proposition}
\begin{proof}
  Substitute the coefficient-removal inverse and elementary tensor
  splitting formula into the actual sphere recurrence, then use the
  inclusion formula for the reduced coproduct.
  \leanok
\end{proof}

\begin{definition}[The reduced sphere boundary equivalence]
  \label{def:h6-reduced-sphere-boundary-equiv}
  \lean{KIP126.Classical.Adams.sphereReducedCooperationTensorEquiv,KIP126.Classical.Adams.sphereReducedBoundaryEquiv}
  \uses{prop:h6-sphere-tensor-coefficient,prop:h6-reduced-boundary-coordinates}
  \leanok
  Remove the normalized sphere coefficient from
  $(\overline A\otimes H_*S)_n$ by degree-zero concentration.
  Compose its inverse with the actual next-stage tensor-coordinate inverse
  to construct $e_1:\overline A_n\simeq H_{n-1}(T_1S)$.
  No Milnor basis or differential comparison is supplied.
\end{definition}

\begin{proposition}[The reduced equivalence is the actual sphere boundary]
  \label{prop:h6-reduced-sphere-boundary-equiv}
  \lean{KIP126.Classical.Adams.sphereHomologyScalarEquiv_cast_symm_one,KIP126.Classical.Adams.sphereReducedCooperationTensorEquiv_symm_apply,KIP126.Classical.Adams.sphereReducedCooperationTensorEquiv_inclusion,KIP126.Classical.Adams.sphereReducedBoundaryEquiv_apply}
  \uses{def:h6-reduced-sphere-boundary-equiv}
  \leanok
  The coefficient-removal inverse inserts $a\otimes1$, commutes with
  inclusion into ordinary tensors, and the normalized coefficient is
  preserved by degree transport. Consequently $e_1(a)=q_S(a\otimes1)$.
\end{proposition}
\begin{proof}
  Expand the concentrated direct-sum equivalence. Apply the next-stage
  tensor coordinate and its proved inverse identity on reduced tensors.
  \leanok
\end{proof}

\begin{definition}[Tensoring a degree-lowering equivalence]
  \label{def:h6-tensor-lowering-equiv}
  \lean{KIP126.Core.Algebra.gradedTensorLowerEquiv}
  \uses{prop:h6-graded-tensor-kernel}
  \leanok
  Degreewise linear equivalences $V_i\simeq W_{i-1}$ induce
  $(A\otimes V)_n\simeq(A\otimes W)_{n-1}$ by tensoring with the
  identity, transporting the displayed degree, and taking the direct sum.
\end{definition}

\begin{proposition}[Agreement with the existing lowering map]
  \label{prop:h6-tensor-lowering-equiv}
  \lean{KIP126.Core.Algebra.gradedTensorLowerEquiv_toLinearMap,KIP126.Core.Algebra.gradedTensorLowerMap_injective}
  \uses{def:h6-tensor-lowering-equiv}
  \leanok
  The underlying linear map is the existing graded tensor lowering map.
  More generally, degreewise injections give an injective lowering map
  whenever each left factor is flat.
\end{proposition}
\begin{proof}
  Check the map formula on elementary tensors. For injectivity use flatness
  of each left factor and injectivity of the direct-sum map.
  \leanok
\end{proof}

\begin{definition}[Double-reduced actual boundary coordinates]
  \label{def:h6-double-reduced-boundary-equiv}
  \lean{KIP126.Classical.Adams.sphereCooperationSquareBoundary,KIP126.Classical.Adams.sphereDoubleReducedBoundaryEquiv,KIP126.Classical.Adams.sphereSecondReducedBoundaryEquiv}
  \uses{prop:h6-reduced-sphere-boundary-equiv,prop:h6-tensor-lowering-equiv,def:h6-double-reduced-cooperations}
  \leanok
  Let $F_n(z)=L_{q_S}\operatorname{assoc}(z\otimes1)$ on ordinary
  cooperation tensors. Tensoring $e_1$ with a reduced left factor gives
  $\overline A\otimes\overline A\simeq\overline A\otimes H_*T_1S$,
  with the required degree shift. Composing with the next actual boundary
  equivalence gives $e_2:(\overline A\otimes\overline A)_n
  \simeq H_{n-2}(T_2S)$.
\end{definition}

\begin{proposition}[Faithful double-reduced boundary transport]
  \label{prop:h6-double-reduced-boundary-equiv}
  \lean{KIP126.Classical.Adams.sphereCooperationSquareBoundary_lof_tmul,KIP126.Classical.Adams.sphereDoubleReducedBoundaryEquiv_lof_tmul,KIP126.Classical.Adams.sphereDoubleReducedBoundary_injective,KIP126.Classical.Adams.sphereDoubleReducedBoundaryEquiv_inclusion}
  \uses{def:h6-double-reduced-boundary-equiv}
  \leanok
  On elementary tensors $F_n$ applies $q_S(b\otimes1)$ to the right
  factor, retaining the left factor and the explicit degree transport.
  On every double-reduced tensor it agrees with the constructed equivalence
  followed by inclusion. This composite is injective: the actual boundary
  step loses no information on the double-reduced subspace.
\end{proposition}
\begin{proof}
  Expand the associator and the lowering map on elementary tensors and
  use invariance of the normalized sphere coefficient under transport.
  Extend linearly; injectivity follows from the equivalence and inclusion.
  \leanok
\end{proof}

\begin{proposition}[First incoming differential and polynomial image criterion]
  \label{prop:h6-first-incoming-polynomial-comparison}
  \lean{KIP126.Classical.Adams.sphereAdamsHomologyD1_reduced_comparison,KIP126.Classical.Adams.sphereAdamsHomologyD1_reduced_polynomial,KIP126.Classical.Adams.sphereAdamsHomologyD1_mem_range_iff_polynomial}
  \uses{prop:h6-double-reduced-boundary-equiv,prop:h6-actual-first-differential-reduced-coproduct,prop:h6-reduced-coproduct-map,prop:h6-milnor-polynomial-faithful}
  \leanok
  Under the explicit low-level K\"unneth and Milnor coproduct
  compatibilities, $d_1e_1=e_2\overline c$ in every integer degree.
  Thus $Q\iota e_2^{-1}d_1e_1(a)=dP(a)$, using the existing polynomial
  differential and cooperation coordinate maps $P,Q$.
  For any actual $y\in H_{n-2}(T_2S)$,
  \[
    y\in\operatorname{im}(d_1:H_{n-1}(T_1S)\to H_{n-2}(T_2S))
    \quad\Longleftrightarrow\quad
    \exists a\in\overline A_n,\ dP(a)=Q\iota e_2^{-1}(y).
  \]
  These are derived conditional theorems, not new comparison axioms.
  They concern the first incoming differential only. The target square
  and source cochains are identified below, followed by a proved
  non-boundary certificate. Comparing the fixed Lin class and higher-page permanence remain separate.
\end{proposition}
\begin{proof}
  Cancel the injective inclusion from the actual first-differential
  formula and the double-reduced boundary transport identity.
  Apply the reduced coproduct polynomial formula. For the image criterion
  use surjectivity of $e_1$ in one direction and faithfulness of $Q\iota$
  together with $e_2$ in the other.
  \leanok
\end{proof}

\begin{definition}[Reduced cooperations as all one-slot cochains]
  \label{def:h6-reduced-cooperation-cochains}
  \lean{KIP126.StableHomotopy.Cohomology.reducedCooperationCochainEquiv}
  \uses{def:h6-reduced-milnor-basis-input,def:h6-word-polynomial-realization,prop:h6-cobar-monomial-basis}
  \leanok
  For every natural internal degree $n$, use the reduced Milnor basis,
  the equivalence of positive monomials with single-slot words, and the
  existing word coordinates to construct
  $\overline A_n\simeq C^{1,n}_{\mathrm{cobar}}$.
  No tower, K\"unneth comparison, or coproduct condition is required.
\end{definition}

\begin{proposition}[The source quantifier exhausts normalized cobar cochains]
  \label{prop:h6-reduced-cooperation-cochains}
  \lean{KIP126.StableHomotopy.Cohomology.reducedCooperationCochainEquiv_val,KIP126.StableHomotopy.Cohomology.reducedCooperation_polynomial_boundary_iff}
  \uses{def:h6-reduced-cooperation-cochains,prop:h6-full-milnor-coordinates}
  \leanok
  The underlying polynomial of the cochain associated to $a$ is exactly
  $P(a)$. For any $y\in C^{2,n}_{\mathrm{cobar}}$, the condition
  $\exists a\in\overline A_n,\ dP(a)=y$ is equivalent to
  $\exists b\in C^{1,n}_{\mathrm{cobar}},\ db=y$.
  In particular, the actual-cooperation quantifier does not omit any
  possible normalized cochain source.
\end{proposition}
\begin{proof}
  Check equality with $P$ on the reduced Milnor basis and extend linearly.
  Use the equivalence and its surjectivity, then equality of cochains is
  equality of their underlying polynomials.
  \leanok
\end{proof}

\begin{proposition}[The specified square in double-reduced boundary coordinates]
  \label{prop:h6-square-double-reduced-coordinate}
  \lean{KIP126.Classical.Adams.sphereSecondReducedBoundaryEquiv_h6_double,KIP126.Classical.Adams.sphereSecondReducedBoundaryEquiv_h6_double_polynomial}
  \uses{prop:h6-double-reduced-boundary-equiv,prop:h6-internal-square-cobar-image,prop:h6-milnor-polynomial-faithful}
  \leanok
  The derived equivalence $e_2$ sends the reduced tensor $a\otimes a$
  in degrees $(64,64)$ to the existing twice-applied actual boundary
  representative with normalized sphere coefficient one. If
  $P(a)=[\xi_1^{64}]$, then its polynomial coordinate
  $Q\iota e_2^{-1}$ is exactly the existing cochain
  $[\xi_1^{64}\mid\xi_1^{64}]$.
  No differential comparison is needed for this coordinate calculation.
\end{proposition}
\begin{proof}
  Apply the next-stage tensor coordinate and its injective inclusion.
  Both sides are the same elementary boundary tensor. The sphere scalar
  is the already normalized unit. Apply the tensor polynomial formula
  and the stated coordinate of $a$.
  \leanok
\end{proof}

\begin{proposition}[Exact pure-cobar criterion for the internal square]
  \label{prop:h6-internal-square-pure-cobar-criterion}
  \lean{KIP126.Classical.Adams.sphereH6DoubleInternalE2_eq_zero_iff_homologyD1,KIP126.Classical.Adams.sphereH6DoubleInternalE2_eq_zero_iff_polynomial,KIP126.Classical.Adams.sphereH6DoubleInternalE2_eq_zero_iff_cobar_boundary,KIP126.Classical.Adams.sphereH6DoubleInternalE2_ne_zero_iff_cobar_nonboundary}
  \uses{prop:h6-internal-square-cobar-image,prop:h6-first-incoming-polynomial-comparison,prop:h6-reduced-cooperation-cochains,prop:h6-square-double-reduced-coordinate}
  \leanok
  Under the explicit low-level compatibility conditions, take the
  constructed internal double class $x$ with $P(a)=[\xi_1^{64}]$ and
  normalized sphere coefficient one. Then
  \[
    x=0\quad\Longleftrightarrow\quad
    \exists b\in C^{1,128}_{\mathrm{cobar}},\ db=h_6^2,
  \]
  and therefore $x\ne0$ is equivalent to the assertion that no such
  cochain exists. The right side uses only the already defined normalized
  Milnor cobar complex, with no categorical or tower source objects.
  This equivalence does not itself prove the non-boundary calculation,
  supply the fixed low-level witnesses, identify the fixed Lin class,
  or assert higher-page permanence.
\end{proposition}
\begin{proof}
  Transport the original internal first-cycle zero criterion through
  the actual first-page homology equivalences. Apply the first incoming
  polynomial image criterion, the target-square calculation, and the
  exhaustive cochain-source equivalence. Negate to obtain the nonzero test.
  \leanok
\end{proof}

\begin{definition}[Seven-coefficient detection of the standard square]
  \label{def:h6-cobar-square-detector}
  \lean{KIP126.Steenrod.Milnor.xiOneProjection,KIP126.Steenrod.Milnor.h6SquarePureDetector,KIP126.Steenrod.Milnor.h6SquareDetector}
  \uses{def:h6-milnor-cobar}
  \leanok
  Project each tensor slot onto its pure $\xi_1$ powers, killing the
  higher Milnor generators. On the resulting polynomial ring
  $\mathbb F_2[X,Y]$ set
  \[
    D(f)=\sum_{i=0}^{6}[X^{128-2^i}Y^{2^i}]f.
  \]
  Compose this linear map with the projection to obtain the detector on
  the original two-slot cobar polynomials. No external table is used.
\end{definition}

\begin{proposition}[Detector normalization and removal of unit terms]
  \label{prop:h6-cobar-square-detector-basic}
  \lean{KIP126.Steenrod.Milnor.xiOneProjection_X,KIP126.Steenrod.Milnor.xiOneProjection_xi,KIP126.Steenrod.Milnor.h6SquarePureDetector_apply,KIP126.Steenrod.Milnor.h6SquareDetector_square,KIP126.Steenrod.Milnor.h6SquarePureDetector_rename,KIP126.Steenrod.Milnor.h6SquareDetector_insertLeft,KIP126.Steenrod.Milnor.h6SquareDetector_insertRight,KIP126.Steenrod.Milnor.h6SquareDetector_differential}
  \uses{def:h6-cobar-square-detector,prop:h6-cobar-projection-generators}
  \leanok
  The detector takes value one on
  $[\xi_1^{64}\mid\xi_1^{64}]$ and vanishes on a polynomial supported
  in only one slot. Both coaugmentation terms are therefore invisible.
  On a one-slot polynomial its cobar differential has the same detector
  value as its projected coproduct.
\end{proposition}
\begin{proof}
  Evaluate the seven coefficients of the square directly. Each tested
  monomial uses both slots with positive exponent, so none can occur in
  a one-slot renaming. Use the projection formulas for the two insertions
  and the definition of the cobar differential.
  \leanok
\end{proof}

\begin{proposition}[Projection reduces all source monomials to two generators]
  \label{prop:h6-cobar-projection-generators}
  \lean{KIP126.Steenrod.Milnor.xiOneProjection_coproductGenerator,KIP126.Steenrod.Milnor.xiOneProjection_split_X,KIP126.Steenrod.Milnor.xiOneProjection_insertLeft,KIP126.Steenrod.Milnor.xiOneProjection_insertRight,KIP126.Steenrod.Milnor.xiOneProjection_split_monomial_high,KIP126.Steenrod.Milnor.singleSlot_exponents_eq_two,KIP126.Steenrod.Milnor.xiOneProjection_split_monomial_two,KIP126.Steenrod.Milnor.singleSlot_two_weight}
  \uses{def:h6-cobar-square-detector}
  \leanok
  The projected coproduct sends $\xi_1$ to $X+Y$, $\xi_2$ to $X^2Y$,
  and every higher generator to zero. It sends $\xi_1^a\xi_2^b$ to
  $(X+Y)^a(X^2Y)^b$, of degree $a+3b$; every monomial containing a
  higher generator has zero image. A one-slot exponent vector with no
  higher generator is exactly the sum of its two indicated components.
  Projection also commutes with inserting an empty slot at either end.
\end{proposition}
\begin{proof}
  Expand Milnor's generator coproduct. For a generator of index at least
  three every summand has a factor killed by the projection. Extend by
  multiplicativity. The exponent-vector identity is pointwise, and the
  weight formula follows by additivity of weight.
  \leanok
\end{proof}

\begin{definition}[Bounded binary binomial parity]
  \label{def:h6-binary-binomial-parity}
  \lean{KIP126.Steenrod.Milnor.binaryChooseParity}
  \leanok
  Define the parity function by the product of the eight binary-digit
  binomial coefficients, reduced modulo two. Eight digits suffice for
  the integers occurring here, all less than $256$.
\end{definition}

\begin{proposition}[Kernel-checked binomial cancellation]
  \label{prop:h6-binomial-parity-certificate}
  \lean{KIP126.Steenrod.Milnor.binaryChooseParity_eq,KIP126.Steenrod.Milnor.h6Square_binary_parity,KIP126.Steenrod.Milnor.h6Square_binomial_parity}
  \uses{def:h6-binary-binomial-parity}
  \leanok
  The binary function equals the actual binomial coefficient modulo two,
  by Lucas's theorem. For every $0\le b\le42$,
  \[
    \sum_{\substack{0\le i\le6\\b\le2^i}}
      {128-3b\choose 2^i-b}=0\quad\text{in }\mathbb F_2.
  \]
  The finite binary identity is checked by the Lean kernel, without a
  native-evaluation axiom or imported computation certificate.
\end{proposition}
\begin{proof}
  Apply Mathlib's proved Lucas theorem, verify the 43 finite binary
  cases using kernel reduction, and transport the resulting parity to
  $\mathbb F_2$.
  \leanok
\end{proof}

\begin{proposition}[The projected monomial contributions are annihilated]
  \label{prop:h6-cobar-detector-binomial}
  \lean{KIP126.Steenrod.Milnor.h6Square_binomial_coefficient,KIP126.Steenrod.Milnor.h6SquarePureDetector_binomial}
  \uses{prop:h6-binomial-parity-certificate,prop:h6-cobar-projection-generators}
  \leanok
  If $a+3b=128$, the coefficient of $X^{128-j}Y^j$ in
  $(X+Y)^a(X^2Y)^b$ is ${a\choose j-b}$ when $b\le j$, and zero otherwise.
  The statement includes all cases where a shifted exponent is too large.
  Hence the detector vanishes on every such projected monomial.
\end{proposition}
\begin{proof}
  Shift coefficients by the monomial $X^{2b}Y^b$ and apply the binomial
  coefficient formula for $(X+Y)^a$. Use binomial symmetry and the
  zero convention outside the range. Sum over $j=2^i$ and apply the
  finite parity certificate, noting $b\le42$.
  \leanok
\end{proof}

\begin{theorem}[The standard square is not a normalized cobar boundary]
  \label{thm:h6-square-cobar-nonboundary}
  \lean{KIP126.Steenrod.Milnor.h6SquareDetector_differential_monomial,KIP126.Steenrod.Milnor.h6SquareDetector_differential_zero,KIP126.Steenrod.Milnor.h6SquareCochain_not_boundary}
  \uses{prop:h6-cobar-square-detector-basic,prop:h6-cobar-detector-binomial,prop:h6-cobar-monomial-basis}
  \leanok
  For every $b\in C^{1,128}_{\mathrm{cobar}}$, one has
  $D(db)=0$ and consequently
  \[
    db\ne[\xi_1^{64}\mid\xi_1^{64}].
  \]
  This is an unconditional theorem of the specified polynomial cobar
  model. Its proof does not import Adams, Lin, or project axioms.
\end{theorem}
\begin{proof}
  On a monomial containing a higher generator the projected coproduct
  is zero. Otherwise its exponents have $a+3b=128$, and the binomial
  cancellation applies. Extend from the actual normalized monomial basis
  to every cochain. Equality with the square would contradict detector
  values zero and one.
  \leanok
\end{proof}

\begin{proposition}[Nonvanishing of the actual internal double class]
  \label{prop:h6-double-internal-nonzero}
  \lean{KIP126.Classical.Adams.sphereH6DoubleInternalE2_ne_zero,KIP126.Classical.Adams.sphereAdamsInternalE2_h6_double_nonzero_exists}
  \uses{prop:h6-internal-square-pure-cobar-criterion,thm:h6-square-cobar-nonboundary}
  \leanok
  Under the explicit low-level structures and coherence conditions, the
  constructed actual internal double class with normalized sphere
  coefficient is nonzero in $E_2^{2,128}$. In particular this page has
  a nonzero element. This does not identify the fixed Lin class, supply
  the fixed low-level witnesses, or prove higher-page survival.
\end{proposition}
\begin{proof}
  Apply the already proved internal nonvanishing equivalence to the
  pure cobar non-boundary theorem. For existence use the unique actual
  cooperation with polynomial coordinate $[\xi_1^{64}]$.
  \leanok
\end{proof}

\begin{proposition}[Polynomial multiplication and the two closed representatives]
  \label{prop:h6-polynomial-cocycles}
  \lean{KIP126.Steenrod.Milnor.cupPolynomial_mem}
  \lean{KIP126.Steenrod.Milnor.h6Polynomial_mem}
  \lean{KIP126.Steenrod.Milnor.h6Polynomial_differential}
  \lean{KIP126.Steenrod.Milnor.h6SquarePolynomial_differential}
  \uses{def:h6-milnor-cobar}
  \leanok
  Concatenation preserves normalization and adds internal degrees.
  The polynomial $[\xi_1^{64}]$ is normalized of internal degree $64$.
  Its polynomial cobar differential and that of
  $[\xi_1^{64}\mid\xi_1^{64}]$ both vanish.
\end{proposition}
\begin{proof}
  Disjoint slot renamings preserve the weight. Augmentation of any slot
  annihilates the corresponding factor of a concatenation. Since
  $\xi_1$ is primitive, the characteristic-two identity
  $(a+b)^{64}=a^{64}+b^{64}$ gives the two differential calculations.
  Preservation of all normalized homogeneous cochains is also proved;
  it does not by itself identify the topological differential with this
  polynomial differential.
  \leanok
\end{proof}

% SOURCE: explicit abstract-foundation boundary confirmed for this task.
% This node is a supplied mathematical structure, not a proved comparison theorem.
\begin{definition}[Milnor coordinates of the mod--2 foundation]
  \label{def:h6-milnor-cooperations}
  \lean{KIP126.Classical.Adams.MilnorCooperations}
  \lean{KIP126.Classical.Adams.sphereFirstDifferential}
  \uses{def:h6-tower-differential,def:h6-milnor-cobar}
  \notready
  The explicit Milnor cooperation input consists of additive isomorphisms
  $\phi^{s,t}:E_1^{s,t}(S^0)\cong C^{s,t}$ for $s,t\ge0$, satisfying
  $\phi^{s+1,t}d_1=d_{\mathrm{cobar}}\phi^{s,t}$.
  Here $E_1$ and $d_1$ are those already constructed from the tower.
  This is a foundational input; its existence is not inferred merely from
  the homotopy groups of $H$.  It contains no later page, named $h$ class,
  selected multiplication, or permanence assertion.
  Proposition~\ref{prop:h6-derived-milnor-coordinates} now constructs its
  vector-space coordinates from below-page inputs. The coproduct comparison
  and the fixed instances needed to replace this whole input remain open.
\end{definition}

\begin{proposition}[The actual initial sphere column survives in internal SSData]
  \label{prop:h6-sphere-initial-survival}
  \lean{KIP126.Classical.Adams.sphereAdamsResolutionBoundary_zero,KIP126.Classical.Adams.sphereAdamsHomologyD1_zero,KIP126.Classical.Adams.sphereAdamsPageD_one_filtration_zero,KIP126.Classical.Adams.sphereAdamsK_zero,KIP126.Classical.Adams.sphereAdamsCycles_filtration_zero,KIP126.Classical.Adams.sphereAdamsDifferential_filtration_zero,KIP126.Classical.Adams.adamsBoundaries_filtration_zero,KIP126.Classical.Adams.sphereAdamsSSData_Z_filtration_zero,KIP126.Classical.Adams.sphereAdamsSSData_B_filtration_zero,KIP126.Classical.Adams.sphereAdams_nonzeroSurvival_filtration_zero_iff}
  \uses{prop:h6-sphere-unit-image,def:h6-tower-internal-data,prop:h6-tower-common-representative,prop:h6-first-differential-homology}
  \leanok
  For the actual sphere tower, the initial connecting map is zero.
  At filtration zero, all finite cycle modules equal the whole first
  page, all boundary modules vanish, and every outgoing differential
  is zero. For its internal SSData this gives $Z_n=\top$ and $B_n=\bot$
  for every $n\in\mathbb N\cup\{\infty\}$. Consequently, for every
  integer $t$ and internal class $x\in E_2^{0,t}$,
  \[
    \operatorname{NonzeroSurvival}(x)\quad\Longleftrightarrow\quad x\ne0.
  \]
  Only the specified Eilenberg--Mac Lane object, stable category, and
  actual cofiber construction are used: no ring, K\"unneth, Milnor,
  Lin-table, or fixed-foundation witness is required. This proves the
  initial column, not the filtration-two square's permanence.
\end{proposition}
\begin{proof}
  Unit surjectivity and exactness kill the initial resolution boundary,
  hence the exact-couple map $k$. Zero lifts through every tower map, so
  every class lies in every cycle module and has zero differential.
  The tower maps from filtration zero to its constant negative extension
  are isomorphisms, so the defining boundary kernels vanish. Restrict to
  the actual $Z_2$ ambient module and take the intersection/union stages.
  In the proved common-representative characterization of survival,
  every representative meets all cycle conditions and the boundary
  condition says precisely that its quotient class is nonzero.
  \leanok
\end{proof}

% SOURCE: consequences of the constant negative-stage extension and the
% the Eilenberg--Mac Lane property; neither Milnor nor Lin coordinates are used.
\begin{proposition}[No incoming differential at the square's bidegree]
  \label{prop:h6-no-incoming}
  \lean{KIP126.Classical.Adams.sphereAdamsPageOne_filtration_zero_subsingleton,KIP126.Classical.Adams.adamsLayerAt_isZero_of_negative,KIP126.Classical.Adams.sphereAdamsInternal_h6_incoming_source_subsingleton,KIP126.Classical.Adams.sphereAdamsInternal_h6_incoming_d_eq_zero}
  \uses{def:h6-em-object,def:h6-fiber-tower,prop:h6-first-page-homology,def:h6-tower-internal-sequence}
  \leanok
  For the chosen Eilenberg--Mac Lane object and every $r\ge2$, the internal sphere tower satisfies
  \[
    E_r^{2-r,129-r}=0.
  \]
  Consequently every incoming differential to $(2,128)$ is zero.
  No Milnor coordinates or Lin table are required. This is not an assertion
  that outgoing differentials vanish.
\end{proposition}
\begin{proof}
  For $r>2$, the source has negative filtration. The tower is constant there,
  so its transition is an isomorphism and its layer is zero. For $r=2$,
  the source is $(0,127)$. The first-page comparison and the right unitor give
  $E_1^{0,127}\cong\pi_{127}(H\wedge S)\cong\pi_{127}H=0$ by the
  Eilenberg--Mac Lane property. Vanishing propagates from the
  first page to later quotient pages. Transport these facts along the proved
  internal page isomorphisms.
  \leanok
\end{proof}

\begin{proposition}[Survival reduced to initial nonvanishing and liftability]
  \label{prop:h6-survival-liftability}
  \lean{KIP126.Classical.Adams.adamsBoundaries_succ_eq_of_source_subsingleton,KIP126.Classical.Adams.sphereAdams_h6_boundaries_eq_two,KIP126.Classical.Adams.sphereAdams_h6_nonzeroSurvival_iff}
  \uses{prop:h6-no-incoming,def:h6-tower-pages,prop:h6-tower-common-representative}
  \leanok
  With the same Eilenberg--Mac Lane foundation, $B_{n+2}^{2,128}=B_2^{2,128}$ for every
  $n\ge0$. Thus an internal class $x\in E_2^{2,128}$ has nonzero survival
  if and only if $x\ne0$ and there is a representative $z\in Z_2^{2,128}$
  of its image in the actual tower quotient such that
  $z\in Z_{n+2}^{2,128}$ for every $n\ge0$.
  The latter condition is liftability through every finite tower stage;
  it is still a proof obligation, not a consequence of the $E_2$ table.
\end{proposition}
\begin{proof}
  Each next-page boundary is the image of a current differential modulo
  current boundaries. Its source is zero by Proposition~\ref{prop:h6-no-incoming},
  so no new boundaries occur. Induct on the page. In the common-representative
  criterion for survival, not being a boundary on every page now reduces to
  not being a boundary on page two, equivalently $x\ne0$.
  \leanok
\end{proof}

\begin{proposition}[Liftability is independent of the chosen quotient representative]
  \label{prop:h6-specified-representative-liftability}
  \lean{KIP126.Classical.Adams.adamsCycles_mem_iff_of_page_eq,KIP126.Classical.Adams.sphereAdams_h6_nonzeroSurvival_iff_representative}
  \uses{prop:h6-tower-common-representative,prop:h6-survival-liftability}
  \leanok
  Representatives of the same actual quotient class have identical cycle
  membership at every tower stage. Consequently, for any specified
  $u\in Z_2^{2,128}$ representing $x$, nonzero survival of $x$ is equivalent
  to $x\ne0$ and $u\in Z_{n+2}^{2,128}$ for every $n\ge0$.
  No search for a different representative is necessary.
\end{proposition}
\begin{proof}
  The difference of equal quotient representatives is a boundary, hence
  lies in the image of $j$. Exactness $kj=0$ places that difference in
  every cycle submodule. Adding or subtracting it preserves membership.
  Apply the no-incoming-differential survival criterion at $(2,128)$.
  \leanok
\end{proof}

% SOURCE: W. M. Singer, Rings of symmetric functions as modules over the
% Steenrod algebra, AGT 8 (2008), p. 548, after (2-13), verified at
% https://msp.org/agt/2008/8-1/agt-v8-n1-p17-p.pdf
\begin{definition}[The constructed $h_6$ and its square]
  \label{def:h6-constructed-square}
  \lean{KIP126.Steenrod.Milnor.h6Cochain}
  \lean{KIP126.Steenrod.Milnor.h6SquareCochain}
  \lean{KIP126.Classical.Adams.Sphere.classOfMilnorCocycle}
  \lean{KIP126.Classical.Adams.Sphere.h6}
  \lean{KIP126.Classical.Adams.Sphere.h6Square}
  \uses{def:h6-milnor-cooperations,def:h6-constructed-sphere-sequence}
  \notready
  Take the cocycle $[\xi_1^{64}]\in C^{1,64}$ and its concatenation square
  $[\xi_1^{64}\mid\xi_1^{64}]\in C^{2,128}$.
  Apply $\phi^{-1}$ to each, take its first-page homology class, and apply
  the constructed page-passage isomorphism.  This defines respectively
  $h_6\in E_2^{1,64}$ and $h_6^2\in E_2^{2,128}$.
\end{definition}

\begin{proposition}[Zero detection for the specified Milnor classes]
  \label{prop:h6-standard-class-zero-criterion}
  \lean{KIP126.Classical.Adams.Sphere.classOfMilnorCocycle_eq_zero_iff}
  \uses{def:h6-constructed-square,def:h6-milnor-cooperations,prop:h6-tower-page-passage}
  \leanok
  Given explicit Milnor coordinates compatible with the constructed $d_1$,
  a cocycle $x\in C^{s+1,t}$ represents zero in the actual second page if
  and only if $x=db$ for some $b\in C^{s,t}$, for $s,t\ge0$.
  This uses the existing specified-class construction; no Lin presentation
  or specified-class comparison is an input.
\end{proposition}
\begin{proof}
  The page-passage isomorphism detects zero. The homology class of a cycle
  is zero precisely when its representative maps to zero in the cokernel
  of the preceding differential, equivalently when it is a boundary.
  Transport this equation through the Milnor coordinate equivalences and
  their $d_1$ compatibility, in both directions.
  \leanok
\end{proof}

\begin{proposition}[Nonvanishing of the specified standard square]
  \label{prop:h6-standard-square-nonzero}
  \lean{KIP126.Classical.Adams.Sphere.h6Square_ne_zero}
  \uses{def:h6-constructed-square,prop:h6-standard-class-zero-criterion,thm:h6-square-cobar-nonboundary}
  \leanok
  For every supplied Milnor cooperation structure on the constructed sphere
  tower, the specified class $h_6^2\in E_2^{2,128}$ is nonzero.
  The assertion is conditional on that structure, not a construction of it.
\end{proposition}
\begin{proof}
  The zero criterion would make the specified square a cobar boundary,
  contradicting the pure polynomial non-boundary theorem.
  \leanok
\end{proof}

\begin{proposition}[Second-page multiplication and the square identity]
  \label{prop:h6-e2-square-identity}
  \uses{def:h6-constructed-square,def:h6-milnor-cobar,prop:h6-tower-page-passage}
  \notready
  The cobar differential satisfies square-zero and the Leibniz identity;
  concatenation descends to homology and corresponds under a multiplicative
  comparison to a second-page product adding both bidegrees. Under this product,
  the class currently named $h_6^2$ equals the product of $h_6$ with itself.
  Bilinearity and the two individual cocycle calculations above do not
  establish this descent or this second-page identity. Neither the
  second-page product nor the asserted identity has a Lean implementation yet.
\end{proposition}

\begin{definition}[Pagewise nonzero permanence]
  \label{def:h6-pagewise-permanence}
  \lean{KIP126.Core.SpectralSequence.nextPageClass}
  \lean{KIP126.Core.SpectralSequence.PageTrajectory}
  \lean{KIP126.Core.SpectralSequence.IsPermanent}
  \uses{def:h6-constructed-sphere-sequence}
  \notready
  A class $x\in E_2^{s,t}$ is nonzero permanent if there are classes
  $x_n\in E_{2+n}^{s,t}$ for every $n\ge0$, with $x_0=x$, such that
  $d_{2+n}x_n=0$, the page-passage image of $[x_n]$ is $x_{n+1}$, and
  $x_n\ne0$ for every $n$.  This definition uses the actual differentials
  and homology quotients and makes no claim of convergence to an abutment.
\end{definition}

\begin{definition}[Fixed standard and computational sphere objects]
  \label{def:h6-fixed-computation}
  \lean{KIP126.Classical.Adams.sphereAdams,KIP126.Classical.Adams.sphereH6Square,KIP126.Classical.Adams.sphereAdamsModel,KIP126.Classical.Adams.sphereAdamsData,KIP126.Classical.Adams.sphereAdams_towerComparison,KIP126.Classical.Adams.survival_comparison,KIP126.Classical.Adams.computedH6Square,KIP126.Core.SpectralSequence.NonzeroSurvival}
  \uses{def:h6-constructed-square,def:h6-pagewise-permanence,def:h6-tower-internal-sequence,prop:h6-tower-internal-comparison,prop:h6-tower-survival-comparison}
  \notready
  Select the stable foundation, $H\mathbb F_2$, and Milnor coordinates once.
  The standard sequence and class specialize the existing tower construction.
  The computational object specializes the proved internal tower construction
  to the same unit and sphere. It is no longer independently postulated.
  Its specified
  class is the image of generator 69 squared in the fixed Lin CSV quotient
  (Zenodo 14875701, v126.3.cw49, internal degree at most 261).
  Its nonzero survival means a common $Z_\infty$ representative whose $E_2$
  image is the specified class and whose $E_\infty$ image is nonzero.
  The fixed all-page tower comparison specializes the proved generic
  comparison; it too is no longer an axiom. The survival comparison is now
  the specialization of the proved generic equivalence. Foundation selection,
  Milnor coordinates, and the range-limited Lin comparison remain named
  development axioms. The specified-class comparison is now proved using
  second-page nonvanishing and the unique nonzero element in bidegree $(2,128)$.
  None of these statements asserts that the
  specified class survives. The internal
  computational proof imports neither KIPBase nor the Mathlib SS adapter.
  The chosen abstract foundation has not been realized as a concrete model
  of spectra. Nor does this construction establish the paper's exclusion
  of differentials on the specified class.
\end{definition}

\begin{proposition}[Fixed standard square: audited second-page nonvanishing]
  \label{prop:h6-fixed-standard-square-nonzero}
  \lean{KIP126.Classical.Adams.sphereH6Square_ne_zero}
  \uses{def:h6-fixed-computation,prop:h6-standard-square-nonzero}
  \notready
  The fixed class \texttt{sphereH6Square} is nonzero on its constructed
  second page. The Lean proof specializes the generic theorem and has no
  local or transitive \texttt{sorry}; its named project assumptions are
  exactly \texttt{standardFoundation} and \texttt{standardMilnorCooperations}.
  This axiom-dependent specialization is not an unconditional completion
  node. It does not use the Lin presentation or the specified-class
  comparison, prove that comparison, or address higher-page survival.
\end{proposition}

\begin{definition}[Fixed Lin algebra and additive catalogue]
  \label{def:lin-e2-catalogue}
  \lean{KIP126.LinE2.E2}
  \lean{KIP126.LinE2.basisRowsAt}
  \lean{KIP126.LinE2.findBasisIndex?}
  \lean{KIP126.LinE2.mulAt}
  \lean{KIP126.LinE2.Compute.Engine.multiplyBasis}
  \uses{def:h6-fixed-computation}
  \notready
  Use the fixed Zenodo 14875701 dataset, version v126.3.cw49, imported from
  PR 110 at commit \texttt{ff39e951}. Its quotient algebra has 2914 algebra
  generators and 231848 relations and is truncated above internal degree 261.
  The separate additive catalogue has 23822 monomials, indexed by $(s,t,k)$;
  these indices are not algebra-generator numbers. The catalogue entry
  $(2,128,0)$ is generator 69 squared. Lookup rejects absent coordinates.
  Multiplication is that of the quotient, with homogeneous closure proved;
  the executable reducer reports invalid input, range violations and fuel
  exhaustion rather than silently returning a value.
\end{definition}

\begin{proposition}[The unique possible square-degree monomial]
  \label{prop:lin-square-degree-unique}
  \lean{KIP126.LinE2.generatorDegree_eq_row,KIP126.LinE2.generatorDegree_low_filtration,KIP126.LinE2.squareDegree_support,KIP126.LinE2.monomialDegree_square_unique}
  \uses{def:lin-e2-catalogue}
  \leanok
  A polynomial monomial of bidegree $(2,128)$ in the actual archived
  generator degrees is exactly generator 69 squared. Every supported
  generator lies among indices $0,1,2,3,7,18,69,324$.
\end{proposition}
\begin{proof}
  A kernel-checked finite certificate exhausts the generator table in
  filtration at most two and internal degree at most 128. Each supported
  generator obeys these bounds because its contribution is at most the
  total degree. For the eight remaining exponents $a_0,\ldots,a_7$, the
  degree equations are $\sum a_i=2$ and
  $\sum 2^i a_i=128$. Nonnegative integer arithmetic forces $a_6=2$
  and every other exponent zero. No relation reduction or additive-basis
  certification is used.
  \leanok
\end{proof}

\begin{proposition}[The square-degree component has at most two elements]
  \label{prop:lin-square-component-two-elements}
  \lean{KIP126.LinE2.homogeneousPart_square_eq_span,KIP126.LinE2.E2At_square_eq_zero_or}
  \uses{prop:lin-square-degree-unique}
  \leanok
  The actual homogeneous submodule of the archived quotient in bidegree
  $(2,128)$ is the $\mathbb F_2$-span of \texttt{dataH6Sq}. Consequently
  each element is either zero or that square. The independent nonvanishing
  certificate makes these two elements distinct, but is not needed for
  the exhaustive alternative itself.
\end{proposition}
\begin{proof}
  Replace every monomial in the defining spanning set by the unique
  monomial just proved. The two possible scalar coefficients are zero and one.
  \leanok
\end{proof}

\begin{proposition}[The fixed internal square is the unique nonzero element]
  \label{prop:h6-internal-square-unique}
  \lean{KIP126.Classical.Adams.sphereAdamsData_square_eq_zero_or,KIP126.Classical.Adams.sphereAdamsData_eq_computedH6Square_of_ne_zero}
  \uses{def:h6-fixed-computation,prop:lin-square-component-two-elements}
  \notready
  Under the existing fixed foundation and Lin presentation, every element
  of the internal $E_2^{2,128}$ is zero or \texttt{computedH6Square}.
  In particular every nonzero element equals that square. This is a
  complete Lean proof by surjectivity of the existing Lin comparison,
  but retains precisely those two fixed project assumptions. It imports
  no Milnor coordinates or Mathlib spectral-sequence adapter.
\end{proposition}

\begin{proposition}[The actual double tensor-boundary class is the computed square]
  \label{prop:h6-computed-double-representative}
  \lean{KIP126.Classical.Adams.sphereH6DoubleInternalE2_eq_computedH6Square,KIP126.Classical.Adams.computedH6Square_double_representative}
  \uses{prop:h6-internal-square-unique,prop:h6-double-internal-nonzero,prop:h6-internal-square-polynomial-provenance}
  \notready
  Supply the ring, K\"unneth, reduced Milnor-basis structures and their
  suspension, diagonal, unit, and coproduct coherence conditions on the
  selected foundation. For the cooperation $a$ with polynomial coordinate
  $[\xi_1^{64}]$ and normalized sphere coefficient, the actual internal
  double tensor-boundary class equals \texttt{computedH6Square}.
  It therefore supplies a $Z_2$ representative of that fixed class whose
  first-page image is the specified double tensor-boundary element and
  whose derived polynomial coordinate is $[\xi_1^{64}\mid\xi_1^{64}]$.
  The proof uses only the fixed foundation and Lin-presentation axioms;
  the lower structures and coherences are explicit, still-unsupplied inputs.
  In particular it neither imports nor assumes the fixed Milnor-coordinate
  comparison axiom, and does not use a Mathlib spectral-sequence adapter.
\end{proposition}
\begin{proof}
  The actual double class is nonzero by the pure cobar non-boundary theorem.
  The internal square-degree uniqueness result makes it the fixed computed
  class. Transport the previously constructed representative and its
  derived polynomial coordinate along this equality.
\end{proof}

\begin{proposition}[The remaining computational target as actual tower lifts]
  \label{prop:h6-computed-double-tower-lifts}
  \lean{KIP126.Classical.Adams.computedH6Square_nonzeroSurvival_iff_double_lifts,KIP126.Classical.Adams.computedH6Square_nonzeroSurvival_iff_double_connecting_lifts}
  \uses{prop:h6-computed-double-representative,prop:h6-specified-representative-liftability,def:h6-tower-pages}
  \notready
  Under those same explicit lower inputs, choose any $Z_2$ lift $u$ of the
  actual double tensor-boundary first-page element. The fixed computational
  permanence statement is equivalent to
  \[
    \forall n\ge0,\quad \exists y_n\in\pi_{125}(T_{n+4}S),\qquad
    (T_{n+4}S\longrightarrow T_3S)_*(y_n)=k(u).
  \]
  Here $k$ is the actual cofiber connecting map and
  $k(u)\in\pi_{125}(T_3S)$. Existence of these lifts is an open mathematical
  obligation, not a new axiom or a consequence of the $E_2$ table. The
  equivalence asks for finite-stage lifts, not a compatible chosen family
  or an inverse-limit element, and asserts no abutment convergence.
\end{proposition}
\begin{proof}
  Every such $u$ represents the computed square, whose nonvanishing also
  follows from the actual double-class calculation. Apply representative
  independence and the no-incoming-differential criterion, then unfold the
  actual definition of $Z_{n+2}$ as the preimage under $k$ of the tower-map
  image. The degrees are $128-2-1=125$ and $2+(n+2)=n+4$.
\end{proof}

\begin{proposition}[The internal second differential is a tower obstruction]
  \label{prop:h6-internal-second-differential-obstruction}
  \lean{KIP126.Classical.Adams.adamsTowerPreSS_d_two_h6_eq,KIP126.Classical.Adams.adamsTower_d_two_h6_comparison,KIP126.Classical.Adams.adamsTower_d_two_h6_representative,KIP126.Classical.Adams.adamsTower_d_two_h6_value_of_lift,KIP126.Classical.Adams.adamsTower_d_two_h6_value_exists,KIP126.Classical.Adams.adamsTower_d_two_h6_eq_zero_iff_lift}
  \uses{def:h6-tower-internal-sequence,prop:h6-tower-internal-laws,def:h6-tower-differential}
  \leanok
  For any tower in the stated stable-category context, let $u\in Z_2^{2,128}$
  represent an internal class $x\in E_2^{2,128}$. There exists
  $y\in\pi_{125}(T_4X)$ mapping to $k(u)\in\pi_{125}(T_3X)$.
  Every such lift computes the existing internal differential as
  $d_2(x)=[j(y)]\in E_2^{4,129}$, using the specified quotient-page
  comparison. Moreover,
  \[
    d_2(x)=0\quad\Longleftrightarrow\quad
    \exists z\in\pi_{125}(T_5X),\quad
    (T_5X\longrightarrow T_3X)_*(z)=k(u).
  \]
  This is independent of the fixed foundation, Lin presentation, Milnor
  coordinates, and Mathlib spectral-sequence adapter.
\end{proposition}
\begin{proof}
  Check the numerical reindexing of the internal differential and transport
  the tower quotient differential through the existing page isomorphisms.
  Membership in $Z_2$ supplies $y$. The lift-independent differential formula
  gives its value, and the kernel criterion identifies vanishing with $Z_3$.
  \leanok
\end{proof}

\begin{proposition}[The computed square's second differential as a concrete obstruction]
  \label{prop:h6-computed-second-differential-obstruction}
  \lean{KIP126.Classical.Adams.computedH6Square_of_double_representative,KIP126.Classical.Adams.computedH6Square_d_two_value_of_double_lift,KIP126.Classical.Adams.computedH6Square_d_two_double_value_exists,KIP126.Classical.Adams.computedH6Square_d_two_eq_zero_iff_double_lift}
  \uses{prop:h6-computed-double-representative,prop:h6-internal-second-differential-obstruction}
  \notready
  Supply the same explicit lower structures and coherences as for the
  actual double tensor-boundary representative. For any $Z_2$ representative
  $u$ of that first-page element, the preceding formula computes
  $d_2(\texttt{computedH6Square})$ in the existing internal sequence.
  Its vanishing is equivalent to lifting $k(u)$ from $T_3S$ to $T_5S$.
  These Lean proofs retain exactly the fixed foundation and Lin-presentation
  project assumptions, with the lower witnesses still explicit and unsupplied.
  No $T_5$ lift or vanishing assertion is provided, and no multiplicativity
  or Leibniz rule for the spectral sequence is assumed by these formulas.
\end{proposition}
\begin{proof}
  Any such representative has second-page class equal to the computed square.
  Substitute this equality in the generic internal differential formula
  and its vanishing criterion.
\end{proof}

\begin{definition}[The actual second differential in bounded Lin coordinates]
  \label{def:h6-lin-second-differential}
  \lean{KIP126.Classical.Adams.sphereE2SecondDifferential,KIP126.Classical.Adams.LinE2Presentation.secondDifferential,KIP126.Classical.Adams.LinE2Presentation.SecondDifferentialLeibniz}
  \uses{def:h6-fixed-computation,def:h6-tower-internal-sequence}
  \notready
  For a supplied Lin presentation $P$, express the existing internal tower
  differential in coordinates as
  \[
    D_P^{s,t}=P_{s+2,t+1}^{-1}\,d_2\,P_{s,t},\qquad t+1\le261.
  \]
  This is a family on homogeneous components, not a freely chosen
  differential on the whole truncated quotient. The predicate
  \texttt{SecondDifferentialLeibniz} requires
  $D_P(xy)=D_P(x)y+xD_P(y)$ for all homogeneous $x,y$ with
  $t(x)+t(y)+1\le261$. The equality is taken in the common Lin quotient.
  The predicate is defined but has no supplied witness: neither the
  presentation axiom nor the imported multiplication table asserts it.
  Establishing it requires compatibility with the actual tower differential.
\end{definition}

\begin{proposition}[Leibniz suffices for the square's first outgoing differential]
  \label{prop:h6-d2-zero-from-leibniz}
  \lean{KIP126.Classical.Adams.sphereE2SecondDifferential_h6_square,KIP126.Classical.Adams.LinE2Presentation.secondDifferential_eq_zero_iff,KIP126.Classical.Adams.linE2_add_self_eq_zero,KIP126.Classical.Adams.LinE2Presentation.secondDifferential_square_eq_zero,KIP126.Classical.Adams.computedH6Square_d_two_eq_zero_of_leibniz,KIP126.Classical.Adams.computedH6Square_double_lift_five_of_leibniz}
  \uses{def:h6-lin-second-differential,prop:h6-computed-second-differential-obstruction}
  \notready
  If the selected Lin presentation satisfies the preceding Leibniz
  condition, then the actual internal $d_2(\texttt{computedH6Square})$ is zero.
  More generally every homogeneous square in the covered range is a
  $D_P$-cycle. No assertion that the unsquared class is a cycle is required.
  Under the additional explicit lower structures identifying the actual
  double representative $u$, this gives a lift of $k(u)$ to $\pi_{125}(T_5S)$.
  These are conditional Lean proofs, not a witness of Leibniz compatibility.
  The fixed conclusions retain the existing foundation and Lin-presentation
  axioms, introduce no new axiom, and imply neither higher outgoing
  differential vanishing nor permanent survival.
\end{proposition}
\begin{proof}
  The actual Lin quotient has $z+z=0$, by its $\mathbb F_2$-module structure;
  no properness of the defining ideal is needed. Commutativity gives
  $D_P(x^2)=D_P(x)x+xD_P(x)=0$.
  Set $x=\texttt{dataH6}$, identify its product with \texttt{dataH6Sq}, and
  transport vanishing through the invertible target coordinate map.
  The previous tower criterion then supplies the $T_5S$ lift.
\end{proof}

\begin{proposition}[Three geometric products suffice for the square's second differential]
  \label{prop:h6-d2-zero-from-long-layer}
  \lean{KIP126.Classical.Adams.SphereH6LongLayerMaps.LinCompatible,KIP126.Classical.Adams.SphereH6LongLayerMaps.LinCompatible.cross_sum_zero,KIP126.Classical.Adams.LinE2Presentation.h6_cross_products_add_eq_zero,KIP126.Classical.Adams.computedH6Square_d_two_eq_zero_of_longLayer}
  \uses{def:h6-three-long-layer-products,def:h6-fixed-computation,prop:h6-long-layer-cross-cancellation,prop:h6-internal-second-differential-obstruction}
  \notready
  Assume the fixed foundation has the indicated tensor and coefficient-ring
  structures, the three actual maps satisfy the preceding compatibility
  and relative-boundary predicates, and their three descended internal
  products equal the corresponding products of the fixed Lin presentation.
  Then $d_2(\texttt{computedH6Square})=0$ for the existing internal
  differential. A global \texttt{SecondDifferentialLeibniz} premise is not
  needed, nor is $d_2(\texttt{computedH6})=0$.
  This conditional implication has a Lean proof; the geometric maps and
  their predicates remain unconstructed for the fixed foundation. Its
  existing foundation and Lin-presentation axioms remain disclosed.
  The result asserts neither higher differential vanishing nor permanence.
\end{proposition}
\begin{proof}
  Apply the internal relative-boundary formula to the two copies of the
  computed $h_6$. The square comparison identifies the input product with
  the computed square. The two cross comparisons turn its differential
  into $D(h_6)h_6+h_6D(h_6)$. Take Lin preimages of the two factors in
  bidegrees $(3,65)$ and $(1,64)$; their products agree by commutativity and
  their sum is zero by characteristic two in the covered degree $(4,129)$.
  Transport this cancellation through the target comparison and its
  injective quotient-page isomorphism.
\end{proof}

\begin{proposition}[The fixed second differential from lower connecting lifts]
  \label{prop:h6-d2-zero-from-boundary-lifts}
  \lean{KIP126.Classical.Adams.computedH6Square_d_two_eq_zero_of_boundaryLifts,KIP126.Classical.Adams.computedH6Square_d_two_eq_zero_of_firstCycleProductRule}
  \uses{def:h6-three-boundary-lifts,prop:h6-filtration-one-boundaries,prop:h6-cross-boundary-from-first-differential,prop:h6-d2-zero-from-long-layer}
  \notready
  Under the same fixed tensor and coefficient-ring assumptions, suppose
  the three lower lifts satisfy \texttt{FitsProducts}, and the resulting
  products satisfy the two \texttt{CrossBoundaryCompatible} conditions, the relative-boundary
  identity, and the three Lin product comparisons. Then
  $d_2(\texttt{computedH6Square})=0$. The long-layer products and their
  projection squares are constructed, rather than additional inputs;
  four of the six original descent boundary conditions are now proved.
  Alternatively, the two local first-differential product rules supply
  the remaining cross conditions. They do not replace the separate
  second-differential relative-boundary identity.
  This implication is proved in Lean without new axioms or \texttt{sorry};
  the existing foundation and Lin-presentation axioms remain. The actual
  lifts and the listed compatibility witnesses are still missing, so this
  is not an unconditional vanishing or permanence theorem.
\end{proposition}
\begin{proof}
  Construct the three products from the prescribed lifts, discharge their
  projection compatibility by construction and four boundary conditions
  by filtration-one boundary vanishing, and apply the preceding
  three-product vanishing theorem with the remaining hypotheses.
\end{proof}

\begin{proposition}[The specified-class comparison is a theorem]
  \label{prop:h6-fixed-class-comparison}
  \lean{KIP126.Classical.Adams.h6Square_comparison}
  \uses{prop:h6-internal-square-unique,prop:h6-fixed-standard-square-nonzero,prop:h6-tower-internal-comparison}
  \notready
  The actual page comparison carries \texttt{computedH6Square} to
  \texttt{sphereH6Square}. This is no longer an independent project axiom:
  the Lean proof has no \texttt{sorry} and depends on exactly the fixed
  foundation, Milnor coordinates, and Lin presentation. Those remaining
  inputs prevent unconditional completion of this fixed node.
\end{proposition}
\begin{proof}
  Pull the standard class back through the actual second-page isomorphism.
  Its preimage is nonzero by standard nonvanishing, hence equals the
  computational square by the exhaustive internal alternative. Map back.
  The proof neither presupposes the class comparison nor asserts permanence.
\end{proof}

\begin{definition}[A finite nonzero-square detector]
  \label{def:lin-square-detector}
  \lean{KIP126.LinE2.SquareDetection.Target,KIP126.LinE2.SquareDetection.evaluate,KIP126.LinE2.SquareDetection.relationCheck,KIP126.LinE2.SquareDetection.allRelationsCheck}
  \uses{def:lin-e2-catalogue}
  \leanok
  Let $T=\mathbb F_2[u]/(u^3)$. Evaluate generator 69 at $u$ and all other
  generators at zero. The executable relation test checks that every
  monomial of each archived relation individually evaluates to zero: it
  contains a positive power of another generator or at least three copies
  of generator 69. The test follows the exact parser conventions of the
  quotient model, including malformed strings. Its value is a Boolean,
  not a theorem asserting that the check succeeds.
\end{definition}

\begin{proposition}[Soundness of the nonzero-square detector]
  \label{prop:lin-square-detector-sound}
  \lean{KIP126.LinE2.SquareDetection.u_square_ne_zero,KIP126.LinE2.SquareDetection.relationCheck_sound,KIP126.LinE2.SquareDetection.definingIdeal_le_ker,KIP126.LinE2.SquareDetection.dataH6Sq_ne_zero_of_check}
  \uses{def:lin-square-detector}
  \leanok
  If the finite test \texttt{allRelationsCheck} equals \texttt{true}, then
  the specified data square is nonzero in the Lin quotient algebra.
  This conditional theorem requires neither additive-basis certification
  nor the admitted soundness of the general Gr\"obner reducer.
\end{proposition}
\begin{proof}
  Induct on the parsed generator/exponent list to verify the evaluation
  test. The checked relations lie in the kernel. A monomial of internal
  degree greater than 261 either involves a killed variable or is a pure
  power of generator 69 with exponent at least five, so the truncation
  ideal also lies in the kernel. Finally $u^2\ne0$, since $X^3$ cannot
  divide $X^2$ in $\mathbb F_2[X]$. Hence generator 69 squared cannot
  belong to the defining ideal.
  \leanok
\end{proof}

\begin{proposition}[A semantics-preserving character-list certificate checker]
  \label{prop:lin-square-parser-equivalence}
  \lean{KIP126.LinE2.SquareDetection.splitOn_singleton_eq_list,KIP126.LinE2.SquareDetection.toNat?_eq_chars,KIP126.LinE2.SquareDetection.relationCheck_eq_chars,KIP126.LinE2.SquareDetection.chunkCheck_eq_chars}
  \uses{def:lin-square-detector}
  \leanok
  For every string, splitting at a singleton separator agrees with splitting
  its character list. The character-list natural-number parser agrees with
  the library parser, including the underscore rules and malformed inputs.
  Hence the character-list relation and chunk checkers return exactly the
  same Booleans as the original archived-string checker.
\end{proposition}
\begin{proof}
  Compare the substring split recursion with the character split recursion:
  matching a one-character separator resets its offset immediately. Use the
  existing characterization of character splitting on lists. Rewrite the
  library numeric parser's character iteration and fold as list operations.
  Compose these equalities through the monomial and relation checks.
  \leanok
\end{proof}

\begin{proposition}[Composition of finite certificates]
  \label{prop:lin-square-certificate-composition}
  \lean{KIP126.LinE2.SquareDetection.charsChunkCheck_join,KIP126.LinE2.SquareDetection.charsChunkCheck_append,KIP126.LinE2.SquareDetection.charsChunkCheck_ofList,KIP126.LinE2.SquareDetection.firstRelations_check,KIP126.LinE2.SquareDetection.allRelationsCheck_of_chunks}
  \uses{prop:lin-square-parser-equivalence}
  \leanok
  Two certified chunks can be joined by a newline. The three separately
  stored initial relations pass the checker by a kernel-checked computation.
  If each remaining archived chunk has a character-list certificate, then
  the original full-table check equals \texttt{true}.
  This statement does not assert that the remaining 227 chunk certificates
  have been constructed.
\end{proposition}
\begin{proof}
  Splitting a list at an inserted separator concatenates the two split lists;
  the Boolean universal check becomes a conjunction. For the full archive,
  decompose membership in the initial list and the flattened chunk list,
  then apply the corresponding certificate and the proved parser comparison.
  \leanok
\end{proof}

\begin{theorem}[Archived square certificate and initial nonvanishing]
  \label{thm:lin-square-certified}
  \lean{KIP126.LinE2.SquareDetection.archivedChunks_check,KIP126.LinE2.SquareDetection.allRelationsCheck_eq_true,KIP126.LinE2.SquareDetection.dataH6Sq_ne_zero}
  \uses{prop:lin-square-certificate-composition,prop:lin-square-detector-sound}
  \leanok
  All 227 archived chunks pass the finite detector. Including the three
  separately stored relations, all 231848 relations pass. Consequently the
  specified square is nonzero in the actual Lin quotient algebra, without
  a finite-check hypothesis or additive-basis certification.
\end{theorem}
\begin{proof}
  Generate small character-list leaf proofs by kernel reduction. Join the
  certificates using the proved list concatenation rule, and transport them
  back with $\operatorname{toList}(\operatorname{ofList}(l))=l$. The resulting
  proof types must match the original archived string literals. Assemble
  eight consecutive batches covering the entire archive and apply the
  soundness theorem. Native execution only generates candidate proof terms;
  it supplies no evaluation axiom and does not replace any archived data.
  \leanok
\end{proof}

\begin{proposition}[Fixed computational nonvanishing and the remaining target]
  \label{prop:h6-certified-nonvanishing}
  \lean{KIP126.Classical.Adams.computedH6Square_ne_zero,KIP126.Classical.Adams.computedH6Square_nonzeroSurvival_iff}
  \uses{thm:lin-square-certified,prop:h6-survival-liftability,def:h6-fixed-computation}
  \notready
  Under the already disclosed fixed foundation and Lin comparison, the
  computational class is nonzero on $E_2$. Using the proved
  no-incoming-differential theorem, its nonzero survival is equivalent to
  the existence of one actual tower $Z_2$ representative of the specified
  class that belongs to every $Z_{n+2}$. No finite-check premise remains;
  the all-stage lifting assertion is not proved by this equivalence.
\end{proposition}
\begin{proof}
  Transport nonvanishing through the injective Lin comparison. Substitute
  this theorem into the already proved survival/liftability equivalence.
  Both proof bodies are implemented and contain no unfinished proof;
  this fixed-object node remains conditional on the named foundation,
  and Lin comparison inputs and is therefore not marked complete. This
  reduction no longer depends on the named Milnor-coordinate assumption.
\end{proof}

The regression module \texttt{Checks/AdamsE2/SquareDetection} audits the
finite certificate and its transfers, rejecting \texttt{sorryAx} and
native-evaluation axioms and listing the permitted fixed comparison inputs
separately. Its additional runtime test is not used as a proof. Initial
nonvanishing neither removes the Lin comparison axiom nor proves higher-page
survival; the computational final Solution remains open.

\begin{theorem}[Additive catalogue certification, pending]
  \label{thm:lin-e2-basis-certification}
  \lean{KIP126.LinE2.basisTable_correct}
  \uses{def:lin-e2-catalogue}
  \notready
  For $t\le261$, the actual listed monomial values of degree $(s,t)$ form an
  $\mathbb F_2$-basis of the specified quotient component. This asserts both
  linear independence and spanning, not merely that the list has the right
  length. The Lean proof remains admitted; executable checks of degrees,
  distinct coordinates and irreducibility do not discharge it.
\end{theorem}

\begin{definition}[Data and internal Adams additive coordinates]
  \label{def:lin-e2-coordinates}
  \lean{KIP126.LinE2.dataBasis}
  \lean{KIP126.LinE2.dataCoordinates}
  \lean{KIP126.Classical.Adams.sphereE2BasisByCSV?}
  \lean{KIP126.Classical.Adams.sphereE2Coordinates}
  \lean{KIP126.Classical.Adams.linToSphere_product_eq}
  \lean{KIP126.Classical.Adams.linToSphere_exists_preimage}
  \lean{KIP126.Classical.Adams.linToSphere_eq_iff}
  \lean{KIP126.Classical.Adams.linToSphere_product_eq_zero}
  \uses{thm:lin-e2-basis-certification,def:h6-fixed-computation}
  \notready
  Under the pending basis certification, obtain unique finite-support
  $\mathbb F_2$ coordinates on each covered component. Transport them through
  the fixed Lin comparison to the internal SSData $E_2$ page, using an
  integer-linear equivalence to retain its existing additive structure.
  Basis vectors have the specified CSV values and are nonzero; coordinates
  reconstruct any element. Product equalities in the data quotient transfer
  to the internal page, including zero products. Every covered page element
  has a data preimage, and equality of images is equivalent to equality of
  data representatives. These interfaces import no Mathlib spectral-sequence
  adapter, and neither they nor the data assert higher-page survival.
\end{definition}

\begin{lemma}[The computed class is the square of computed $h_6$]
  \label{lem:lin-h6-square-product}
  \lean{KIP126.Classical.Adams.computedH6_mul_self}
  \uses{def:h6-fixed-computation,def:lin-e2-catalogue}
  \notready
  Under the fixed Lin comparison, the internal $E_2$ product of the image of
  data $h_6$ with itself is the computational target class. This replaces
  PR 110's historical \texttt{h6Sq\_eq\_comparison}, with no additional
  admitted theorem. It neither asserts permanence nor eliminates the
  existing presentation axiom.
\end{lemma}

\begin{theorem}[Coordinate-check soundness, pending]
  \label{thm:lin-e2-coordinate-check}
  \lean{KIP126.LinE2.Automation.coordinateCheck_sound}
  \uses{def:lin-e2-catalogue}
  \notready
  If the concrete reducer successfully checks equal CSV coordinates for two
  expressions within range, their interpretations in the quotient algebra
  are equal. The migrated \texttt{e2\_mul} tactic executes the check but uses
  this admitted theorem to close goals. This algorithm-correctness debt is
  separate from additive-basis certification and from the comparison with
  the actual sphere Adams page. On Lean 4.32, the tactic's native evaluation
  also generates evaluation axioms; the strict project audit rejects them.
  Kernel-checkable reduction certificates remain a separate required step.
\end{theorem}
